\documentclass[11pt,oneside]{book}
\usepackage[margin=1.2in]{geometry}
\usepackage[toc,page]{appendix}
\usepackage{graphicx}
\usepackage{natbib}
\usepackage{lipsum}
\usepackage{caption}
\usepackage[T1]{fontenc}
\usepackage{titlesec, blindtext, color}
\usepackage{xcolor,tikz}
\usetikzlibrary{patterns}
\usetikzlibrary{decorations.markings}
\usetikzlibrary{decorations.pathmorphing}
\usepackage{amsmath,amssymb,amsthm,mathrsfs,amsfonts,xfrac,pifont,bbold,physics,enumitem}
\usepackage[utf8]{inputenc}
\usepackage{amsthm}
\usepackage[breakable, theorems, skins]{tcolorbox}
\usepackage[colorlinks = true,
            linkcolor = red,
            urlcolor  = blue,
            citecolor = red,
            anchorcolor = red]{hyperref}
\usepackage{cleveref}

\usepackage{simplewick}

\usepackage{soul}
\usepackage{frcursive}
\usepackage{booktabs}


\usepackage{array}
\newcolumntype{C}[1]{>{\centering\arraybackslash}m{#1}}


% -------------------------------------------------------------------
% Theorem Styles
% -------------------------------------------------------------------

\theoremstyle{definition} % Define theorem styles here based on the definition style (used for definitions and examples)
\newtheorem*{definition}{Definition}

\theoremstyle{plain} % Define theorem styles here based on the plain style (used for theorems, lemmas, propositions)
\newtheorem{theorem}{Theorem}[section]
\newtheorem{axiom}{Axiom}
\newtheorem{corollary}[theorem]{Corollary}
\newtheorem{lemma}[theorem]{Lemma}
\newtheorem{proposition}[theorem]{Proposition}

\theoremstyle{remark} % Define theorem styles here based on the remark style (used for remarks and notes)
\newtheorem*{notation}{Notation}
\newtheorem*{solution}{Solution}

\newtheoremstyle{underline}% name
{}        % Space above, empty = `usual value'
{}              % Space below
{}              % Body font
{}    % Indent amount (empty = no indent, \parindent = para indent)
{}              % Thm head font
{.}             % Punctuation after thm head
{1.5mm}         % Space after thm head: \newline = linebreak
{{\underline{\textit{\thmname{#1}\thmnumber{ #2}}~\thmnote{(#3)}\unskip}}}  % Thm head spec

\theoremstyle{underline}

\newtheorem{remark}[theorem]{Remark}
\newtheorem{example}[theorem]{Example}
\newtheorem{claim}[theorem]{Claim}
% -------------------------------------------------------------------
% Chapter Headings
% -------------------------------------------------------------------

\setcounter{chapter}{-1}

\makeatletter
\renewcommand{\@chapapp}{Lecture}
\makeatother
\definecolor{lightergray}{rgb}{0.9,0.9,0.9}

\usepackage{titlesec}
\titleformat{\section}{\large\bfseries\raggedright}{}{0em}{\colorsection}[\titlerule]
\titleformat{name=\section,numberless}{\large\scshape\bfseries\raggedright}{}{0em}{\colorsectionnonumber}[\titlerule]

\titleformat{\subsection}{\bfseries\raggedright}{}{0em}{\colorsubsection}
\titleformat{name=\subsection,numberless}{\bfseries\raggedright}{}{0em}{\colorsubsectionnonumber}

\newcommand{\colorsection}[1]{%
    \colorbox{lightergray}{\parbox{\dimexpr\textwidth-2\fboxsep}{\thesection\ \ #1}}}
\newcommand{\colorsectionnonumber}[1]{%
    \colorbox{lightergray}{\parbox{\dimexpr\textwidth-2\fboxsep}{#1}}}
    
\newcommand{\colorsubsection}[1]{%
    \colorbox{lightergray}{\parbox{\dimexpr\textwidth-2\fboxsep}{\thesubsection\ #1}}}
\newcommand{\colorsubsectionnonumber}[1]{%
    \colorbox{lightergray}{\parbox{\dimexpr\textwidth-2\fboxsep}{#1}}}
    
\definecolor{gray75}{gray}{0.75}
\newcommand{\hsp}{\hspace{20pt}}
\titleformat{\chapter}[hang]{\Huge\bfseries}{\thechapter\hsp\textcolor{gray75}{|}\hsp}{0pt}{\Huge\bfseries}

% Begin/end Shortcuts
\def\bc{\begin{corollary}}
\def\ec{\end{corollary}}
\def\bcl{\begin{claim}}
\def\ecl{\end{claim}}
\def\bd{\begin{definition}}
\def\ed{\end{definition}}
\def\ben{\begin{enumerate}}
\def\een{\end{enumerate}}
\def\be{\begin{equation}}
\def\ee{\end{equation}}
\def\bse{\begin{equation*}}
\def\ese{\end{equation*}}
\def\bex{\begin{example}}
\def\eex{\end{example}}
\def\bit{\begin{itemize}}
\def\eit{\end{itemize}}
\def\bl{\begin{lemma}}
\def\el{\end{lemma}}
\def\bnn{\begin{notation}}
\def\enn{\end{notation}}
\def\bn{\begin{note}}
\def\en{\end{note}}
\def\bp{\begin{proposition}}
\def\ep{\end{proposition}}
\def\bq{\begin{proof}}
\def\eq{\end{proof}}
\def\br{\begin{remark}}
\def\er{\end{remark}}
\def\bs{\begin{solution}}
\def\es{\end{solution}}
\def\btab{\begin{table}}
\def\etab{\end{table}}
\def\btb{\begin{tabular}}
\def\etb{\end{tabular}}
\def\bt{\begin{theorem}}
\def\et{\end{theorem}}
\def\btik{\begin{tikzpicture}}
\def\etik{\end{tikzpicture}}

% Letter Shortcuts
% Greek letters
\def\a{\alpha}
\def\b{\beta}
\def\del{\delta}
\def\g{\gamma}
\def\l{\lambda}
\def\n{\nabla}
\def\sig{\sigma}

% Mathbb
\def\b1{\mathbb{1}} % identity
\def\C{\mathbb{C}} % complex numbers
\def\F{\mathbb{F}}
\def\N{\mathbb{N}} % Natural
\def\Q{\mathbb{Q}} % Rational
\def\R{\mathbb{R}} % Reals
\def\Z{\mathbb{Z}} % Integers

% Mathcal
\def\cA{\mathcal{A}}
\def\cB{\mathcal{B}}
\def\cD{\mathcal{D}}
\def\cF{\mathcal{F}}
\def\cH{\mathcal{H}}
\def\cJ{\mathcal{J}}
\def\cK{\mathcal{K}}
\def\cL{\mathcal{L}}
\def\cM{\mathcal{M}}
\def\cN{\mathcal{N}}
\def\cO{\mathcal{O}}
\def\cP{\mathcal{P}}
\def\cT{\mathcal{T}}

% Mathfrak 
\DeclareMathOperator{\GL}{GL}
\def\gl{\mathfrak{gl}}
\DeclareMathOperator{\Ort}{O}
\def\ort{\mathfrak{o}}
\DeclareMathOperator{\SL}{SL}
\def\sl{\mathfrak{sl}}
\DeclareMathOperator{\SO}{SO}
\def\so{\mathfrak{so}}
\DeclareMathOperator{\SU}{SU}
\def\su{\mathfrak{su}}
\def\fF{\mathfrak{F}}

% Misc. 
\def\diag{\text{diag}}
\def\p{\partial}
\def\ss{\subset}
\def\se{\subseteq}
\def\sm{\setminus}
\def\ve{\varepsilon}
\def\vn{\varnothing}
\def\cl{\colon}
\def\lra{\overset{\leftrightarrow}}
\def\la{\langle}
\def\ra{\rangle}

% Boxes 
\DeclareRobustCommand{\mybox}[2][colframe=blue!10!black,before skip=10pt]{%
\begin{tcolorbox}[
        breakable,
        width=\dimexpr\textwidth\relax, 
        boxsep=5pt,
        arc=5pt,outer arc=5pt,
        ]
        #2
\end{tcolorbox}
}

\def\bbox{\begin{tcolorbox}[colback=white!5,colframe=blue!75!black,title=Exercise]}
\def\ebox{\end{tcolorbox}}

% Feynman Diagram Stuff 
\def\midarrow{\draw[thick, decoration={markings, mark=at position 0.5 with {\arrow{>}}}, postaction={decorate}]}
\def\aftermidarrow{\draw[thick, decoration={markings, mark=at position 0.7 with {\arrow{>}}}, postaction={decorate}]}
\def\beforemidarrow{\draw[thick, decoration={markings, mark=at position 0.4 with {\arrow{>}}}, postaction={decorate}]}
\def\wavey{\draw[thick, decorate, decoration={snake, segment length=1.5mm, amplitude=0.5mm}]}


\begin{document}

\captionsetup[figure]{margin=1.5cm,font=small,labelfont={bf},name={Figure},labelsep=colon,textfont={it}}
\captionsetup[table]{margin=1.5cm,font=small,labelfont={bf},name={Table},labelsep=colon,textfont={it}}
\setlipsumdefault{1}

\frontmatter

\begin{titlepage}
	\centering
	    \scshape % Use small caps for all text on the title page
        \vspace*{\baselineskip} % White space at the top of the page
        
	    \rule{\textwidth}{1.6pt}\vspace*{-\baselineskip}\vspace*{2pt} % Thick horizontal rule
	    \rule{\textwidth}{0.4pt} % Thin horizontal rule
	    
	    \vspace{0.75\baselineskip} % Whitespace above the title
	    
	    {\LARGE Introduction To Field Theory} % Title
	    
	    \vspace{0.75\baselineskip} % Whitespace below the title
	    
	    \rule{\textwidth}{0.4pt}\vspace*{-\baselineskip}\vspace{3.2pt} % Thin horizontal rule
	    \rule{\textwidth}{1.6pt} % Thick horizontal rule
	    
        \vspace{5\baselineskip} % Whitespace after the title block

	    Course delivered in 2019 by 
	
	    \vspace{0.5\baselineskip} % Whitespace before the editors
	
	    {\scshape\Large Professor Michael Spannowsky} % Lecturer Name
	
	    \vspace{0.5\baselineskip} % Whitespace below the editor list
	
	    \textit{Durham University} % Lecturer Institution
	    
	    \vspace{5\baselineskip} % Whitespace after the title block
	    
	    \includegraphics[width=8cm]{images/DurhamLogo.png}\\[1cm] % Logo 
	    
	    \vspace{3\baselineskip}

	    Notes taken by 
	
	    \vspace{0.5\baselineskip} % Whitespace before my name
	
	    {\scshape\Large Richie Dadhley} % Lecturer Name
	   
	    \vspace{0.5\baselineskip} % Whitespace below my name
	    \textit{richie@dadhley.com} % Email
	
	    \vfill % Whitespace between editor names and publisher logo
\end{titlepage}



% -------------------------------------------------------------------
% Acknowledgements
% -------------------------------------------------------------------

\newpage
\section*{Acknowledgements}

These are my notes on the 2019 lecture course "Introduction to Field Theory" taught by Professor Michael Spannowsky at Durham University as part of the Particles, Strings and Cosmology Msc. For reference, the course lasted 4 weeks and was lectured over 24 hours. \\

I have tried to correct any typos and/or mistakes I think I have noticed over the course. I have also tried to include additional information that I think supports the taught material well, which sometimes has resulted in modifying the order the material was taught. Obviously, any mistakes made because of either of these points are entirely mine and should not reflect on the taught material in any way. \\

% For any other notable sources used:
I have also used/borrowed from Professor David Tong's notes throughout the course (particularly at the start). This notes are a must for any new students to QFT and can be found via 

\begin{center}
    \href{http://www.damtp.cam.ac.uk/user/tong/qft/qft.pdf}{http://www.damtp.cam.ac.uk/user/tong/qft/qft.pdf}
\end{center}

The material is presented in the order taught. This generally coincides with Professor Tong's notes, however there are small deviations here and there. Probably the biggest difference is the fact that we present the LSZ theorem in this course, whereas that is saved for the Advanced QFT course at Cambridge. Equally, normal ordering is not introduced until lecture 7, where Wick's theorem is discussed, whereas Prof. Tong introduces it much earlier (in order to drop the infinities in expectation values). 

I would like to extend a message of thanks to Professor Spannowsky for teaching this course. I would also like to thank Professor Tong for uploading his notes to the internet. \\

If you have any comments and/or questions please feel free to contact me via the email provided on the title page. \\

These notes are currently a work in progress, so for updated notes, as well as a list of other notes/works I have available, visit my blog site

\begin{center}
    \href{https://richie291.wixsite.com/theoreticalphysics}{https://richie291.wixsite.com/theoreticalphysics}
\end{center}

These notes are not endorsed by Professor Spannowsky or Durham University.

\vspace{1cm}

\begin{flushright}
    \Huge{{\cursive\setul{0.1ex}{}\ul{~~Richie Dadhley~~}}}
\end{flushright}

% -------------------------------------------------------------------
% Contents
% -------------------------------------------------------------------

\tableofcontents

% -------------------------------------------------------------------
% Main sections (as required)
% -------------------------------------------------------------------

\mainmatter

\chapter{Introduction}

Stuff to come. Thinking I might combine the notes on IFT, QFT1, QFT2, QED, QCD, SM etc into one big document. If I do do that, put an explanation here. If not just delete this chapter.
\chapter{Introduction \& Motivation}

\section{The Battle Between Quantum Mechanics \& Special Relativity}

The first question we should ask ourselves is "why does quantum field theory (QFT) even exist and what is it used for?" The answer to this question comes from the fact that special relativity and quantum mechanics (QM) don't get along, and we need some method to make them consistent. Let's see why.

\subsection{Causality}

In the pre-Einstein ages, we often just had formulae for the forces between objects and just got on with using them. For example, let's say we have two particles of charges $q_1$ and $q_2$. Let's plonk them down somewhere and label their separation by $d$. Then the rules of electromagnetism tell us that the force between these two particles is given by 
\be 
\label{eqn:EMForce}
    F_{EM} \propto \frac{q_1q_2}{d^2}.
\ee
We were happy with this, until Einstein came along and told us about a little thing called \textit{causality}. This is basically the statement that information cannot travel faster than the speed of light. In the language of special relativity that is `information cannot travel between two spacelike separated points'. This wreaks havoc for \Cref{eqn:EMForce}. Why? Well because there is no time dependence! Therefore if we were to move one of the charges, $q_1$ say, then the value of $d$ would change, and \Cref{eqn:EMForce} tells us that $q_2$ would feel this change in force \textit{instantly}. In other words, the information contained in the fact that we have moved $q_1$ is transmitted to $q_2$ (which could be arbitrarily far away) instantaneously; not good! We can fix this problem if instead we think of the force being mediated by some field that travels at the speed of light, and this is exactly what we do.

\br 
    Note that when $q_1$ and $q_2$ are close by, the fact that light travels so fast means that the difference between the result obtained by \Cref{eqn:EMForce} and the field theory description will be negligible. This is the usual story of non-relativistic vs. relativistic theories. 
\er 

Ok, that's all very well and good, but what on earth does that have to do with QM? Well in quantum mechanics we have operators which act on the states of our Hilbert space and do something to them. This is information! If we couple this idea with the fact that we want to be able to act with operators at any point in space and time\footnote{From now on I shall only say spacetime, and not `space and time' as the latter tends to promote the idea of thinking of space and time as different things.} we see that we are in a very similar situation to the electromagnetism case. 

Recall that one of the weird and wonderful things about QM is that it is not always possible to simultaneously measure two observables. Mathematically we write this as the statement that their commutator is non-vanishing:
\bse 
    [\hat{\cO}_A,\hat{\cO}_B] \neq 0
\ese 
for non-compatible operators. This essentially tells us that the action of $\cO_A$ on our state speaks to $\cO_B$ and alters its action and visa versa. However if we are going to preserve causality, it is important that two spacelike separated operators must not speak to each other. In terms of our commutator relations, this is\footnote{We are working with a field theorist signature of $(+,-,-,-)$.}
\be
\label{eqn:CausualOperators}
    [\hat{\cO}_A(x),\hat{\cO}_B(y)] = 0  \qquad \text{if } (x-y)^2<0,
\ee 
for \textit{any} two operators $\cO_A$ and $\cO_B$. To emphasise, the above condition holds even if $\cO_A$ and $\cO_B$ are non-compatible. 

The collection of observables form a quantum field in spacetime and signals propagate as disturbances in this field. We interpret these disturbances as \textit{particles}.\footnote{On a technical note, the particle interpretation in QFT is actually much more subtle. A global notion of particle requires a global notion of time. As we are working in Minkowski spacetime we of course have a global time coordinate and so this is not a problem. However for a curved spacetime things become more subtle and you get some highly non-trivial results, namely that black holes have temperature; Hawking Radiation.}

\subsection{Multiparticle States}

Another problem that arises when trying to marry QM with special relativity is the conservation of particle number. In regular quantum mechanics, we work hard to define our Hilbert space\footnote{For \textit{a lot} more info on this see Simon Rea and my joint \href{https://richie291.wixsite.com/theoreticalphysics/post/dr-frederic-schuller-s-course-of-quantum-theory}{notes on Dr. Frederic Schuller's Quantum Theory course}.} and then we simply act on it. The game is a bit different in QFT, though. 

Recall that in QM we have the famous \textit{Heisenberg uncertainty relation}, which tells us that you cannot know both the position and momentum to arbitrary accuracy, and that the lower bound is 
\be 
\label{eqn:UncertaintyRelation}
    \Delta L \Delta \mathbf{p} \geq \frac{\hbar}{2} \qquad \implies \qquad \Delta \mathbf{p} \geq \frac{\hbar}{2\Delta L}.
\ee 
This tells us that if we constrain a particle to a small region of space (say by putting it into a really small box) then the uncertainty in momentum goes through the roof. 

Recall separately that special relativity tells us that energy and mass are essentially the same thing $E=mc^2$, and that in a general frame the energy of a particle is given by its mass and 3-momentum:
\bse 
    E^2 = m^2c^4 + \mathbf{p}^2 c^2.
\ese 
If we consider a very light particle, this relation can be approximated as 
\be
\label{eqn:E=pc}
    E \approx |\mathbf{p}|c
\ee

Now we have a problem! Imagine constraining the particle into a box of size $L = \hbar/2mc$ then \Cref{eqn:UncertaintyRelation} tells us that the uncertainty in the momentum is 
\bse 
    \Delta \mathbf{p} \geq mc,
\ese 
and then \Cref{eqn:E=pc} tells us that the uncertainty in energy is 
\bse 
    \Delta E \geq mc^2.
\ese 
This is enough energy to create a new particle!

So we see that in QFT we need to adapt our Hilbert space so that it accounts for these multiparticle states. This construction is known as a Fock space, and mathematically can be written as 
\be 
\label{eqn:FockSpace}
    \cF = \bigoplus_{n=0}^{\infty} \cH^{\otimes n},
\ee
where $\cH$ is our `one particle state'. To make the above seem a little less abstract, the two particle state is an element of the space $\cH^{\otimes 2} := \cH \otimes \cH$, and so each $n$ above corresponds to allowing $n$-particle states. The $\oplus$ symbol just means we need to `sum over' these different combinations to account for the fact that our system can transition from a $n$-particle state to a $(n+1)$-particle state.\footnote{For more info on the symbols used, see either the QM notes mentioned in footnote 4 or any decent linear algebra textbook.}

\br 
    Anyone familiar with the basic idea of QFT will know that what we never really just create a single particle in this way, but instead they always appear is so-called \textit{pair production}. This is essentially just the statement that a particle is always produced with its antiparticle, but we'll return to this later. 
\er 

QFT will allow us to study how such particle creation processes occur and interpret them in terms of interacting field. Being a quantum theory, what we will obtain is the \textit{probability} that such a process occurs, known as \textit{cross sections}. 

\subsection{Fields, Not Particles}

The above argument should actually give us a reasonably intuitive idea to why we want to consider fields instead of particles. We have just seen that even in the simple case of Minkowski spacetime where the notion of a particle makes sense (see footnote 3), the notion of how many particles we have is very blurred. This is not something that lends well to our understanding, and so we would really like to think of things differently, and this is where the fields come in. As mentioned above, the fields are the fundamental objects and we obtain a particle interpretation by the disturbances in the field. To quote Weinberg:

\begin{center}
    "Quantum fields are the basic ingredients of the Universe, particles are just bundles of energy and momentum made out of them."
\end{center}

So QFT theory is, as the name suggests, a study of \textit{fields}, not particles. This might seem like a pedantic point to make, however this switch in mentality is actually very important for a solid understanding of the material that is about to come.

\section{Units \& Range}

\subsection{Units}

Experienced physicists will know that physical units (or dimensions) are important, and must be kept track of. As you may be aware, it turns out that \textit{all} physical units can be characterised by seven others, known as SI units. The three most commonly encountered (and the only ones used in this course) are mass, length and time. We write the dimensions of a quantity by square brackets and use the letters $[M]$, $[L]$ and $[T]$ for the three units, respectively.

The constants that appear in our equations will, of course, sometimes carry units. The three dimensionful constants in particle physics are: the reduced Planck's constant, $\hbar$; the speed of light $c$; and Newton's gravitational constant $G_N$. These have dimensions
\bse 
    [c] = \frac{[L]}{[T]}, \qquad [\hbar] = \frac{[M][L]^2}{[T]}, \qand [G_N] = \frac{[L]^3}{[M][T]^2}.
\ese 
It is the former two that will show up a lot (as they already have above) and particle physicists decided one day they had had enough of writing them out and so decided to switch to a different set of units, known as \textit{natural units}. These are defined by setting $\hbar=1=c$. From the above relation this results in 
\bse 
    [L] = [T] = [M]^{-1}.
\ese 
We can therefore just characterise everything in terms of one of these three units. It is standard to give this privilege to mass, and we obtain \textit{mass dimensions}. These are just numbers, and correspond to the power of $[M]$ appearing. For example 
\bse 
    [G_N] = [M]^{-3} [M]^{-1} [M]^{2} = -2.
\ese 

Now recall \Cref{eqn:UncertaintyRelation}. This tells us that position (which has length dimension) has the opposite dimension to 3-momentum. Then \Cref{eqn:E=pc} tells us that energy has the same dimension as 3-momentum. We therefore have 
\be 
    [E] = 1.
\ee 
This result is also trivially obtained from $E=m$ (as $c=1$). We can therefore express these dimensions in units of energy. Due to the energy scales that appear in particle physics, we do not choose to use the Joule but instead use the much more convenient electron volt, eV. These often appear with the prefixes M, G, T for `mega', `tera' and `giga', respectively.

\bex 
    As an example, the mass of an electron in these units is 
    \bse 
        m_{e} = 0.511 MeV.
    \ese
\eex 

\subsection{Range}

Perhaps an obvious question to now ask is "what energy scale does QFT apply to?" 

%\textcolor{red}{Fill this in later. Maybe just reference Tong.}

\section{Classical Particle Mechanics}

In order to make the discussion of QFT easier, let's first look at classical field theory. 

\subsection{Configuration Space \& Generalised Coordinates}

We can describe the configuration of a system by a set of parameters. For example, the position of a particle on $\R^2$ is simply given by its coordinates $(x,y)$, say. Perhaps a less obvious example is the colour of something. This is described by its RGB value, and so is parameterised by 3 numbers. We can think of these parameters as coordinates for our system. This coordinate space need not be continuous, as it would be for both the above examples, but could form a lattice. For example if we whether something is `right way up' or 'upside down' would form two points on a configuration space. 

The question is, do constraints on the system effect the parameters, and therefore effect the coordinate systems? The answer is yes. To see why lets contrast a free particle with one constrained to some surface, a circle, say. The free particle can have any position and so we need to consider the entirety of $\R^2$, that is it can have any $(x,y)$ values. The constrained particle, however, can only take values such that $x^2+y^2=r^2$. If instead we used polar coordinates, we could completely specify the particle's position just by its angle around the circle. So this constraint has reduced the number of parameters needed, and thus altered our configuration space. A similar example for the colour of an object would be if we fixed the R value but allowed G and B to vary still. 

Note all of these constructions were made relative to some `reference point': for the particle's position we reference the values to the origin; for the colours we reference it to white, which has RGB=$000$.

\bd[Generalised Coordinates]
    We call the minimum number of independent parameters needed to describe a system relative to some reference point the \textbf{generalised coordinates}. We often denote these with a $\{q_a\}$.
\ed 

We use the name `generalised' coordinates, because, as we have seen, we don't need the same number of generalised coordinates as we do `real' coordinates (i.e. those used to define a coordinate system). Note that the configuration space is given by the span of the generalised coordinates. 

\bex 
    The generalised coordinates for the free particle are $(x,y)$, whereas for the constrained particle we just have one, $\theta$.
\eex 

\bex 
    The values $(q_1,q_2,q_3)$ could be the generalised coordinates for either a single particle in 3-dimensions or for 3 particles in 1-dimension. This example highlights that in order to differentiate two systems, we need to know how the configuration space is constructed.
\eex 

\subsection{Principle Of Least Action \& Euler-Lagrange}

The configuration space tells us the configuration of a system, but in physics we also want to know about the \textit{dynamics} of a system. That is, we want to know how the configuration varies through time. For the case of a particle's position, this corresponds to knowing its momentum. In the language of generalised coordinates, we call these momenta the \textit{generalised velocities}. We often denote these by $\{\dot{q}_a\}$. The span of $\{q_a\}\cup\{\dot{q}_a\}$ is called the \textit{phase space}. Points in the phase space correspond to \textit{states} of the system, and the evolution of states are given by paths on phase space.

In classical physics, the dynamics of a particle is given by a phase space where $q_a$ and $p_i$ are the position and momentum of the particle. For non-accelerating particles, we know that we can simply specify the position and momentum of a particle at some point and from that extract out its evolution. In other words, each point in phase space only has one path going through it, and so by specifying the state of a system we know its evolution. 

The question is "how do we find this evolution?" The answer comes from the \textit{principle of least action}. These tell us the equations of motion for a system (i.e. how it evolves through time) and so will give us exactly what we want.

The procedure is as follows: we obtain the \textit{Lagrangian} $L$ for our system in terms of $q$ and $\dot{q}$, integrate it from initial to final time to obtain the \textit{action}, then we extremise this action under the constraint that the start and end points remain fixed, and obtain the shortest path between those points. The calculation is as follows: 

\bse 
    \begin{split}
        \del S & = \int_{t_i}^{t_f} dt \del L\big(q_a(t),\dot{q}_a(t)\big) \\
        & = \int_{t_i}^{t_f} dt \bigg( \frac{\p L}{\p q_a} \del q_a + \frac{\p L}{\p \dot{q}_a} \del \dot{q}_a \bigg) \\
        & = \int_{t_i}^{t_f} dt \bigg[ \frac{\p L}{\p q_a} - \frac{d}{dt}\bigg(\frac{\p L}{\p \dot{q}_a}\bigg)\bigg] \del q_a + \bigg[\frac{\p L}{\p \dot{q}_a} \del q_a\bigg]^{t_f}_{t_i} \\
        & = \int_{t_i}^{t_f} dt \bigg[ \frac{\p L}{\p q_a} + \frac{d}{dt}\bigg(\frac{\p L}{\p \dot{q}_a}\bigg)\bigg] \del q_a,
    \end{split}
\ese 
where we have used 
\bse 
    \frac{d}{dt}\bigg(\frac{\p L}{\p \dot{q}_a}\del q_a \bigg) = \frac{\p L}{\p \dot{q}_a} \del \dot{q}_a + \frac{d}{dt}\bigg(\frac{\p L}{\p \dot{q}_a}\bigg)\del q_a,
\ese 
integrated by parts, and used the fact that $\del q_a =0$ at $t_i$ and $t_f$. Now setting the variation to zero, and using the fact that $\del q_a$ was arbitrary (up the constraints) we can conclude 
\be 
\label{eqn:EulerLagrange}
    \frac{\p L}{\p q_a} - \frac{d}{dt}\bigg(\frac{\p L}{\p \dot{q}_a}\bigg) = 0,
\ee 
known as the \textit{Euler-Lagrange equations}. These are exactly the equations of motion of the system we set out to obtain. 

Practically speaking, the Lagrangian is given by the kinetic energy, $T$, minus the potential energy, $U$. We shall now give a few examples to show that \Cref{eqn:EulerLagrange} does indeed give us the expected equations of motion.

\bex 
    Consider the motion of single particle in 1-dimension with a potential $U(x)$. The Lagrangian is 
    \bse 
        L(x,\dot{x}) = \frac{1}{2}m \dot{x}^2 - U(x),
    \ese 
    and so out Euler-Lagrange equations are:
    \bse 
        m\ddot{x} = -U'(x),
    \ese 
    which is exactly Newton's law of motion.
\eex 

\section{Lagrangian Field Theory}

The above example was to do with a single \textit{particle}, but what we're trying to study here are \textit{field} theories. We therefore need to adapt the above procedure to this case. 

The first change we need to make is to the generalised coordinates. For the particle case, $a$ was an index which labelled the different $q(t)$s. However in field theory, we must also introduce a label for the position spatial of the field, and so our generalised coordinates are replaced by 
\bse 
    q_a(t) \longrightarrow \phi(x) := \phi_a(\vec{x},t),
\ese 
where $\vec{x}$ is a treated as a label. Note that for the field, we are dealing with a system with an infinite number of degrees of freedom: one for each value of $\vec{x}$.

We can also show by considering a chain of connected springs \textcolor{red}{(Do this later from his notes)} that we must also replace the Lagrangian with a Lagrangian \textit{density}, $\cL$, related by 
\be 
\label{eqn:LtocL}
    L(t) = \int d^3\vec{x} \, \cL\big(\phi_a(x),\p_{\mu}\phi_a(x)\big).
\ee 
Our Euler-Lagrange equations then become
\mybox{ 
    \be 
    \label{eqn:EulerLagrangeDensity}
        \frac{\p \cL}{\p \phi_a} - \p_{\mu}\bigg(\frac{\p \cL}{\p(\p_{\mu}\phi_a)}\bigg) = 0
    \ee 
}

The symmetries of the physics are inserted at the Lagrangian level, by which we mean if we want a theory that is spatial-rotationally invariant then we require 
\bse 
    R(\varphi_x,\varphi_y,\varphi_z) \cL = \cL,
\ese
where $R(\varphi_x,\varphi_y,\varphi_z)$ is a general spatial rotation. 

In this course we will only consider Lagrangians that obey the following conditions
\ben[label=(\roman*)]
    \item Lorentz invariance; i.e. $\cL(\Lambda x) = \cL(x)$ where $\Lambda$ is a general Lorentz transformation,
    \item Local Lagrangians; i.e. all of our fields will be evaluated at the same event in spacetime,
    \item Only contain single derivatives of the fields; i.e. we don't want anything of the form $\p_{\mu}\p_{\nu}\phi$, 
    \item Renormalisable Lagrangians; this is something that will not really be covered explained in this course, but basically the coefficients in our Lagrangian must have positive dimension. This essentially amounts to not having more than $d$ powers of the field, where $d$ is the dimension of the spacetime. 
\een 
These four conditions will actually massively reduce the number of possible Lagrangians we can consider. We now give a brief justification for each: 
\ben[label=(\roman*)]
    \item We are considering relativistic particles and so we really want our theories to be Lorentz invariant, 
    \item This massively helps with avoiding non-causal theories, 
    \item There are two reasons: one is that we would need to alter our Euler-Lagrange equations to account for higher derivatives, and also adapt our phase space\footnote{Basically you get equations of motion that contain second order derivatives, and so it is not enough to just specify $q$ and $\dot{q}$ on an initial time surface, but we must also specify $\ddot{q}$.}; secondly because we are fundamentally interested in 4-dimensional spacetime, a Lagrangian that included a $\p_{\mu}\p_{\nu}\phi$ term would be non-renormalisable. This comes from the fact that the action has vanishing mass dimension\footnote{Basically because $S \sim \hbar$.} and so, in order to cancel the $[d^dx]=-d$\footnote{Note that integration increases the dimension and so has units of length.} contribution from the integral we require $[\cL]=d$. Now if we want to study kinematics, we need a kinetic energy term in the Lagrangian, this takes the form
    \bse 
        \cL = (\p_{\mu}\phi)(\p^{\mu}\phi) + ...,
    \ese
    but the derivatives have dimension $[\p_{\mu}]=1$\footnote{They take away length dimension.} and so we get $4 = 2 + 2[\phi]$, or $[\phi]=1$. If we were then to include a second order derivative term, Lorentz invariance would require it appears as
    \bse 
        \cL = a(\p_{\mu}\p_{\nu}\phi)(\p^{\mu}\p^{\nu}\phi),
    \ese 
    but this would then require $[a]=-2$, which gives a non-renormalisable theory. 
    \item We don't need to understand this here, but basically it causes a huge pain in perturbation theory. 
\een 

\bnn 
    As we have just seen in the explanation of (iii) above, Lorentz invariance requires that any derivative term comes with the corresponding `raised' derivative term. It is therefore very useful to introduce a common notion to simply the expressions. We define 
    \bse 
        (\p_{\mu}\phi)^2 := \p_{\mu}\phi\p^{\mu}\phi = \eta^{\mu\nu}\p_{\mu}\phi\p_{\nu}\phi. 
    \ese 
    This notation can appear a bit confusion because it appears to break Einstein summation convention (the left-hand side seems to have a $\mu$ index but on the right-hand side it is summed over). For this reason some people (and I am one of these people) prefer to define the notation
    \be 
    \label{eqn:PartialSquaredNotation}
        (\p\phi)^2 := \p_{\mu}\phi\p^{\mu}\phi.
    \ee
    This is the notation that will be used in these notes from now on, unless some potential confusion might arise, in which case we shall write the full expression (i.e. like the right-hand side above). Similarly we shall use 
    \be
    \label{eqn:ClassicalKleinGordan}
        \p^2 := \p_{\mu}\p^{\mu}.
    \ee 
\enn 

\bex 
    One of the simplest Lagrangians that meets the above conditions is 
    \be 
    \label{eqn:KleinGordanLagrangian}
        \begin{split}
            \cL & = \frac{1}{2}(\p\phi)^2 - \frac{1}{2} m^2\phi^2 \\
            & = \frac{1}{2} \dot{\phi}^2 - \frac{1}{2}\nabla^2\phi - \frac{1}{2}m^2\phi^2,
        \end{split}
    \ee 
    where $\phi = \phi(x)$ is a real scalar field. The Euler-Lagrange equations for this are simply 
    \bse 
        (\p^2 - m^2)\phi = 0.
    \ese 
    This is known as the \textit{Klein-Gordan equation}. It is very important to note here that this is the \textit{classical} Klein-Gordan equation, not the quantum version. As we will see later it has an quantum analogue (which looks exactly the same) which gives the equations of motion of spineless particles, i.e. quantum scalars.
\eex 

\bex 
    The Klein-Gordan equation above takes the form of a wave equation (i.e. it is second order in both temporal and spatial derivatives). This is nice,but we know in QM that the wave equation is replaced by the Schr\"{o}dinger equation, which is linear in time derivatives and quadratic is spatial derivatives. We can get a \textit{classical} analogue to this by considering a complex scalar field $\psi$. The Lagrangian is
    \bse 
        \cL = \frac{i}{2}(\psi^*\dot{\psi} -\dot{\psi}^*\psi) - \nabla\psi^*\cdot\nabla\psi - m\psi^*\psi.
    \ese 
    You treat $\psi$ and $\psi^*$ as independant fields. For the $\psi^*$ field we have,
    \bse 
        \frac{\p\cL}{\p\psi^*} = \frac{i}{2}\dot{\psi} - m\psi, \qquad \frac{\p\cL}{\p\dot{\psi}^*} = -\frac{i}{2}\psi, \qand \frac{\p\cL}{\p\nabla\psi^*} = -\nabla\psi,
    \ese 
    and so the equation of motion for this field is 
    \bse 
        i\frac{\p\psi}{\p t} = -\nabla\psi + m\psi.
    \ese 
    The $\psi$ field gives a similar equation of motion. 
    
    This \textit{looks} a lot like the Schr\"{o}dinger equation, but as was emphasised above this is a \textit{classical} theory. It does not have any of properties of the Schr\"{o}dinger equation, e.g. it doesn't have a probabilistic interpretation, as is the key to quantum theory. 
\eex 
\chapter{Symmetries \& Hamiltonian Field Theory}

\section{Lorentz Invariance}

As we explained last lecture, we want to consider theories that are Lorentz invariant. Will therefore be testing for this often and so we want to clarify the notion and conventions here. 

As in (i) last lecture, we denote a general Lorentz transformation via the symbol $\Lambda$. We think of Lorentz transformations as acting on the spacetime coordinates $\{x^{\mu}\}$\footnote{In fact this interpretation is a bit misleading, but that's not important here. For an explanation of what I mean, see my \href{https://richie291.wixsite.com/theoreticalphysics/post/the-we-heraeus-international-winter-school-on-gravity-and-light}{notes on Dr. Schuller's Winter School on Gravity and Light}, Section 13.3.} and so it is useful to express them as matrices. We therefore have 
\bse 
    x^{\mu} \longrightarrow \widetilde{x}^{\mu} = {\Lambda^{\mu}}_{\nu} x^{\nu}
\ese 
as our transformation. We also have the condition 
\be 
\label{eqn:LorentzTransformationsOnMetric}
    {\Lambda^{\mu}}_{\tau} \eta^{\tau\sig} {\Lambda^{\nu}}_{\sig} = \eta^{\mu\nu}. 
\ee 

Lorentz transformations are essentially spatial rotations and boosts, and we can work out the entries of the matrix ${\Lambda^{\mu}}_{\nu}$ by consider the action on the $\{x^{\mu}\}$. We give an example of each below. 

\bex
    Let's set $(x^0,x^1,x^2,x^3)=(t,x,y,z)$. First let's consider a rotation around the $z$ axis. This does nothing to the $t$ axis and obviously doesn't effect $z$. We therefore just set the SO(2) rotation in the $xy$-plane by angle $\theta$. In other words we have 
    \bse 
        \begin{pmatrix}
            t \\
            x \\
            y \\
            z 
        \end{pmatrix} \longrightarrow \begin{pmatrix}
            t \\
            x\cos\theta - y\sin\theta \\
            y\sin\theta + x\cos\theta \\
            z 
        \end{pmatrix},
    \ese 
    and so we can conclude\footnote{In these notes I shall try and be consistent and put brackets around indexed objects to indicate that we mean the matrix.}
    \bse 
        ({\Lambda^{\mu}}_{\nu}) = \begin{pmatrix}
            1 & 0 & 0 & 0 \\
            0 & \cos\theta & -\sin\theta & 0 \\
            0 & \sin\theta & \cos\theta & 0 \\
            0 & 0 & 0 & 1
        \end{pmatrix}.
    \ese 
\eex

\bex 
    Following a similar idea to the previous example, we can show that for a boost along the $x$ axis we have 
    \bse 
        ({\Lambda^{\mu}}_{\nu}) = \begin{pmatrix}
            \g & -\g v & 0 & 0 \\
            -\g v & \g & 0 & 0 \\
            0 & 0 & 1 & 0 \\
            0 & 0 & 0 & 1
        \end{pmatrix},
    \ese 
    where $\g = 1/\sqrt{1-v^2}$.
\eex 

\subsection{Active vs. Passive Transformations}

There is a subtlety in taking transformations that can be wonderfully confusing the first time you hear it, and this is the difference between an active and passive transformation. This short section just aims to clear up the differences so that the notation that follows in these notes is not confusing. 

Let's say we want shift the temperature field in a room to the right. That is the temperature your left-hand neighbour is currently feeling you will feel after the transformation. There are essentially two ways to achieve this: firstly we could actually move the air particles themselves and so physically move the temperature to the right; secondly we could leave the air where it is and instead we, the people,\footnote{Not to be confused with \href{https://en.wikipedia.org/wiki/We_the_People_(disambiguation)}{this}.} could move to the left. In both cases you will end up feeling the temperature your left-hand neighbour previously felt, the question is "what's the difference?" 

The answer is that the former is an \textit{active} transformation whereas the latter is a \textit{passive} one. Put in a more mathematically meaningful way, an active transformation is one where the thing itself (in this case the temperature field) is moved, whereas a passive transformation is one where we simply shift the underlying reference frame/coordinate system (in this case, us). 

Those passionate about the concepts of relativity should now throw their arms up in protest. Why? Well because a passive transformation clearly depends on the choice of coordinate system and so should not be a physical thing. In contrast, the active transformation makes perfect sense without reference to any coordinate system. Put in terms of the example above, the passive transformation only makes sense if we are in the room and move ourselves and explain the shift, whereas the active transformation makes perfect sense even if we're not in the room at all. 

It is for this reason that we actually consider active transformations in field theory. That is we actually consider shifting the \textit{fields} in our problems and not the coordinate system. As we have just seen, when we \textit{measure} something (i.e. we personally measure the change in temperature), we can think of the active transformation as a passive one going in the other direction (i.e. we move left so that `the temperature moves right'). This is why in what is to follow we shall write the transformation on fields as follows
\be 
\label{eqn:ActiveFieldTransformation}
    \phi(x) \longrightarrow \widetilde{\phi}(x) = \phi\big(\Lambda^{-1}x\big).
\ee 

The above equation is actually a particular case for the action of the \textit{representation} of the Lorentz transformation on fields. In this case we are acting on a scalar field and so no $D[\Lambda]$ factors appear. For a general field we have 
\bse 
    \phi^{\mu\nu...}(x) \longrightarrow D[\Lambda]^{\mu}_{\sig} \big(D[\Lambda]^{\nu}_{\tau} ...\big) \phi^{\sig\tau...}\big(\Lambda^{-1}x\big).
\ese 

\section{Noether's Theorem}

\bt[Noether's Theorem]
\label{thrm:Noether}
    Every continuous symmetry of the Lagrangian gives rise to a conserved current, which we label $j^{\mu}(x)$, such that the equations of motion imply 
    \be 
    \label{eqn:pj=0}
        \p_{\mu}j^{\mu} = 0.
    \ee 
\et 

As we will see, symmetries and their conserved currents are incredibly important in field theory, and so Noether's theorem is an invaluable tool to field theorists. 

\bq 
    In order to prove Noether's theorem it is very helpful to work infinitesimally. As we are considering continuous symmetries, we can always do this. 
    
    Now the first thing we note is that the equations of motion, \Cref{eqn:EulerLagrangeDensity}, were obtained by varying the action and dropping surface terms. We therefore see that these will be completely unaffected for any transformation that takes the form 
    \be 
    \label{eqn:LagrangianBoundaryChange}
        \cL \longrightarrow \cL + \epsilon\p_{\mu}F^{\mu},
    \ee 
    where $F^{\mu}$ is some function of the fields and $\epsilon$ is some small parameter.
    
    Ok, so let's consider some infinitesimal transformation of the field
    \bse 
        \phi_a(x) \longrightarrow \phi_a(x) + \epsilon\del\phi_a(x).
    \ese 
    This gives a change in the action, of the general form\footnote{We use the fact that $\epsilon$ is a constant to take it outside the derivative.} 
    \bse 
        \begin{split}
            \epsilon\del\cL & = \frac{\p\cL}{\p\phi_a} \epsilon\del\phi_a + \frac{\p\cL}{\p(\p_{\mu}\phi_a)}\p_{\mu}(\epsilon\phi_a) \\
            & = \bigg[ \frac{\p\cL}{\p \phi_a} -\p_{\mu}\bigg(\frac{\p\cL}{\p(\p_{\mu}\phi_a)}\bigg)\bigg] \epsilon\del\phi_a + \epsilon\p_{\mu} \bigg(\frac{\p\cL}{\p(\p_{\mu}\phi_a)}\del\phi_a\bigg) \\
            & = \epsilon\p_{\mu} \bigg(\frac{\p\cL}{\p(\p_{\mu}\phi_a)}\del\phi_a\bigg),
        \end{split}
    \ese 
    where we have used \Cref{eqn:EulerLagrangeDensity} to set the square bracket term to zero. Now if this transformation is a symmetry of the system then we also have 
    \bse 
        \epsilon\del\cL = \p_{\mu}F^{\mu},
    \ese
    and so taking these two results away from each other gives us 
    \mybox{
        \be 
        \label{eqn:NoetherCurrent}
            j^{\mu} = \frac{\p\cL}{\p(\p_{\mu}\phi_a)}\del\phi_a - F^{\mu}(\phi_a) \qquad \implies \qquad \p_{\mu}j^{\mu} = 0.
        \ee 
    }
    So this is our conserved current. We get one of these for each continuous $\del\phi$ transformation and so we have proved Noether's theorem. 
\eq 

Before doing some examples, first there is an important comment to make. As we have been careful to say above, Noether's theorem tells us that we get a conserved \textit{current}. This is a stronger result then saying we get a conserved \textit{charge}. Indeed a conserved current gives rise to a conserved charge, and we get this charge as follows: 
\bse 
    \p_{\mu}j^{\mu} = 0 \qquad \implies \qquad \frac{\p j^0}{\p t} = - \nabla \cdot \Vec{j}.
\ese 
Integrating both sides of this over all of space, we have 
\bse 
    \begin{split}
        \frac{\p}{\p t} \int_{\R^3} d^3 \vec{x} \, j^0 & = -\int_{\R^3}d^3\vec{x} \, \nabla\cdot\vec{j} \\
        & = - \oint d\vec{s} \cdot \vec{j} \\
        & = 0,
    \end{split}
\ese 
where the last line follows from the fact that we assume $\vec{j}$ falls off sufficiently fast as $|\vec{x}|\to\infty$. We can therefore define our conserved charge by 
\be 
\label{eqn:NoetherCharge}
    Q := \int_{\R^3} d^3 \vec{x} \, j^0.
\ee 
This is a \textit{globally} conserved charge. 

The reason that the existence of a conserved current is stronger than the existence of a globally conserved charge is that the current tells us that charge is conserved \textit{locally}. This is easily seen by repeating the above procedure but now just integrating over some finite volume, $V$. We obtain 
\bse 
    \frac{dQ_V}{dt} = - \oint_{\p V} \vec{j}\cdot d\vec{S},
\ese 
where $\p V$ is the boundary of $V$. This equation tells us the the charge leaving a volume $V$ in time $t$ (left-hand side) is equal to the current flowing through the boundary of the volume (right-hand side). 

This is a much more powerful statement than simply that charge is conserved globally. For further clarity, if charge is conserved locally everywhere then we know that it is conserved globally (simply take your local region to be infinitely big). However if we only know that charge is conserved globally it could be possible that a charge disappears at some point $x\in\R^3$ and miraculously reappears at some other point $y\in\R^3$ far away from $x$. In other words the charge `teleported' through space. 

\subsection{The Energy-Momentum Tensor}

An important example of Noether's theorem comes from considering translations in spacetime. Recall that in particle mechanics, translations in space gave rise to momentum conservation, whereas translations in time gave rise to energy conservation. We are now dealing with a relativistic theory and so want to avoid this splitting of spacetime as much as possible, and so we want to find a 4-dimensional generalisation of the above conservation laws. This is exactly the \textit{energy-momentum tensor}. We denote it by a $T$. 

\br 
    Some people also refer to the energy-momentum tensor as the stress-energy tensor or the stress-momentum tensor, or even the stress-energy-momentum tensor. I will try to use energy-momentum tensor everywhere here, but I might sometimes just call it the stress tensor. Apologies for this in advance.
\er 

Ok, let's derive this lovely chappy.

Consider a continuous translation of the spacetime coordinates. As above, we shall work infinitesimally, and so we have 
\bse 
    x^{\mu} \longrightarrow x^{\mu} - \epsilon^{\mu},
\ese 
where we use a minus sign so that the transformation of the field is positive, see the discussion of active vs. passive transformations above. We can Taylor expand the field to obtain the transformation
\bse 
    \phi_a(x) \longrightarrow \phi_a(x) + \epsilon^{\mu}\p_{\mu}\phi_a(x),
\ese 
and similarly the Lagrangian transforms as 
\bse 
    \cL(x) \longrightarrow \cL(x) + \epsilon^{\mu}\p_{\mu}\cL(x).
\ese 
Comparing this to \Cref{eqn:LagrangianBoundaryChange}, we see that $F^{\mu}=\epsilon^{\mu}\cL$. Our conserved current is therefore 
\bse
    \begin{split}
        j^{\mu} & = \frac{\p \cL}{\p(\p_{\mu}\phi_a)} \epsilon^{\nu}\p_{\nu}\phi_a - \epsilon^{\mu}\cL \\
        & = \epsilon^{\nu}\bigg(\frac{\p \cL}{\p(\p_{\mu}\phi_a)} \p_{\nu}\phi_a - \del^{\mu}_{\nu}\cL\bigg),
    \end{split}
\ese 
and so we have a conserved current for each $\epsilon^{\nu}$, i.e. for the translations in each direction. These are the components of our energy-momentum tensor.
\mybox{
\be
\label{eqn:EnergyMomentumTensor}
    {T^{\mu}}_{\nu} := (j^{\mu})_{\nu} = \frac{\p \cL}{\p(\p_{\mu}\phi_a)} \p_{\nu}\phi_a - \del^{\mu}_{\nu}\cL.
\ee 
} 
Being a conserved current it satisfies 
\be 
\label{eqn:EnergyMomentumConservation}
    \p_{\mu}{T^{\mu}}_{\nu} = 0 \qquad \forall \nu = 0,1,...,d.
\ee 

In spacetime (i.e. $d=4$) we have four currents and so we have 4 charges. The $\nu$ index above tells us which direction we translated in, and so in correspondance with the comment made at the start of this section we want $\nu=0$ to correspond to energy conservation and $\nu=1,2,3$ to correspond to momentum conservation. We therefore define\footnote{Note we have raised the $\nu$ index here. As we are working in Minkowksi spacetime this really isn't a big deal, however as these are definitions you should be careful in more general spacetimes as some non-trivial factors will appear.}
\be 
\label{eqn:EnergyAndMomentumCharges}
    E := \int d^3\vec{x} \, T^{00}, \qand P^i := \int d^3\, \vec{x}T^{0i}
\ee 
to be the total energy and the ($i$-th component of the) total momentum of the field configuration, respectively. 

\br 
    Note in deriving \Cref{eqn:EnergyMomentumTensor}, we didn't say anything about the actual form of our Lagrangian, and so this result holds for generic $\cL$.
\er 

\subsection{Angular Momentum \& Boost Symmetries}

So we have done spacetime translations, but recall that we have also restricted ourselves to Lagrangians that are Lorentz invariant. Given that Lorentz transformations are continuous, Noether's theorem tells us that we should have some conserved currents to go along with them. Keeping in line with the logic applied to the translations, what kind of conservation do we expect? Well in particle mechanics spatial rotations give rise to conservation of angular momentum, so we expect some generalisation of this. As we are considering the whole set of Lorentz transformations, we will also derive the symmetries corresponding to boosts, whatever these may be. 

As always, we want to work infinitesimally, and its a fact that infinitesimal Lorentz transformations can be written as
\bse 
    {\Lambda^{\mu}}_{\nu} = {\del^{\mu}}_{\nu} + {\omega^{\mu}}_{\nu},
\ese 
for some infintesimal ${\omega^{\mu}}_{\nu}$, and where \Cref{eqn:LorentzTransformationsOnMetric} tells us 
\be 
\label{eqn:AntisymmetricOmega}
    \omega^{\mu\nu} = - \omega^{\nu\mu}.
\ee 

From \Cref{eqn:ActiveFieldTransformation}, we have, after Taylor expanding
\bse 
    \phi(x) \longrightarrow \phi(x^{\mu}) - {\omega^{\mu}}_{\nu} x^{\nu} \p_{\mu}\phi(x) \qquad \implies \qquad \del\phi = -{\omega^{\mu}}_{\nu}x^{\nu}\p_{\mu}\phi.
\ese 
Similarly we obtain 
\bse 
    \del\cL = -{\omega^{\mu}}_{\nu} x^{\nu} \p_{\mu}\cL.
\ese 
Now we use a clever trick. Firstly we note that ${\omega^{\mu}}_{\nu}$ is a constant so we can take it inside the derivative. Next we note that 
\bse 
    \p_{\mu}x^{\nu} = \del^{\nu}_{\mu},
\ese 
along with ${\omega^{\mu}}_{\mu}=0$ by \Cref{eqn:AntisymmetricOmega}. We can therefore also take the $x^{\nu}$ inside the derivative and obtain 
\bse 
    \del\cL = -\p_{\mu}\big( {\omega^{\mu}}_{\nu} x^{\nu}\cL\big).
\ese 
This is now of the form \Cref{eqn:LagrangianBoundaryChange}, with $F^{\mu} = {\omega^{\mu}}_{\nu} x^{\nu}\cL$. So our conserved currents are 
\bse
    \begin{split}
        j^{\mu} & = - \frac{\p \cL}{\p(\p_{\mu}\phi)}{\omega^{\sig}}_{\nu}x^{\nu}\p_{\sig}\phi + {\omega^{\mu}}_{\nu} x^{\nu}\cL \\
        & = - {\omega^{\sig}}_{\nu} \bigg[ \frac{\p \cL}{\p(\p_{\mu}\phi)}\p_{\sig}\phi + {\del^{\mu}}_{\sig} \cL \bigg]x^{\nu} \\
        & = - {\omega^{\sig}}_{\nu} {T^{\mu}}_{\sig} x^{\nu}.
    \end{split}
\ese 

As above we can split this into the individual currents, one for each ${\omega^{\sig}}_{\nu}$, and obtain 
\be 
\label{eqn:LorentzCurrents}
    (\cJ^{\mu})^{\rho\sig} = x^{\rho}T^{\mu\sig} - x^{\sig}T^{\mu\rho}, \qquad \text{with} \qquad \p_{\mu}(\cJ^{\mu})^{\rho\sig} = 0 \qquad \forall \rho,\sig=0,..,3.
\ee 
Each one of these has a corresponding conserved charge. We get the particle result, i.e. angular momentum, when $\rho,\sig=1,2,3$: 
\bse 
    Q^{ij} = \int d^3\vec{x} \, \big(x^iT^{0j} - x^jT^{0i}\big).
\ese 
The question is "what charge do the boosts give us?" Well these correspond to $\rho=0, \sig=1,2,3$, which give 
\bse 
    Q^{0i} = \int d^3\vec{x} \, \big(x^0T^{0i} - x^iT^{00}\big).
\ese 
What is this? Well we know that its temporal derivative vanishes, so using $x^0=t$, we have 
\bse 
    \begin{split}
        \frac{d Q^{0i}}{dt} & = \int d^3\vec{x} \, T^{0i} + \int d^3\vec{x} \, t \frac{\p T^{0i}}{\p t} - \frac{d}{dt}\int d^3\vec{x} \, x^i T^{00} \\
        0 & = P^i + \frac{d P^i}{dt} - \frac{d}{dt}\int d^3\vec{x} \, x^i T^{00},
    \end{split}
\ese 
where we've used \Cref{eqn:EnergyAndMomentumCharges}. But we've already seen that $P^i$ is a constant and so we simply get 
\bse 
    \frac{d}{dt}\int d^3\vec{x} \, x^i T^{00} = \text{constant},
\ese 
which is the statement that the center of energy of the field travels with constant velocity. 

\br 
    Note that \Cref{eqn:AntisymmetricOmega} agrees with the number of Lorentz transformations. That is, it is an antisymmetric, $4\times4$ matrix and so has $4\times 3/2=6$ independant entries. So it has $6$ basis elements, $3$ of these correspond to spatial rotations and the other $3$ are our boosts.
\er 

\subsection{Internal Symmetries}

In the above calculations we considered transformations that did something to the spacetime itself, that is the $x^{\mu}$s changed. The question is "is this the only kind of symmetry we can have?" The answer is no, but in order to see it we need to have at least two fields in our Lagrangian. When this is the case we can consider transformations in the plane spanned by these two fields. If our transformations give rise to a symmetry, then we call them \textit{internal symmetries}. Let's look at some examples. 

\bex 
    Consider a theory with 2 real scalar fields $\phi_1$ and $\phi_2$. We can package these together into a two column representation simply as 
    \bse 
        \vec{\phi} := \begin{pmatrix}
            \phi_1 \\
            \phi_2
        \end{pmatrix}.
    \ese 
    We can consider these spanning some 2-dimensional plane as indicated diagrammatically below. Let's then consider the Lagrangian 
    \bse 
        \begin{split}
            \cL & = \frac{1}{2}(\p_{\mu}\vec{\phi})\cdot(\p^{\mu}\phi) - \frac{1}{2}m^2 \vec{\phi}\cdot \vec{\phi} \\
            & = \frac{1}{2}\big[(\p\phi_1)^2 + (\p\phi_2)^2\big] - \frac{1}{2}m^2 \big[ \phi_1^2 + \phi_2^2 \big]
        \end{split}
    \ese 
    where on the first line we have written out the multiple of the derivatives explicitly to make the dot-product clear. This Lagrangian is invariant under rotations in the $\phi_1-\phi_2$ plane, i.e. under 
    \bse 
        \begin{split}
            \phi_1 \longrightarrow \phi_1' & = \phi_1\cos\theta + \phi_2\sin\theta \\
            \phi_2 \longrightarrow \phi_2' & = - \phi_1\sin\theta + \phi_2\cos\theta.
        \end{split}
    \ese 
    This is a continuous transformation and so we can work infinitesimally and find the conserved current. We have 
    \bse 
        \begin{split}
            \phi_1' & \approx \phi_1 + \theta \phi_2 \qquad \implies \qquad \del\phi_1 = \theta\phi_2 \\
            \phi_1' & \approx -\theta\phi_1 + \phi_2 \qquad \implies \qquad \del\phi_2 = -\theta\phi_1.
        \end{split}
    \ese 
    
    \begin{center}
        \btik 
            \draw[thick] (-1,0) -- (3,0);
            \node at (3,-0.3) {\large{$\phi_1$}};
            \draw[thick] (0,-1) -- (0,3);
            \node at (-0.3,3) {\large{$\phi_2$}};
            \draw[thick, dashed, rotate around={20:(0,0)}] (0,0) -- (3,0);
            \node at (3.2,1) {\large{$\phi_1'$}};
            \node at (1.2,0.2) {\large{$\theta$}};
            \draw[thick, dashed, rotate around={20:(0,0)}] (0,0) -- (0,3);
            \node at (-1.3,2.7) {\large{$\phi_2'$}};
            \node at (-0.15,1) {\large{$\theta$}};
        \etik 
    \end{center}
    
    The Lagrangian doesn't change, so our conserved current is simply 
    \bse 
        \begin{split}
            j^{\mu} & = \frac{\p \cL}{\p (\p_{\mu}\phi_1)} \del\phi_1 + \frac{\p \cL}{\p (\p_{\mu}\phi_2)} \del\phi_2 \\
            & = (\p^{\mu}\phi_1)\phi_2 - (\p^{\mu}\phi_2)\phi_1.
        \end{split}
    \ese 
    We can check that this obeys $\p_{\mu}j^{\mu}=0$. First we need the Euler-Lagrange equations for our action. We have 
    \bse 
        0 = \p_{\mu}\bigg(\frac{\p\cL}{\p(\p_{\mu}\phi_1)}\bigg) - \frac{\p \cL}{\p \phi_1} = (\p^2 + m^2)\phi_1
    \ese
    and similarly for $\phi_2$. Now take the derivative of our current:
    \bse
        \begin{split}
            \p_{\mu}j^{\mu} & = (\p^2\phi_1)\phi_2 + \p^{\mu}\phi_1\p_{\mu}\phi_2 - (\p^2\phi_2)\phi_1 - \p^{\mu}\phi_2\p_{\mu}\phi_1 \\
            & = m^2\phi_1\phi_2 - m^2\phi_2\phi_1 \\
            & = 0,
        \end{split}
    \ese 
    where we used the Euler-Lagrange equations to get to the third line, and then used the fact that we have real scalars so $\phi_1\phi_2=\phi_2\phi_1$. 
    
    This symmetry is known as a \textit{global SO(2) symmetry}. The name makes sense: its global (i.e. the whole plane is rotated) and its a rotation in 2-dimensions. It is actually a specific case of the more general global SO($N$) symmetry, which has 
    \bse 
        \frac{N(N-1)}{2}
    \ese 
    conserved currents for $N$ fields. 
\eex 

\bex 
    We can reformulate the previous example by considering the complex scalar field
    \bse 
        \psi = \frac{1}{\sqrt{2}}\big( \phi_1 + i\phi_2\big).
    \ese 
    The above Lagrangian then becomes\footnote{Note that we don't have a factor of $1/2$ in this expression as was the case for the real scalar field. This is just because of the $1/\sqrt{2}$ factor above.} 
    \be 
    \label{eqn:ComplexLagrangianClassical}
        \cL = \p_{\mu}\psi^*\p^{\mu}\psi - m^2 \psi^*\psi.
    \ee 
    We are now rotating in the complex plane so instead of considering a SO(2) rotation we consider the complex U(1) rotation
    \bse 
        \psi \longrightarrow e^{i\a}\psi.
    \ese 
    We already know that this is a symmetry because its exactly the same as the previous example, but what does the current look like in this complexified case? Well simple calculation will give 
    \be 
    \label{eqn:ComplexCurrent}
        j^{\mu} = i\big[ (\p^{\mu}\psi^*)\psi - \psi^*(\p^{\mu}\psi)\big]
    \ee 
\eex 

\bbox 
    Derive the above conserved current. \textit{Hint: Note that $\psi^* \longrightarrow e^{-i\a}\psi^*$ and then work infinitesimally.} 
\ebox 

\br 
    \textcolor{red}{Maybe put a comment here in line with the non-abelian comment made in Tong, page 18.}
\er 

Internal symmetries will prove to be vital for the study of particle physics as a QFT. We will see that the charges arising from such symmetries correspond to things such as electric charge and particle number. 

\section{Hamiltonian Field Theory}

Above we have constructed classical field theory in terms of Lagrangians. It is true that we can extend this Lagrangian approach to QFT by using so-called \textit{path integrals}, and there are advantages to doing that (most notably that the Lorentz invariance is manifest throughout), however in this course we shall use the Hamiltonian approach instead. This uses so-called \textit{canonical quantisation} to promote the fields to operators. 

In particle mechanics, we define the Hamiltonian to be 
\bse 
    H(p,q) = \sum_a p_a\dot{q}_a - L(q,\dot{q}),
\ese
where 
\bse 
    p_a := \frac{\p L}{\p \dot{q}_a}
\ese 
is known as the \textit{conjugate momentum}. The idea is to eliminate $\dot{q}$ wherever we can and replace it with $p$. 

In field theory we do a very similar thing, but now we have to use the Lagrangian density, $\cL$, and we define the \textit{conjugate momentum density} and \textit{Hamiltonian density}
\mybox{
\be 
\label{eqn:ConjugateMomentumDensity}
    \pi(x) := \frac{\p \cL}{\p \dot{\phi}(x)}
\ee 
\be 
\label{eqn:HamiltonianDensity}
    \cH = \pi(x)\dot{\phi}(x) - \cL(x).
\ee 
}
Again we favour replacing any $\dot{\phi}(x)$ dependence with $\pi(x)$ once we know the relation. As with \Cref{eqn:LtocL}, we define the Hamiltonian as the spatial integral over the Hamiltonian density, 
\be 
    H(t) = \int d^3\vec{x} \, \cH(x).
\ee 

Note in this line we appear to have broken Lorentz symmetry as we have picked a preferred time to define our Hamiltonian. However, as long as we are careful not to do anything silly, we should be alright because we started with a Lorentz invariant theory. This is what we mean by the Lorentz symmetry not being manifest: it is not obvious just by looking at the equations that we have Lorentz symmetry, in contrast to the Lagrangian formalism where spacetime indices were always summed over and so it was clear.

\bex 
    As an example let's consider the real scalar field mentioned above with Lagrangian\footnote{It is incredibly common in the field theory world to forget to say `density', as I have done here (and probably above too). The symbols should tell you what we mean: densities are normally given in fancy curly font.}
    \bse 
        \begin{split}
            \cL & = \frac{1}{2}(\p\phi_a)^2 - \frac{1}{2}m^2 \phi_a^2 \\
            & = \frac{1}{2}\dot{\phi}_a^2 - \frac{1}{2}(\nabla\phi_a)^2 - \frac{1}{2}m^2 \phi_a^2.
        \end{split}
    \ese 
    The conjugate momentum is therefore 
    \bse 
        \pi_a(x) = \dot{\phi}_a(x),
    \ese 
    and the Hamiltonian density is 
    \bse
        \begin{split}
            \cH & = \dot{\phi}_a^2 - \cL \\
            & = \frac{1}{2}\dot{\phi}_a^2 + \frac{1}{2}(\nabla\phi_a)^2 + \frac{1}{2}m^2\phi_a^2 \\
            & = \frac{1}{2}\pi_a^2 + \frac{1}{2}(\nabla\phi_a)^2 + \frac{1}{2}m^2\phi_a^2,
        \end{split}
    \ese 
    where in the last line we have done our procedure of swapping out the $\dot{\phi}$s for $\pi$s using the relation above.\footnote{For clarity, you don't just swap $\dot{\phi}\longrightarrow\pi$, but you use the relation for $\pi$ in terms of $\dot{\phi}$ to eliminate the $\dot{\phi}$s. In this particular case, $\pi=\dot{\phi}$, and it does correspond to just swapping them.} The integral over these terms give three contributions to the Hamiltonian, they are 
    \bse 
        H = \int d^3 \vec{x} \, \bigg(\underbrace{\frac{1}{2}\pi_a^2}_{\text{kinetic}} + \underbrace{\frac{1}{2}(\nabla\phi_a)^2}_{\text{shear}} + \underbrace{\frac{1}{2}\phi_a^2}_{\text{mass}}\bigg).
    \ese 
\eex
\chapter{Canonical Quantisation \& Free Klein Gordan Field}

\section{Canonical Quantisation}

Recall the idea behind quantisation in QM is to promote the generalised coordinates in the classical theory and `promote' them to operators acting on the Hilbert space of the quantum theory. This recipe is called \textit{canonical quantisation}. The Poisson bracket\footnote{\textcolor{red}{Note to self: maybe put something about Poisson brackets above when discussing generalised coordinates.}} relation between the generalised coordinates magically transformed into a commutation relation between the operators. That is\footnote{There are factors of $\hbar$s in these equations, but in natural units they go bye-bye.} 
\begin{equation*}
    \begin{split}
        \{q_a,a_b\} = \{p^a,p^b\} = 0 \qquad & \longrightarrow \qquad [\hat{q}_a,\hat{q}_b] = [\hat{p}^a,\hat{p}^b] = 0 \\
        \{q_a,p^b\} = \del^b_a \qquad & \longrightarrow \qquad [\hat{q}_a,\hat{p}^b] = i\del^b_a
    \end{split}
\end{equation*} 
As we are working in Minkowski spacetime, we can lower the indices on the $p$s easily and obtain 
\bse 
    [\hat{q}_a,\hat{q}_b] = [\hat{p}_a,\hat{p}_b] = 0 \qand [\hat{q}_a,\hat{p}_b] = i\del_{ab}.
\ese 

We adopt the same philosophy for fields, and promote them to \textit{operator valued functions}. However we have a small problem: for the particle mechanics case we just had a finite number of generalised coordinates and so it was easy to do, whereas for the fields there's an infinite number, one for each point $\Vec{x}$ in space. We treat this spatial dependence as a label and so need something analogous to the $\del_{ab}$ above. The answer is obviously the usual delta function.  We use the Schr\"{o}dinger picture, where all time dependence appears in the states $\ket{\psi}$ which obey the Schr\"{o}dinger equation
\bse 
    i\frac{d\ket{\psi}}{dt} = H\ket{\psi},
\ese 
and the operators themselves are time independent, and write\footnote{You can work in the Heisenberg picture and define what are known as `equal time' commutation relations.}
\be 
\label{eqn:FieldsCommutation}
    [\phi_a(\Vec{x}),\phi_b(\Vec{y})] = [\pi_a(\Vec{x}),\pi_b(\Vec{y})] = 0, \qand [\phi_a(\Vec{x}), \pi_b(\Vec{y})] = i \del^{(3)}(\Vec{x}-\Vec{y}) \del_{ab}.
\ee 

\br 
\label{rem:ButcheredLorentzInv}
    Note we have really butchered our manifest Lorentz invariance here: we completely separated space and time and made the operators only functions of space! Of course it must be true that we're alright to do this and still get a Lorentz invariant theory, otherwise we wouldn't be doing it in these notes. However this choice of doing things will introduce some factors here and there (e.g. to the measures in integrals) to ensure Lorentz invariance.
\er 

\section{Free Theories}

Recall that the typical goal of QM is to find the spectrum (i.e. eigenvalues) of operators, in particular the Hamiltonian. We now want to do a similar thing for QFTs, however this turns out to be an incredibly hard thing to do, as we now have an uncountably infinite number of degrees of freedom (one for each $\Vec{x}$ value)! The question is "can we somehow get around this?" The answer is "yes, but it restricts the type of theories we consider." What we do is consider theories in which each degree of freedom evolved independently to all others. This is essentially the statement that $\phi(\vec{x})$ and $\phi(\vec{y})$ don't talk to each other unless $\vec{x}=\vec{y}$. For a reasonably self explanatory reason, we call these theories \textit{free theories}.

\br
    Of course free theories are boring from a physical perspective (as nothing interacts so there are essentially no forces!), and we want to study interacting theories. We will return to these later, and study them as perturbations using Feynman diagrams, but first we need to develop the mathematical tools of free theory. So hold tight, more interesting stuff is coming. 
\er 

Ok, so how do we go about quantising the free fields and finding their spectrum? Well the answer is to consider our lovely friend the classical free Klein-Gordan field. Recall that the Euler-Lagrange equations gave us the Klein-Gordan equation, \Cref{eqn:ClassicalKleinGordan}:
\bse 
    (\p^2 + m^2)\phi = 0.
\ese 
As we are working in Minkowski spacetime, we have a global notion of time\footnote{Again this comment is just made because things are different in general spacetimes, see footnote 3 from the first lecture.} and so we can decompose these fields in terms of their Fourier transform
\bse 
    \begin{split}
        \phi(\vec{x},t) & = \int \frac{d^3\vec{p}}{(2\pi)^3} \, \frac{1}{2}\Big[  \widetilde{\phi}(\vec{p},t)e^{i\vec{p}\cdot\vec{x}} + \widetilde{\phi}^*(\vec{p},t)e^{-i\vec{p}\cdot\vec{x}}\Big] \\
        & = \int \frac{d^3\vec{p}}{(2\pi)^3} \, \widetilde{\phi}(\vec{p},t)e^{i\vec{p}\cdot\vec{x}},
    \end{split}
\ese 
where to get to the second line we have identified $\widetilde{\phi}^*(-\vec{p},t) = \widetilde{\phi}(\vec{p},t)$. Using this coordinate system, our Klein-Gordan equation reduces to (dropping the tilde and using the argument to indicate which function it is)
\bse 
    \bigg( \frac{\p^2}{\p t^2} + (\vec{p}+m^2)\bigg)\phi(\vec{p},t) = 0.
\ese

Then the keen-eyed person notes that \textit{for each value} of $\vec{p}$ this corresponds to its own harmonic oscillator with frequency 
\be
\label{eqn:KGHarmonicFrequency}
    \omega_{\vec{p}} = \sqrt{\vec{p}^2+m^2}.
\ee 
To stress the point, we get a harmonic oscillator each \textit{every single} value of $\vec{p}$ independently from any other value. The general field $\phi(\vec{x},t)$ is then simply given by a linear superposition of (an infinite number of) harmonic oscillators. We therefore have, at least classically, achieved the goal above to frame the theory in such a way that each degree of freedom (each $\vec{p}$) evolves independently from the others. Now what we want to do is quantise it. 

\br 
    Note that in order to take the Fourier expansion above we need to assume that the field $\phi(\vec{x},t)$ die off sufficiently quickly as $|\vec{x}|\to \infty$. In order to guarantee this, we use fields from the so-called \textit{Schwartz space}. A precise definition (with a lot more context for why they're useful in QM) can be found in Simon Rea and my \href{https://richie291.wixsite.com/theoreticalphysics/post/dr-frederic-schuller-s-course-of-quantum-theory}{notes on Dr. Schuller's course on quantum theory}, but for here we shall just say they are functions that die off at infinity as do their derivatives. 
\er 

\subsection{A Quantum Field Theorist's Best Friend: The Harmonic Oscillator}

\bnn 
    In this section, and most likely in everything to follow from now on, I am going to drop the hats on quantum operators. I might reinstate them at some points for clarity, but we shall see. 
\enn 

As we have been doing above, let's forget about field theory for a minute and study regular old quantum mechanics. Recall that the quantum harmonic oscillator (QHO) has the Hamiltonian (in 1-dimension)
\be
\label{eqn:QHOHamiltonian}
    H = \frac{1}{2}p^2 + \frac{1}{2}\omega^2q^2,
\ee 
where $p$ and $q$ are the canonical operators obeying 
\be 
\label{eqn:pqcommutator}
    [q,p]=i.
\ee 
We can rewrite this Hamiltonian in a nicer form by introducing the \textit{creation} and \textit{annihilation operators}: $a^{\dagger}$ and $a$, respectively. These are defined such that 
\be 
\label{eqn:qpcreationannihilation}
    q = \frac{1}{\sqrt{2\omega}}\big( a + a^{\dagger} \big), \qand p = -i\sqrt{\frac{\omega}{2}}\big( a - a^{\dagger} \big).
\ee 

\bbox
    Show that \Cref{eqn:pqcommutator} and \Cref{eqn:qpcreationannihilation} give 
    \be 
    \label{eqn:creationqpcreationannihilationcommutator}
        [a,a^{\dagger}] = 1.
    \ee 
    Then using this result show that the Hamiltonian \Cref{eqn:QHOHamiltonian} can be rewritten as 
    \be
    \label{eqn:QHOHamiltonianaadagger}
        H = \omega\bigg(a^{\dagger}a + \frac{1}{2}\bigg).
    \ee 
    Finally show that 
    \be 
    \label{eqn:HaadaggerCommutation}
        [H,a^{\dagger}] = \omega a^{\dagger}, \qand [H,a] = -\omega a.
    \ee 
\ebox

We can now begin to look at the spectrum of the QHO. The \textit{ground state} $\ket{0}$ is defined via 
\be 
\label{eqn:QHOGroundState}
    a\ket{0} = 0
\ee 
and so we see from \Cref{eqn:QHOHamiltonianaadagger} that the ground state has energy 
\bse 
    E_0 = \frac{\omega}{2}.
\ese 
Now using \Cref{eqn:HaadaggerCommutation}, we see that 
\bse 
    \begin{split}
        H(a^{\dagger}\ket{0}) & = a^{\dagger}(H\ket{0}) + \omega (a^{\dagger}\ket{0}) \\
        & = (E_0+\omega)(a^{\dagger}\ket{0}),
    \end{split}
\ese 
and so we see that $a^{\dagger}\ket{0}$ is an eigenvector of $H$ with energy $E=E_0+\omega$. We therefore define \textit{excited states} as\footnote{We're ignoring normalisation here, i.e. $\braket{n}{n}\neq 1$. This is not important for our present discussion.} 
\be 
\label{eqn:QHOExcitedStates}
    \ket{n} := (a^{\dagger})^n\ket{0}.
\ee 
Extending the calculation above, these states have energy 
\be 
\label{eqn:QHOEn}
    E_n = \bigg(n+\frac{1}{2}\bigg)\omega,
\ee 
and so we build a spectrum for our theory as integer steps in $\omega$, as indicated pictorially below. 
\begin{center}
    \btik 
        \draw[thick] (0,0) -- (3,0);
        \node at (-0.5,0) {\large{$\ket{0}$}};
        \node at (3.5,0) {\large{$\frac{1}{2}\omega$}};
        \draw[thick] (0,1) -- (3,1);
        \node at (-0.5,1) {\large{$\ket{1}$}};
        \node at (3.5,1) {\large{$\frac{3}{2}\omega$}};
        \draw[thick] (0,2) -- (3,2);
        \node at (-0.5,2) {\large{$\ket{2}$}};
        \node at (3.5,2) {\large{$\frac{5}{2}\omega$}};
        \node at (1.5,2.75) {\Huge{$\mathbf{\vdots}$}};
    \etik 
\end{center}

As the notation suggests (and as can easily be seen from \Cref{eqn:qpcreationannihilation}) the creation and annihilation operators are Hermitian conjugates to each other. This translates into bra-ket notation in terms of left and right actions: 
\be 
\label{eqn:HermitionLeftRightAction}
    \big(a\ket{\psi}\big)^{\dagger} = \bra{\psi}a^{\dagger}.
\ee 

\subsection{Spectrum Of Quantum Klein-Gordan Field}

Ok so we now know how to get the spectrum of a single QHO. We have also seen that the Klein-Gordan field is essentially an infinite sum (i.e. an integral) of QHOs, so now all we need to do to get the spectrum of the Hamiltonian is integrate over an infinite number of creation and annihilation operators, labelled by $\vec{p}$ --- $a^{\dagger}_{\vec{p}}$ and $a_{\vec{p}}$. To do this we express $\phi$ and $\pi$ as
\be 
\label{eqn:phipicreationannihilation}
    \begin{split}
        \phi(\vec{x}) & = \int \frac{d^3\vec{p}}{(2\pi)^3} \frac{1}{\sqrt{2\omega_{\vec{p}}}} \Big[ a_{\vec{p}} \, e^{i\vec{p}\cdot\vec{x}} + a^{\dagger}_{\vec{p}} \, e^{-i\vec{p}\cdot\vec{x}}\Big] \\
        \pi(\vec{x}) & = \int \frac{d^3\vec{p}}{(2\pi)^3} (-i)\sqrt{\frac{\omega_{\vec{p}}}{2}} \Big[ a_{\vec{p}} \, e^{i\vec{p}\cdot\vec{x}} - a^{\dagger}_{\vec{p}} \, e^{-i\vec{p}\cdot\vec{x}}\Big],
    \end{split}
\ee 
where the form of these expressions makes sense when comparing to \Cref{eqn:qpcreationannihilation}.

\br 
    Note that it is only at this point that we are now considering the \textit{quantum} Klein-Gordan field. This is in contrast why I was careful to emphasise before that we were considering the classical Klein-Gordan field. The equations of motion are still the same, but now they are quantum expressions. 
\er 

This looks nice, but what about the commutators? 
\bcl 
\label{claim:FieldCommutator}
    \Cref{eqn:FieldsCommutation} hold, if and only if, 
    \be 
    \label{eqn:FieldCreationAnnihilationCommutator}
        [a_{\vec{p}} \, , a_{\vec{q}}] = \Big[a^{\dagger}_{\vec{p}} \, , a^{\dagger}_{\vec{q}}\Big] = 0, \qand \Big[a_{\vec{p}} \, , a^{\dagger}_{\vec{q}}\Big] = (2\pi)^3 \del^{(3)}(\vec{p}-\vec{q}).
    \ee 
\ecl 
This claim is not complicated to prove, but a bit of a pain to write out and so I, lovingly, decided to set them as an exercise.

\bbox 
    Prove \Cref{claim:FieldCommutator}. Note it is an if and only if statement so you need to show it both ways, i.e. the $\phi/\pi$ commutators imply the $a/a^{\dagger}$ ones and visa versa. \textit{Hint: If you get very stuck, Prof. Tong sketches one on page 24 of his notes.}
\ebox 

\br
\label{rem:Sandwiching}
    \Cref{eqn:FieldCreationAnnihilationCommutator} seem strange: the left-hand side is a commutator between operators whereas the right-hand side is a delta distribution? The way we wrap our heads around this is to remember that any physical results in QM appear in inner products $\bra{\psi}A\ket{\psi}$, which can be written as integrals and so the delta function makes sense. 
    
    Using \Cref{eqn:HermitionLeftRightAction}, and the fact that our states are orthogonal, we see that it is only the terms that contain $a_{\vec{p}}\, a^{\dagger}_{\vec{p}}$ or $a^{\dagger}_{\vec{p}} \,  a_{\vec{p}}$ that contribute to expectation values. We will see an example of this below when finding the momentum of the first excited state. 
    
    This idea, and similar ones where we get $\C$-numbers on the right-hand side, comes up again and again in QFT. What we have to keep telling ourselves is that "remember we're sandwiching this between states!" and then go from there.
\er 

Ok, so we have already derived that the Hamiltonian for this system is 
\bse 
    H = \frac{1}{2}\int d^3\vec{x} \, \big(\pi^2 + (\nabla\phi)^2 + m^2\phi^2\big).
\ese 
We can substitute \Cref{eqn:phipicreationannihilation} in and obtain 
\be 
\label{eqn:KGHamiltonian}
    \begin{split}
        H & = \int \frac{d^3\vec{p}}{(2\pi)^3} \frac{\omega_{\vec{p}}}{2}\Big[ a_{\vec{p}} \, a^{\dagger}_{\vec{p}} + a^{\dagger}_{\vec{p}} \, a_{\vec{p}}\Big] \\
        & = \int \frac{d^3\vec{p}}{(2\pi)^3} \, \omega_{\vec{p}} \bigg[ a^{\dagger}_{\vec{p}}\, a_{\vec{p}} + \frac{1}{2}(2\pi)^3\del^{(3)}(0)\bigg],
    \end{split}
\ee 
where to get to the last line we have used \Cref{eqn:FieldCreationAnnihilationCommutator}. 

\bbox 
    Another lovely exercise for you: Obtain the first line of \Cref{eqn:KGHamiltonian}. \textit{Hint: Again if you get really stuck, Prof. Tong has done this in his notes, also on page 24.}
\ebox  

\br 
    All jokes aside about the above exercises, it is actually a really beneficial exercise to do and will test your understanding of a reasonable amount of the information leading up to here. So honestly at least give them a go.
\er 

\subsection{A Quantum Field Theorist's Least Favourite Friend: Infinities}

We can extend \Cref{eqn:QHOGroundState} to the field theory case to define the ground state as
\be 
\label{eqn:KGGroundState}
    a_{\vec{p}}\ket{0} = 0 \qquad \forall \vec{p}.
\ee 
We can then use \Cref{eqn:KGHamiltonian} to find the ground state energy: 
\bse 
    H\ket{0} = \bigg[\int d^3\vec{p} \, \frac{1}{2} \omega_{\vec{p}} \del^{(3)}(0)\bigg] \ket{0},
\ese 
which is... infinite?! Uh oh this does not seem good at all, but what did we do wrong? The answer is actually nothing. QFT is filled with infinities and, as Prof. Tong explains "each tells us something important, usually that we're doing something wrong, or asking the wrong question". What we need to do when infinities arise is ask where they come from and what that implies. So let's investigate this one. 

Firstly the bad news... it's actually two infinities! The first one comes from the fact that we're considering the entirety of $\R^3$. So instead let's put our theory inside a box of size $L$. We adopt the usual idea where we prescribe periodic boundary conditions to the sides, and so get the \textit{flat torus}.\footnote{If you're confused why I say flat torus, either google it or think about what happens geometrically when you make these periodic boundary conditions.} Finally we just let $L\to\infty$ to get our result back. We can then use 
\bse 
    (2\pi)^3 \del^{(3)}(0) = \lim_{L\to\infty} \int_{-L/2}^{L/2} d^3\vec{x} \, e^{i\vec{x}\cdot\vec{p}} \bigg|_{\vec{p}=0} = \lim_{L\to\infty} \int_{-L/2}^{L/2}d^3\vec{x} = V,
\ese 
where $V$ is the volume of our theory. Ok so the $(2\pi)^3\del(0)$ infinity is just because we're considering the total energy $E_0$ instead of the \textit{energy density}
\bse 
    \varepsilon_0 := \frac{E_0}{V} = \int \frac{d^3\vec{p}}{(2\pi)^3} \frac{1}{2}\omega_{\vec{p}}.
\ese 
This is still infinite though! Why? well because we take the positive root in \Cref{eqn:KGHarmonicFrequency} and so we're summing an infinite number of positive numbers! This seems like a more tricky beast to tame, however then a light bulb goes off in our heads and we realise that this is just the ground state energy, and so, from the extension of \Cref{eqn:QHOEn}, it will appear in \textit{all} the energies. Then remembering that all we can measure physically is energy \textit{differences} we realise that we are free to just `take this away' from every energy in our system (as the difference wont be effected), so this is what we shall do. 

You might not be very comfortable with `subtracting infinity', and if that is the case, allow me to provide a calming remark: we are mere mortals pretending to be Gods. In less poetic (and perhaps offensive) words: we have assumed that the theory we have written down is valid up to arbitrarily large momentum. This corresponds to arbitrarily large energies, or arbitrarily low length scales. Nature is highly unlikely to agree with us and so plays the trump card of "you need to cut-off your theory at high momenta!" If we do this, the ground state energy density will become finite and then we can comfortably take it away from all energies and proceed as if nothing ever happened. Infinities of this kind are called \textit{ultra-violet divergences}. 

\subsection{Finally, The Spectrum}

Now that we have tamed our infinities, we can proceed to finding the spectrum of the Klein-Gordan field. Continuing with the extension of the QHO, we have 
\be 
\label{eqn:KGHamiltonianAAdaggerCommutators}
    [H,a_{\vec{p}}\,]\ket{0} = -\omega_{\vec{p}}\, a_{\vec{p}}\ket{0}, \qand [H,a^{\dagger}_{\vec{p}}\,] = \omega_{\vec{p}}\, a^{\dagger}_{\vec{p}}\ket{0},
\ee 
which encourages us to define 
\bse 
    \ket{\vec{p}} := a^{\dagger}_{\vec{p}}\, \ket{0}.
\ese 
But now we have different creation/annihilation operators and so can excite the ground state in different ways. We therefore use the notation 
\bse 
    \ket{\vec{p}, \vec{q}, ...} := \Big(a^{\dagger}_{\vec{p}} \, a^{\dagger}_{\vec{q}} \,  ... \Big)\ket{0}.
\ese 
Then finally using \Cref{eqn:KGHamiltonianAAdaggerCommutators} gives us 
\bse 
    H\ket{\vec{p}, \vec{q},...} = (\omega_{\vec{p}} + \omega_{\vec{q}} + ... )\ket{0}.
\ese 
This exhausts our spectrum.\footnote{And perhaps it has exhausted you getting to this point.}

\br 
    Just as with the states $\ket{n}$ defined in \Cref{eqn:QHOExcitedStates}, our states $\ket{\vec{p}}$ are not normalised. They are orthogonal though. We shall return to the normalisation later.
\er 

\subsection{Interpreting The Eigenstates: Particles}

So we have our spectrum, now we want to interpret what they mean physically. Well let's focus on the first excited state
\bse 
    \ket{\vec{p}} = a^{\dagger}_{\vec{p}}\ket{0}.
\ese 
We have already seen that this has energy 
\bse 
    E_{\vec{p}} = \omega_{\vec{p}} = \sqrt{\vec{p}^2 +m^2},
\ese 
and have already remarked that this is the energy of a relativistic particle. But could this just be some coincidence? Well let's look at the \textit{physical} momentum given by \Cref{eqn:EnergyAndMomentumCharges}. Straight forward calculation using the Lagrangian for our system gives 
\bse 
    T^{\mu\nu} = \p^{\mu}\phi\p^{\nu}\phi - \eta^{\mu\nu}\cL,
\ese
and so the momentum is 
\bse 
    P^i = \int d^3\vec{x} \, \dot{\phi}(x) \p^i \phi(x),
\ese 
which as a operator we can write\footnote{Recall that we're using the field theorist's convention of $(+,-,-,-)$.}
\bse 
    \vec{P} = - \int d^3\vec{x} \, \pi(\vec{x})\nabla\phi(\vec{x})
\ese
If we then use \Cref{eqn:phipicreationannihilation}, we get 
\bse 
    \begin{split}
        \vec{P} & = \int \frac{d^3\vec{p}}{(2\pi)^3} \frac{\vec{p}}{2}\Big[ a^{\dagger}_{\vec{p}} \, a_{\vec{p}} +  a_{\vec{p}} \, a^{\dagger}_{\vec{p}}  + a_{\vec{p}}\, a_{\vec{-p}} + a^{\dagger}_{\vec{p}}\, a^{\dagger}_{\vec{-p}} \Big] \\
        & = \int \frac{d^3\vec{p}}{(2\pi)^3}\vec{p} a^{\dagger}_{\vec{p}} \, a_{\vec{p}},
    \end{split}
\ese 
where to get to the last line we have used the commutation relation between $a/a^{\dagger}$ and then dropped the $\del(0)$ term as we did for the energy, and then `dropped' the $aa$ and $a^{\dagger}a^{\dagger}$ terms using the argument of \Cref{rem:Sandwiching}. So we see that the first excited state has momentum 
\bse 
    \vec{P}\ket{\vec{p}} = \vec{p}\ket{\vec{p}},
\ese 
which is exactly what we expect for a particle. We therefore \textit{interpret} the state $\ket{\vec{p}}$ to be a particle with 3-momentum $\vec{p}$. It is important to note that this is simply an interpretation, what we're really dealing with is an excitation of a field. 

\bbox 
    Using the definition 
    \bse 
        J^i = \epsilon^{ijk} \int d^3\vec{x} \, (\cJ^0)^{jk},
    \ese 
    where $\epsilon^{ijk}$ is the Levi-Civita tensor density and $\cJ$ is the classical angular momentum defined in \Cref{eqn:LorentzCurrents}, show that the state $\ket{\vec{p}=0}$ has 
    \bse 
        J^i\ket{\vec{p}=0} = 0.
    \ese 
\ebox 

The result of the above exercise tells us that the our interpreted particle has no spin (i.e. no internal angular momentum). That is, our quantisation of the Klein-Gordan field gives rise to a spin-0 particle. 

\subsection{Multiparticle States}

Above we just considered the first excited state, but we have already seen that we can apply the creation operators multiple times, i.e. we have 
\bse 
    \ket{\vec{p}_1,...,\vec{p}_n} = a^{\dagger}_{\vec{p}_1}...a^{\dagger}_{\vec{p}_n}\ket{0}.
\ese 
As we have used $n$ creation operators, we refer to these states are $n$-particle states and interpret it as $n$ particles with 3-momenta $(\vec{p}_1,...,\vec{p}_n)$. As the creation operators commute with each other we have 
\bse 
    \ket{\vec{p},\vec{q}} = \ket{\vec{q},\vec{p}},
\ese 
and so the particles are symmetric under interchange. We therefore see that these particles are \textit{bosons}. This also agrees with the fact that they are spin-0, as per the above exercise. 

Now recall that in the first lecture we showed when we unite special relativity with QM we have to account for particle number changing. We said that we do this by changing our Hilbert space to be a Fock space, \Cref{eqn:FockSpace}. We can clarify a bit what we meant there now. If $\cH$ is the Hilbert space of our 1-particle states (i.e. $\ket{\vec{p}}\in\cH$), then we have
\bse 
    \ket{\vec{p}_1,...,\vec{p}_n} \in \cH^{\otimes n}.
\ese 
We then take the direct sum over all these different states so that our total Hilbert space accounts for the particle creation process. That is, let's say we start off with a state $\ket{\psi}\in\cH^{\otimes n}$ and then something happens to it and causes the production of a new particle. The new state would then be an element of $\cH^{\otimes (n+1)}$. If we don't take the Fock space construction as above, this process would mean that we leave our Hilbert space, which is a big no-no.

\subsection{Interpretation Of $\phi(\vec{x})\ket{0}$}

The next thing we can ask is how does the action of $\phi(\vec{x})$ fit into our particle interpretation? Well let's calculated it and see. We have 
\bse 
    \begin{split}
        \phi(\vec{x})\ket{0} & = \int \frac{d^3\vec{p}}{(2\pi)^3} \frac{1}{\sqrt{2\omega_{\vec{p}}}} \Big[ a_{\vec{p}} \, e^{i\vec{p}\cdot\vec{x}} + a^{\dagger}_{\vec{p}} \, e^{-i\vec{p}\cdot\vec{x}} \Big]\ket{0} \\
        & = \int \frac{d^3\vec{p}}{(2\pi)^3} \frac{1}{\sqrt{2\omega_{\vec{p}}}} e^{-i\vec{p}\cdot\vec{x}} \ket{\vec{p}}.
    \end{split}
\ese 
Now compare this to the result from QM\footnote{This result comes from inserting a complete set of states, $\b1 = \int dp \ket{p}\bra{p}$. See any decent QM book for details.}
\bse 
    \ket{x} = \int dp e^{-ipx} \ket{p}.
\ese 
We can therefore interpret $\phi(\vec{x})\ket{0}$ as a particle at position $\vec{x}$. 

\subsection{Number Operator}

We have just talked about constructing a whole Fock space to take into account the fact that particle number can change. Well in our free theory this construction is actually not important because nothing is interacting and so there's no chance for more particles to be produced. We can show this explicitly by defining the \textit{number operator}
\be 
\label{eqn:NumberOperator}
    N := \int \frac{d^3p}{(2\pi)^3} a^{\dagger}_{\vec{p}} \, a_{\vec{p}},
\ee 
which satisfies 
\bse 
    N\ket{\vec{p}_1,...,\vec{p}_n} = n\ket{\vec{p}_1,...,\vec{p}_n}.
\ese 
The proof that particle number is conserved in a free theory is the content of the next exercise. 

\bbox 
    Using the above definition and \Cref{eqn:KGHamiltonian} (with the delta term dropped) show that 
    \be
    \label{eqn:NumberOperatorConserved}
        [H,N] = 0.
    \ee 
\ebox  

This tells us that once we are in a particular $\cH^{\otimes n}$ in the Fock space, we will never leave it in a free theory. This makes them a bit boring to study. Fortunately later we will allow for interactions in our theory and see that the above result no longer holds in general. The sad news is free theories are the only ones we can solve exactly, and so in order to study interacting theories we are going to have to develop the mathematics of Feynman diagrams. These are essentially life saving tools to calculate things in perturbation expansions of interacting QFTs. 
\chapter{Restoring Lorentz Invariance \& Causality}

\section{Operator Valued Distributions}

Recall that last lecture we gave an interpretation of our results as particles. There is a problem with this interpretation: they are momentum eigenstates and so, by Heisenberg uncertainty principle, are not localised in position space at all. This problem essentially stems from the fact that in QM the position and momentum eigenstates are not good elements of the Hilbert space as the normalise to delta-functions. The QFT equivalent is that the operators $\phi(\vec{x})$ and $a_{\vec{p}}$ are not good operators on the Fock space as the resulting states are not normalisable. Explicitly we have 
\bse 
    \bra{0}a_{\vec{p}}a^{\dagger}_{\vec{p}}\ket{0} = \braket{\vec{p}}{\vec{p}} = (2\pi)^3 \del^{(3)}(0), \qand \bra{0}\phi(\vec{x})\phi(\vec{x})\ket{0} = \braket{\vec{x}}{\vec{x}} = \del(0).
\ese 

These are known as \textit{operator valued distributions}, and the way we fix this problem is by smearing them over space. For example the wavepacket 
\bse 
    \ket{\varphi} = \int \frac{d^3\vec{p}}{(2\pi)^3} e^{-i\vec{p}\cdot\vec{x}} \varphi(\vec{p})\ket{\vec{p}},
\ese 
is partially localised in both momentum and position space. A typical state is an element of the Schwarz space and is of the form 
\bse 
    \varphi(\vec{p}) = e^{-\vec{p}^2/2m^2}.
\ese

\section{Relativistic Normalisation}

This problem has come from the fact that we haven't been careful to maintain Lorentz invariance. This is what we were warning about in \Cref{rem:ButcheredLorentzInv}. So what's broken our Lorentz invariance? The answer is that 
\bse 
    \braket{\vec{p}}{\vec{p}} = (2\pi)^3 \del^{(3)}(0),
\ese 
which is a $3$-delta function, i.e. we have lost a coordinate! Now it's possible that we're lucky and this object is still Lorentz invariant, however this is not the case. We see this easily by taking a Lorentz transformation
\bse 
    p^{\mu} \longrightarrow {\Lambda^{\mu}}_{\nu} p^{\nu}.
\ese 
If our expectation value is going to be unchanged, we would need the transformation on the states to be unitary, i.e. 
\bse 
    \ket{\vec{p}} \longrightarrow U(\Lambda)\ket{\vec{p}},
\ese 
but there is absolutely no reason why this would be the case. 

So what do we do? We need to normalise our one-particle states so that we somehow get a Lorentz invariant inner product. We do this using two tricks. 

\subsection{Trick 1}

The first thing we note is that the integral $\int d^4p$ is manifestly Lorentz invariant, provided the integrand is also Lorentz invariant. Then we note that the relativistic dispersion relation
\bse 
    p_{\mu}p^{\mu} = m^2 
\ese 
is also manifestly Lorentz invariant. This can be written as 
\bse 
    p_0^2 = \vec{p}^2 + m^2 = \omega_{\vec{p}}^2,
\ese 
which has two branches, $p_0 = \pm\omega_{\vec{p}}^2$, as depicted below.

\begin{center}
    \btik 
        \draw[thick] (-2,-2) -- (2,2);
        \draw[thick] (2,-2) -- (-2,2);
        %
        \draw[thick, ->] (0,0) -- (2.5,0);
        \node at (2.5,-0.5) {\large{$\vec{p}$}};
        \draw[thick] (0,0) -- (0,1.8);
        \node at (-0.5,3) {\large{$p_0$}};
        %
        \draw[thick, fill = gray!40, opacity = 0.8] (1.5,2) .. controls (0,0.5) .. (-1.5,2) arc (-180:360: 1.5 and 0.2);
        \draw[thick] (1.5,2) .. controls (0,0.5) .. (-1.5,2) arc (-180:360: 1.5 and 0.2);
        \draw[thick] (-1.5,2) arc (180:360: 1.5 and 0.2);
        %
        \draw[thick, dashed] (-1.5,-2) arc (-180:360: 1.5 and 0.2);
        \draw[thick, fill = gray!40, opacity = 0.8] (1.5,-2) .. controls (0,-0.5) .. (-1.5,-2) arc (180:360: 1.5 and 0.2);
        \draw[thick] (1.5,-2) .. controls (0,-0.5) .. (-1.5,-2) arc (180:360: 1.5 and 0.2);
        % 
        \draw[thick, ->] (0,1.8) -- (0,3);
    \etik 
\end{center}

The choice of branch is Lorentz invariant, so w.l.o.g. we take $p_0>0$. We we know the integral 
\bse 
    \int \frac{d^4p}{(2\pi)^4} 2\pi \del(p^2-m^2) \Theta(p_0)
\ese 
where $\Theta$ is the Heaviside function, is Lorentz invariant. We can rewrite this in polar coorindates as
\bse 
    \int \frac{d^3\vec{p}}{(2\pi)^3} \int_0^{\infty} dp_0 \, \del(p_0^2 - \vec{p}^2 -m^2) = \int \frac{d^3\vec{p}}{(2\pi)^3} \int_0^{\infty} \frac{dp_0}{2p_0} \del(p^0 - \omega_{\vec{p}}) = \int \frac{d^3\vec{p}}{(2\pi)^3} \frac{1}{2\omega_{\vec{p}}}
\ese 
where the middle term comes from
\bse 
    \del\big(f(x) - f(x_0)\big) = \frac{1}{|f'(x_0)|} \del(x-x_0).
\ese 
The thing we started with was Lorentz invariant and so we know that that the measure 
\be 
\label{eqn:LorentzInvariantMeasure}
    \int \frac{d^3\vec{p}}{(2\pi)^3}\frac{1}{2\omega_{\vec{p}}}
\ee 
is a Lorentz invariant measure.

\subsection{Second Trick}

Now we use 
\bse 
    \int \frac{d^3\vec{p}}{(2\pi)^3} \bra{0}a_{\vec{q}}\, a_{\vec{p}}^{\dagger} \ket{0} = \int \frac{d^3\vec{p}}{(2\pi)^3} \, (2\pi)^3 \del^{(3)}(\vec{p}-\vec{q}) = 1,
\ese 
which is clearly Lorentz invariant (it's just a number). We want to write this in terms of our Lorentz invariant measure, \Cref{eqn:LorentzInvariantMeasure}. That is:
\bse 
    \int \frac{d^3\vec{p}}{(2\pi)^3} \frac{1}{2\omega_{\vec{p}}} \bra{0} \big(\sqrt{2\omega_{\vec{q}}} \, a_{\vec{q}}\big) \big(\sqrt{2\omega_{\vec{p}}} \, a_{\vec{p}}^{\dagger}\big),
\ese 
and conclude that if we normalise our one-particle states as 
\be 
\label{eqn:RelativisticNormalisation}
    \ket{\vec{p}} = \sqrt{2\omega_{\vec{p}}} \, a_{\vec{p}}^{\dagger} \ket{0} = \sqrt{2E_{\vec{p}}} \, a_{\vec{p}}^{\dagger} \ket{0} 
\ee 
where on the second equality we have just used $E$, for energy, as this is standard notation, then our inner products become Lorentz invariant.

\section{Heisenberg Picture}

As we remarked above, we slaughtered our Lorentz invariance through our quantisation procedure. Even besides the above mix of the measure, we could have easily seen (or at least been cautious) by the fact that we are working in the Schr\"{o}dinger picture, so our operators $\phi(\vec{x})$ only depend on space, and not time. Equally our one states evolve in time, with the evolution given by the Schr\"{o}dinger equation, which implies
\bse 
    \ket{\psi(t)}_S = e^{-iH(t-t_0)} \ket{\psi(t_0)}_S.
\ese 

So we would be closer to having manifest Lorentz invariance if we stripped the time dependence from the states and put it in the operators. This is exactly what the Heisenberg picture does. The relation between states is 
\be 
\label{eqn:HeisenbergStates}
    \ket{\psi}_H = e^{iH(t-t_0)} \ket{\psi(t)}_S \qquad \implies \qquad \frac{d}{dt}\ket{\psi}_H = 0.
\ee  
The operators are related by 
\be 
\label{eqn:HeisenbergOperators}
    \cO_H = e^{iHt}\cO_S e^{-iHt} 
\ee 
which tells us 
\be 
\label{eqn:HeisenbergOperatorsEvolution}
    \frac{d}{dt}\cO_H = i [H,\cO_H].
\ee 

So in our field theory our Schr\"{o}dinfer picture operators $\phi(\vec{x})/\pi(\vec{x})$ become the Heisenberg operators 
\be 
    \phi(x) = e^{iHt}\phi(\vec{x}) e^{-iHt}, \qand \pi(x) = e^{iHt}\pi(\vec{x}) e^{-iHt},
\ee 
where, as is standard in field theory, we have dropped the $H/S$ subscripts and used the arguments to tell us which picture we're in (i.e. we've used $\phi(x)=\phi(\vec{x},t)$). We adapt the commutation relations \Cref{eqn:FieldsCommutation} to be \textit{equal time} commutations:
\be 
\label{eqn:EqualTimeCommutators}
    [\phi_a(\vec{x},t),\phi_b(\vec{y},t)] = [\pi_a(\vec{x},t),\pi_b(\vec{y},t)] = 0, \qand [\phi_a(\vec{x},t),\pi_b(\vec{y},t)] = i \del^{(3)}(\vec{x}-\vec{y})\del_{ab}.
\ee 

We can now study to see if the Heisenberg equations, \Cref{eqn:HeisenbergOperatorsEvolution}, for $\phi(x)$ and $\pi(x)$ are equivalent to the Klein-Gordan equation. Using 
\bse 
    H = \int d^3\vec{x} \, \frac{1}{2}\big(\pi^2 + (\nabla\phi)^2 + m^2\phi^2\big),
\ese 
and the commutations above
\bse 
    \begin{split}
        \frac{d}{dt}\phi & = \frac{i}{2} \Big[ \int d^3\vec{y}\,  \pi^2(y) , \phi(x) \Big] \\
        & = \frac{i}{2} \int d^3y \Big( \pi(y)[\pi(y), \phi(x)] + [\pi(y),\phi(x)]\pi(y)\Big) \\
        & = i \int d^3 \vec{y} \, \pi(y) (-i)\del^{(3)}(\vec{y}-\vec{x}) \\
        & = \pi(x).
    \end{split}
\ese 
That is, 
\be 
\label{eqn:dotphiHeisenberg}
    \dot{\phi}(x) = \pi(x).
\ee 

\bbox 
    Using \Cref{eqn:HeisenbergOperatorsEvolution} show that 
    \be
    \label{eqn:dotpiHeisenberg}
        \dot{\pi}(x) = (\nabla^2-m^2)\phi(x).
    \ee
\ebox  

Combining \Cref{eqn:dotphiHeisenberg,eqn:dotpiHeisenberg}, we have 
\bse 
    (\nabla^2-m^2)\phi(x) = \dot{\pi} = \Ddot{\phi},
\ese 
or more familiarly, 
\bse 
    (\p^2 + m^2)\phi = 0,
\ese 
the Klein-Gordan equation! We refer to this as the \textit{correspondence principle}.

\bbox 
    Show that 
    \bse  
        e^{iHt}a_{\vec{p}} \, e^{-iHt} = e^{-iE_{\vec{p}}} a_{\vec{p}}, \qand e^{iHt}a_{\vec{p}}^{\dagger} \, e^{-iHt} = e^{iE_{\vec{p}}} a_{\vec{p}}^{\dagger}.
    \ese 
    \textit{Hint: Use} 
     \bse  
        e^{iHt}a_{\vec{p}}^{\dagger} \, e^{-iHt} = \sum_{n=0}^{\infty} \frac{(iHt)^n}{n!} a_{\vec{p}}^{\dagger} \, e^{-iHt}, \qand [H, a_{\vec{p}}^{\dagger}\, ] = E_{\vec{p}}\, a_{\vec{p}}^{\dagger}.
    \ese
\ebox  

We can find the mode expansion of $\phi(x)$ in the Heisenberg picture by starting with the Schr\"{o}dinger expression
\bse 
    \phi(\vec{x}) =  \int \frac{d^3 \vec{p}}{(2\pi)^3} \frac{1}{\sqrt{2E_{\vec{p}}}} \Big( e^{i\vec{p}\cdot\vec{x}}a_{\vec{p}} \, + e^{-i\vec{p}\cdot\vec{x}}a_{\vec{p}}^{\dagger}\Big)
\ese 
and then use \Cref{eqn:HeisenbergOperators} and the result of the previous exercise to obtain
\be 
\label{eqn:HeisenbergPhi}
    \begin{split}
        \phi(x) & = \int \frac{d^3 \vec{p}}{(2\pi)^3} \frac{1}{\sqrt{2E_{\vec{p}}}} \Big( e^{-iE_{\vec{p}}t + i\vec{p}\cdot\vec{x}} a_{\vec{p}} \, + e^{iE_{\vec{p}}t - i\vec{p}\cdot\vec{x}} a_{\vec{p}}^{\dagger}\Big) \\
        & = \int \frac{d^3 \vec{p}}{(2\pi)^3} \frac{1}{\sqrt{2E_{\vec{p}}}} \Big( a_{\vec{p}} \, e^{-ipx} + a_{\vec{p}}^{\dagger}\, e^{ipx}\Big)\Big|_{p_0=E_{\vec{p}}}.
    \end{split}
\ee 

You can then quickly arrive at the expression for $\pi(x)$ by using \Cref{eqn:dotphiHeisenberg}, giving 
\be 
\label{eqn:HeisenbergPi}
    \pi(x) = \int \frac{d^3\vec{p}}{(2\pi)^3} i\sqrt{\frac{E_{\vec{p}}}{2}}\Big( a_{\vec{p}}^{\dagger} \, e^{ipx} - a_{\vec{p}} \, e^{-ipx}\Big)\Big|_{p_0=E_{\vec{p}}}.
\ee 
These look a lot like \Cref{eqn:phipicreationannihilation}, but now in terms of $4$-vectors. This seems more manifestly Lorentz invariant!

So we see that $\phi(x)$ is a superposition of plane waves $e^{\pm ip_0t}$, but the coefficients create and destroy particles. We can therefore think of the field as a `hammer' which bashes the vacuum and shakes quanta out of it. That is the operator acting on the vacuum creates a particle:
\bse 
    \phi(\vec{x},0) \ket{0} = \int \frac{d^3\vec{p}}{(2\pi)^3} \frac{1}{2E_{\vec{p}}} e^{-i\vec{p}\cdot\vec{x}} \ket{\vec{p}}, \qquad \implies \qquad  \bra{\vec{p}}\phi(\vec{x},0)\ket{0} = e^{-i\vec{p}\cdot\vec{x}}.
\ese 
If we act on an $n$ particle state instead, because $\phi(x)$ contains both creation and annihilation operators, it has a amplitude to create both $(n+1)$ and $(n-1)$ particle states. 

\section{Causality}

As we have mentioned above, in order for our theory to be causal, we want spacelike separated operators to commute, that is 
\bse 
    [\cO_1(x),\cO_2(y)] = 0 \qquad \forall (x-y)^2 <0.
\ese 
This statement is obviously made in the Heisenberg picture, as the operators are functions of spacetime, not just space. So we need to study the commutator between two fields $[\phi(x),\phi(y)]$, but where now we don't impose equal time, i.e. $x^0 \neq y^0$ in general. We define 
\be 
\label{eqn:Delta(x-y)}
    \Delta(x-y) := [\phi(x),\phi(y)]
\ee
Plugging in the definition \Cref{eqn:HeisenbergPhi}, we have 
\bse 
    \begin{split}
        \Delta(x-y) & = \int \frac{d^3\vec{p}}{(2\pi)^3} \frac{d^3\vec{q}}{(2\pi)^3} \frac{1}{2\sqrt{E_{\vec{p}}E_{\vec{q}}}} \Big[ a_{\vec{p}}\, e^{ipx} + a_{\vec{p}}^{\dagger} \, e^{-ipx} \, ,  a_{\vec{q}} \, e^{iqy} + a_{\vec{q}}^{\dagger}\, e^{-iqy} \big] \\
        & = \int \frac{d^3\vec{p}}{(2\pi)^3} \frac{d^3\vec{q}}{(2\pi)^3} \frac{1}{2\sqrt{E_{\vec{p}}E_{\vec{q}}}} \Big( e^{-ipx+ipy}  + e^{ipx-ipy}\Big) (2\pi)^3 \del^{(3)}(\vec{p}-\vec{q}) \\
        & = \int \frac{d^3\vec{p}}{(2\pi)^3} \frac{1}{2E_{\vec{p}}} \Big( e^{-ip(x-y)}  - e^{ip(x+y)}\Big)
    \end{split}
\ese 

The first thing we note is that this expression is Lorentz, as we have the correct measure, \Cref{eqn:LorentzInvariantMeasure}, and the exponential contain $4$-vectors. We can therefore choose a specific reference frame in order to evaluate its timelike/spacelike properties. 

\ben[label=(\roman*)]
    \item First look at the spacelike separations, i.e. $(x-y)^2>0$. Let's work in a reference frame where $(x-y) = (0,\vec{a})$, then we have 
    \bse 
        \Delta(0,\vec{a}) = \int \frac{d^3\vec{p}}{(2\pi)^3} \frac{1}{E_{\vec{p}}} \Big( e^{i\vec{p}\cdot \vec{a}} - e^{-i\vec{p}\cdot \vec{a}}\Big).
    \ese 
    Now note that $E_{\vec{p}} = \sqrt{\vec{p}^2+m^2}$ and so $E_{-\vec{p}} = E_{\vec{p}}$, and also that the measure is invariant under $\vec{p}\longrightarrow -\vec{p}$, as you get a minus sign but the integration limits flip. So we can freely change the sign in the second exponential, and then the two terms cancel, i.e. 
    \bse 
        \Delta(0,\vec{a}) = 0.
    \ese 
    So our theory is causal! 
    \item Let's also check that it doesn't vanish for timelike separations (which would be an obvious problem). Again can pick a frame, so simply choose $(x-y) = (t,0)$, then we have 
    \bse 
        \Delta(t,0) = \int \frac{d^3\vec{p}}{(2\pi)^3} \Big( e^{-iE_{\vec{p}}\,t} - e^{iE_{\vec{p}}\,t} \Big) \sim \int d^3\vec{p}\, \sinh\big(\sqrt{\vec{p}^2+m^2} \, t\big),
    \ese 
    which does not vanish (its a Bessel function).
\een

\br 
    Note that the objects on the right-hand side of \Cref{eqn:Delta(x-y)} are operators. However we've just shown that it comes out to be a $\C$-number. Essentially what we have done is sandwich it in between two vacuum states, as per \Cref{rem:Sandwiching}. In other words we really should have written 
    \bse 
        \bra{0}[\phi(x),\phi(y)]\ket{0},
    \ese 
    but its common field theory notation to drop the bra and ket. We will continue to do this next lecture, but this is just another remark to go with \Cref{rem:Sandwiching} to remind us that everything needs to be sandwiched. 
\er 
\chapter{Complex Scalar Field \& Propagators}

\section{Complex Scalar Field}

Recall the \textit{classical} Lagrangian for the complex scalar field, \Cref{eqn:ComplexLagrangianClassical}, 
\bse 
    \cL = \p_{\mu}\psi^*\p^{\mu}\psi - m^2 \psi^*\psi.
\ese 
We treat $\psi$ and $\psi^*$ as independent fields and get the equations of motion 
\bse 
    \big(\p^2 + m^2 \big)\psi = 0, \qand \big(\p^2 + m^2\big) \psi^* = 0.
\ese 
When we quantise this theory we again treat each field as an independent degree of freedom. In quantising we replace the complex conjugate with the Hermitian conjugate. We want to obtain an expansion equivalent to that of \Cref{eqn:phipicreationannihilation}, but we have to note something: the field $\psi$ is \textit{not} real and so the corresponding quantum operator will \textit{not} be Hermitian. We therefore need to have different labels for our creation and annihilation operators in the expansion for $\psi$. We then get the expansion for $\psi^{\dagger}$ by taking the Hermitian conjugate. The result is 
\be 
\label{eqn:ComplexFieldExpansion}
    \begin{split}
        \psi & = \int \frac{d^3\Vec{p}}{(2\pi)^3} \frac{1}{\sqrt{2E_{\vec{p}}}} \Big( b_{\vec{p}} \,  e^{i\vec{p}\cdot\vec{x}} + c_{\vec{p}}^{\dagger} \, e^{-i\vec{p}\cdot\vec{x}}\Big) \\
        \psi^{\dagger} & = \int \frac{d^3\Vec{p}}{(2\pi)^3} \frac{1}{\sqrt{2E_{\vec{p}}}} \Big( c_{\vec{p}} \,  e^{i\vec{p}\cdot\vec{x}} + b_{\vec{p}}^{\dagger} \, e^{-i\vec{p}\cdot\vec{x}}\Big) 
    \end{split}
\ee 

\br 
    Note that the dagger comes with the negative exponential. This is because we want the particles we create to have positive frequency, which corresponds to $e^{-i\vec{p}\cdot\vec{x}}$. 
\er 

Classically we had the result $\pi = \p\cL/\p\dot{\psi}$, \Cref{eqn:ConjugateMomentumDensity}. We turn this into a quantum relation, and using $\p\cL/\p\dot{\psi} = \dot{\psi}^*$, we obtain the operator expressions
\be 
\label{eqn:ComplexPiExpansion}
    \begin{split}
        \pi & = \int \frac{d^3\vec{p}}{(2\pi)^3} i \sqrt{\frac{E_{\vec{p}}}{2}} \Big( - c_{\vec{p}}e^{i\vec{p}\cdot\vec{x}} + b_{\vec{p}}^{\dagger} e^{-i\vec{p}\cdot\vec{x}}\Big) \\
        \pi^{\dagger} & = \int \frac{d^3\vec{p}}{(2\pi)^3} (-i) \sqrt{\frac{E_{\vec{p}}}{2}} \Big(  b_{\vec{p}}e^{i\vec{p}\cdot\vec{x}} - c_{\vec{p}}^{\dagger} e^{-i\vec{p}\cdot\vec{x}}\Big)
    \end{split}
\ee 

\br 
    I have chosen to write the above equations so that the creation operator always appears to the far most right. This means that the $b$s and $c$s move around, obviously. There are two reasons I am writing them like this: it's how Prof. Spannowsky writes them, and it then they take the same form as \Cref{eqn:phipicreationannihilation}. Obviously it doesn't matter which order I write them in, I just make this remark because Prof. Tong writes them so that the $b$ term is always to the left (see pages 33/34 of his notes).
\er 

We want to keep the same form for the \textit{equal time} commutation relations as for the real scalar field, i.e.
\be 
    [\psi(\vec{x}), \pi(\vec{y})] = i \del^{(3)}(\vec{x}-\vec{y}), \qquad [\psi(\vec{x}),\pi^{\dagger}(\vec{y})] = 0,
\ee 
and similarly with Hermitian conjugates everywhere,\footnote{Note you can simply use $[A,B]^{\dagger}=-[A^{\dagger},B^{\dagger}]$, along with the $i$s on the right-hand, to get these results.} and the vanishing ones 
\bse 
    [\psi(\vec{x}),\psi(\vec{y})] = 0 = [\pi(\vec{x}),\pi(\vec{y})].
\ese 

\bbox 
    Show that the above commutation relations hold if, and only if, we also have 
    \be 
        [b_{\vec{p}} \, b_{\vec{p}}^{\dagger}\,] = (2\pi)^3 \del^{(3)}(\vec{p}-\vec{q}), \qand [c_{\vec{p}} \, c_{\vec{p}}^{\dagger}\,] = (2\pi)^3 \del^{(3)}(\vec{p}-\vec{q}),
    \ee 
    and all others vanishing, i.e. 
    \bse 
        [b_{\vec{p}} \, b_{\vec{p}}\,] = 0 = [b_{\vec{p}} \, c_{\vec{p}}\,] 
    \ese 
    etc.
\ebox 

\bbox 
    Derive the following expression for the Hamiltonian, 
    \be 
    \label{eqn:ComplexFieldHamiltonian}
        H = \frac{1}{2} \int \frac{d^3\vec{p}}{(2\pi)^3} E_{\vec{p}} \, \big( b_{\vec{p}}^{\dagger} \, b_{\vec{p}} + c_{\vec{p}}^{\dagger} \, c_{\vec{p}} \big).
    \ee 
    \textit{Hint: Derive the Hamiltonian density and use \Cref{eqn:ComplexFieldExpansion,eqn:ComplexPiExpansion}. You will need to drop an infinity term as we did for the real scalar field.} 
\ebox 

In our particle interpretation, \Cref{eqn:ComplexFieldHamiltonian} tells us that the complex theory contains two particles of the same energy, and therefore mass. We can show that they also both have vanishing spin. The states of our Hilbert space are made up of products of these two particle types, we use a semi-colon in the bra/ket to indicate the different types. For example 
\bse 
    \ket{\vec{p}_1,\vec{p}_2 ; \vec{q}_1, \vec{q}_2} := b_{\vec{p}_1}^{\dagger}  b_{\vec{p}_2}^{\dagger}  c_{\vec{q}_1}^{\dagger} c_{\vec{q}_2}^{\dagger} \ket{0}.
\ese 
We refer to the particles created by $c$ as \textit{anti-particles} and the particles created by $b$ as \textit{particles}. Note that for the real scalar field we have $b=c$ and so the `particle is its own antiparticle'.

The classical complex field conserved current, \Cref{eqn:ComplexCurrent}, can be upgraded to the quantum theory and from that we can obtain the conserved charge 
\be
\label{eqn:ParticleAntiparticleNumberConserved}
    Q = \int \frac{d^3\vec{p}}{(2\pi)^3} \big( c_{\vec{p}}^{\dagger}\,  c_{\vec{p}} - b_{\vec{p}}^{\dagger}\, b_{\vec{p}}\big) = N_c - N_b,
\ee 
which says that the number of antiparticles minus the number of particles is conserved \textit{locally}. Of course in the free theory we have that $N_c$ and $N_b$ are separatly conserved, so this statement is nothing new. However, as we keep teasing, in the interacting theory the number of each type of particle will no longer be a conserved charge, but \Cref{eqn:ParticleAntiparticleNumberConserved} will still hold. The locality of this result will translate into the fact that particles and antiparticles are always created in pairs, and they can annihilate to produce something different (e.g. a photon). 

\br 
    As we just said, these are spin-0 particles and so do not correspond to Fermions (i.e. electrons etc), but they give us a taste of what's to come. We can refer to this complex scalar field as a `poor man's Fermion'.
\er

\section{Propagators}

At the end of the last lecture we checked to see if our theory was causal by checking that the operators commuted for spacelike separation. We could have worded this question differently, and perhaps more `particle physicsy'. We could have asked "what is the probability for a particle to propagate from point $y$ to point $x$?" If $x$ and $y$ are spacelike separated, we would want the answer to be a big fact zero, otherwise the particles would be travelling faster than light. The propagation is given by the 2-point function $\bra{0}A(x)A(y)\ket{0}$ where $A$ is our field in the Heisenberg picture. So we need to consider this calculation theory by theory. 

\subsection{Real Scalar Field}

First let's consider the real scalar field. Our propagator is 
\bse 
    \begin{split}
        D(x-y) & = \bra{0}\phi(x)\phi(y)\ket{0} \\
        & = \int \frac{d^3 \vec{p}}{(2\pi)^3} \frac{d^3\vec{q}}{(2\pi)^3} \frac{1}{\sqrt{2E_{\vec{p}}}} \frac{1}{\sqrt{2E_{\vec{q}}}} \bra{0} \Big(a_{\vec{p}}\, e^{-ipx} + a_{\vec{p}}^{\dagger} \, e^{ipx} \Big) \Big(a_{\vec{q}}\, e^{-iqy} + a_{\vec{q}}^{\dagger} \, e^{iqy} \Big) \ket{0} \\
        & = \int \frac{d^3 \vec{p}}{(2\pi)^3} \frac{d^3\vec{q}}{(2\pi)^3} \frac{1}{\sqrt{2E_{\vec{p}}}} \frac{1}{\sqrt{2E_{\vec{q}}}} \bra{0} a_{\vec{p}}\,  a_{\vec{q}}^{\dagger} \ket{0} e^{-ipx + i qy} \\
        & = \int \frac{d^3 \vec{p}}{(2\pi)^3} \frac{1}{2E_{\vec{p}}} e^{-ip(x-y)},
    \end{split}
\ese 
where we have used 
\bse 
    \bra{0}a_{\vec{p}}^{\dagger} = 0 = a_{\vec{q}}\ket{0}, \qand \bra{0} a_{\vec{p}}\,  a_{\vec{q}}^{\dagger} \ket{0} = (2\pi)^3 \del^{(3)}(\vec{p}-\vec{q}).
\ese 
Again this result is Lorentz invariant and so we can pick any frame we like. We want to check spacelike separations $(x-y)^2<0$, so choose a frame with $(x-y)=(0,\vec{r})$, and work in polar coordinates, $d^3\vec{p} = p^2 \sin^2\theta d\theta d\varphi dp$. This gives us 
\bse 
    \begin{split}
        D(x-y) & = \frac{2\pi}{(2\pi)^3} \int_0^{\infty} dp \frac{p^2}{2E_p} \frac{e^{ipr} - e^{-ipr}}{ipr}  \\
        & = - \frac{i}{2(2\pi)^2r} \int_{-\infty}^{\infty} dp \frac{p e^{ipr}}{\sqrt{p^2+m^2}} 
    \end{split}
\ese 
This is an integral in the complex plane with two poles at $p=\pm im$. We therefore need to take two branch cuts as indicated in the figure below. We do the integral over the shaded region and extend out to infinity.  

\begin{center}
    \btik 
        \begin{scope}
            \clip (-3,3) -- (-0.3,3) -- (-0.3,1) .. controls (-0.15,0.7) and (0.15,0.7) .. (0.3,1) -- (0.3,3) -- (3,3) -- (3,-3) -- (0.3,-3) -- (0.3,-1) .. controls (0.15,-0.7) and (-0.15,-0.7) .. (-0.3,-1) -- (-0.3,-3) -- (-3,-3) -- (-3,3);
            \draw[fill = gray!40, opacity = 0.8] (-3,3) -- (-0.3,3) -- (-0.3,1) .. controls (-0.15,0.7) and (0.15,0.7) .. (0.3,1) -- (0.3,3) -- (3,3) -- (3,-3) -- (0.3,-3) -- (0.3,-1) .. controls (0.15,-0.7) and (-0.15,-0.7) .. (-0.3,-1) -- (-0.3,-3) -- (-3,-3) -- (-3,3);
        \end{scope}
        \draw[thick, ->] (-3.5,0) -- (3.5,0);
        \node at (3.5,-0.5) {\large{Re($p$)}};
        \draw[thick, decorate, decoration={snake, segment length=1.5mm, amplitude=0.5mm}] (0,1) -- (0,3);
        \draw[fill=black] (0,1) circle [radius=0.07cm];
        \node at (0,0.5) {\large{$im$}};
        \draw[thick, decorate, decoration={snake, segment length=1.5mm, amplitude=0.5mm}] (0,-1) -- (0,-3);
        \draw[fill=black] (0,-1) circle [radius=0.07cm];
        \node at (0,-0.5) {\large{$-im$}};
        % 
        \draw[ultra thick] (-0.3,3) -- (-0.3,1) .. controls (-0.15,0.7) and (0.15,0.7) .. (0.3,1) -- (0.3,3);
        \draw[ultra thick] (-0.3,-3) -- (-0.3,-1) .. controls (-0.15,-0.7) and (0.15,-0.7) .. (0.3,-1) -- (0.3,-3);
    \etik
\end{center}

We then define $\rho = -ip$ and obtain 
\bse 
    D(x-y) = \frac{1}{4\pi^2 r} \int_m^{\infty} d \rho \frac{\rho e^{-\rho r}}{\sqrt{\rho^2-m^2}},
\ese 
which in the limit $r\to\infty$ (which is \textit{very} spacelike separated) gives 
\bse 
    D(x-y) \sim e^{-m|\vec{x}-\vec{y}|},
\ese 
where we've put $r=|\vec{x}-\vec{y}|$ back in. 

Ah this does not appear good... we've shown that the amplitude for a particle to propagate on a spacelike curve decays exponentially, which is small but non-zero. What was the difference to what we did at the end of last lecture? Well last lecture we had the (sandwiched) \textit{commutator}, which in terms of the propagators is 
\bse 
    [\phi(x),\phi(y)] = D(x-y) - D(y-x).
\ese 
But the above formula has $|\vec{x}-\vec{y}|$, so it doesn't change when we switch $x\longleftrightarrow y$. This is why the result vanished last lecture. So what does this correspond to in our particle propagation terms? Well because we're considering a spacelike path, their is no Lorentz invariant way to order events, so we have to consider both the propagation $x\to y$ and $y\to x$, but these have the same amplitude, and so cancel. 

\subsection{Complex Scalar Field}

What about for the scalar field? Well first consider the exercise 
\bbox 
    Define 
    \be
    \label{eqn:ComplexPropagators}
        \begin{split}
            D_b(x-y) & := \bra{0} \psi(x)\psi^{\dagger}(y) \ket{0}, \\
            D_c(x-y) &:= \bra{0} \psi^{\dagger}(x)\psi(y) \ket{0}.
        \end{split}
    \ee 
    Then using the relevant definitions and commutators show these both give 
    \bse 
        \int \frac{d^3\vec{p}}{(2\pi)^3} \frac{1}{2E_{\vec{p}}} e^{-ip(x-y)}
    \ese
    \textit{Hint: The subscripts $b/c$ should become justified relatively quickly.}
\ebox  

For the complex scalar case we therefore have 
\bse 
    \big[\psi(x),\psi^{\dagger}(y)\big] = D_b(x-y) - D_c(y-c),
\ese
and a similar calculation to the real scalar gives that this vanishes for spacelike separations. So we give the same particle propagation tale, but now with an interesting difference. As we see in \Cref{eqn:ComplexPropagators}, $D_b(x,y)$ corresponds to an \textit{antiparticle} propagating from $y\to x$ whereas $D_c(y-x)$ corresponds to a \textit{particle} propagating from $x\to y$. So now the cancellation in amplitude comes from a particle going one way while an antiparticle going the other way! So causality requires that every particle have a corresponding antiparticle with the same mass but opposite quantum numbers. This isn't actually a new interpretation, it's just that for the real scalar field the particle is its own antiparticle so we didn't notice. 

\subsection{Feynman Propagator}

As we shall see, for interacting field theories perhaps the most important object is the so-called \textit{Feynman propagator}:

\mybox{
\be 
\label{eqn:FeynmanPropagator}
    \Delta_F(x-y) = \bra{0}\cT\phi(x)\phi(y)\ket{0} = \begin{cases}
        D(x-y) & x^0 > y^0 \\
        D(y-x) & y^0 > x^0
    \end{cases},
\ee 
}
\noindent where the cases on the right hand side defines $\cT$, the \textit{time ordering} operator:
\be 
\label{eqn:TimeOrdering}
    \begin{split}
        \cT\big(\cO_1(t_1)\cO_2(t_2)\big) & = \cO_1(t_1),\cO_2(t_2) \Theta(t_1-t_2) + \cO_2(t_2),\cO_1(t_1) \Theta(t_2-t_1) \\
        & = \begin{cases}
        \cO_1(t_1),\cO_2(t_2) & t_1 > t_2 \\
        \cO_1(t_2),\cO_2(t_1) & t_2 > t_1
    \end{cases}.
    \end{split}
\ee 

Physically what this tells us is that first the particle is created and then destroyed. It will be incredibly useful for us to write this expression as a $4$-integral, and this comes in the form of \textit{Feynman's trick}. 

\bcl 
    We can write the Feynman propagator as 
    \mybox{
    \be 
    \label{eqn:FeynmanPropagatorIntegral}
        \Delta_F(x-y) = \int \frac{d^4 p}{(2\pi)^4} \frac{i}{p^2-m^2 + i\epsilon} e^{-ip(x-y)},
    \ee 
    }
    where $\epsilon >0$ but infinitesimal.
\ecl 

\bq 
    We prove this claim by taking a complex contour integral and using Cauchy's theorem. First let's rewrite the right-hand side of \Cref{eqn:FeynmanPropagatorIntegral} slightly: 
    \bse
        \begin{split}
            \frac{1}{p^2-m^2+i\epsilon} & = \frac{1}{p_0^2 - \vec{p}^2-m^2 +i\epsilon} \\
            & = \frac{1}{p_0^2 - E_{\vec{p}}^2+i\epsilon} \\
            & = \frac{1}{p_0^2 - (E_{\vec{p}}-i\epsilon)^2} \\
            & = \frac{1}{p_0 - (E_{\vec{p}}-i\epsilon)} \frac{1}{p_0 + (E_{\vec{p}}-i\epsilon)}
        \end{split}
    \ese 
    where to get to the penultimate line we used  the fact that $\epsilon$ is infinitesimal, and relabelled $\epsilon = 2E_{\vec{p}}\epsilon$. Explicitly, we have 
    \bse 
        (E_{\vec{p}} - i\epsilon)^2 = E_{\vec{p}}^2 - 2iE_{\vec{p}}\, \epsilon - \epsilon^2
    \ese 
    then we drop the $\epsilon^2$ term and, because $E_{\vec{p}}>0$, we can redefine $\epsilon\to 2E_{\vec{p}}\epsilon$ which is still positive and infinitesimal. 
    
    We now have an integral with poles 
    \bse 
        p_0^{\pm} =  \pm (E_{\vec{p}}-i\epsilon)
    \ese 
    as indicated in the following diagram. 
    \begin{center}
        \btik 
            \draw[fill = blue, opacity = 0.5] (-3,0) -- (3,0) -- (3,-2) -- (-3,-2) -- (-3,0);
            \draw[fill = red, opacity = 0.5] (-3,0) -- (3,0) -- (3,2) -- (-3,2) -- (-3,0);
            %
            \draw[ultra thick, blue, ->] (-1,0) -- (-0.99,0);
            \draw[ultra thick, blue, ->] (1,0) -- (1.01,0);
            \draw[ultra thick, blue, ->] (3,0) -- (3,-1.5);
            \draw[ultra thick, blue] (3,-1.5) -- (3,-2);
            \draw[ultra thick, blue, ->] (3,-2) -- (1.5,-2);
            \draw[ultra thick, blue, ->] (1.5,-2) -- (-1.5,-2);
            \draw[ultra thick, blue] (-1.5,-2) -- (-3,-2);
            \draw[ultra thick, blue, ->] (-3,-2) -- (-3,-1);
            \draw[ultra thick, blue] (-3,-1) -- (-3,0);
            %
            \draw[ultra thick, red, ->] (-2,0) -- (-1.99,0);
            \draw[ultra thick, red, ->] (2,0) -- (2.01,0);
            \draw[ultra thick, red, ->] (3,0) -- (3,1);
            \draw[ultra thick, red] (3,1) -- (3,2);
            \draw[ultra thick, red, ->] (3,2) -- (1.5,2);
            \draw[ultra thick, red, ->] (1.5,2) -- (-1.5,2);
            \draw[ultra thick, red] (-1.5,2) -- (-3,2);
            \draw[ultra thick, red, ->] (-3,2) -- (-3,1);
            \draw[ultra thick, red] (-3,1) -- (-3,0);
            %
            \draw[thick, ->] (-3.5,0) -- (3.5,0);
            \node at (4.5,0) {\large{Re($p_0$)}};
            \draw[thick, ->] (0,-2.5) -- (0,2.5);
            \node at (0,3) {\large{Im($p_0$)}};
            \draw[white, fill=white] (-2,0.5) circle [radius=0.15cm];
            \draw[white, fill=white] (2,-0.5) circle [radius=0.15cm];
            \draw[fill=black] (-2,0.5) circle [radius=0.07cm];
            \draw[fill=black] (2,-0.5) circle [radius=0.07cm];
            \draw[thick, ->] (3.5,-1) -- (2.2,-0.55);
            \node at (3.8,-1) {\large{$p_0^+$}};
            \draw[thick, ->] (-3.5,1) -- (-2.2,0.55);
            \node at (-3.8,1) {\large{$p_0^-$}};
            \node at (2,1.5) {\large{$y^0>x^0$}};
            \node at (2,-1.5) {\large{$x^0>y^0$}};
        \etik 
    \end{center}
    The shaded regions are meant to indicate how we close our contours for the relevant $(x^0-y^0)$ sign. The arrows are meant to indicate which way round we close each contour (so $x^0>y^0$ is clockwise and $y^0>x^0$ is anticlockwise). We pick up the relative poles as indicated. 
    
    Let's consider the case $x^0>y^0$, then we pick up $p_0^+$ pole, and using Cauchy's theorem, 
    \bse 
        \oint dz \frac{f(z)}{(z-z_0)} = (2\pi i) f(z_0),
    \ese 
    we get the residue $-1/2E_{\vec{p}}$, where the minus sign comes from that fact that we're doing a clockwise integral,\footnote{If this doesn't make sense, basically in Cauchy's theorem you use the anticlockwise contour integral.} and obtain 
    \bse 
        \int \frac{d^4}{(2\pi)^4} \frac{i}{p^2 - m^2 +i\epsilon} e^{-ip(x-y)} = \int \frac{d^3p}{(2\pi)^3} \frac{2\pi i}{2\pi} \frac{i}{2E_{\vec{p}}}e^{-iE_{\vec{p}}(x^0-y^0) + i \vec{p}\cdot (\vec{x}-\vec{y})},
    \ese 
    but the right-hand side (after cancelling) is just $D(x-y)$.
\eq 

\bbox 
    Finish the proof above. That is show 
    \bse 
        \Delta_F(x-y) = D(y-x) \qquad \text{if} \qquad y^0>x^0.
    \ese 
\ebox 

\br 
    If you are not 100\% comfortable with the idea of putting the $i\epsilon$ in the denominator, you can obtain the same result but now the poles like on the real axis. The contour you need to take to get the Feynman propagator is drawn below. 
    \begin{center}
        \btik 
            \draw[thick, ->] (-3.5,0) -- (3.5,0);
            \node at (4.2,0) {\large{Re($p_0$)}};
            \draw[thick, ->] (0,-1.5) -- (0,1.5);
            \node at (-0.7,1.5) {\large{Im($p_0$)}};
            \draw[fill=black] (-2,0) circle [radius=0.07cm];
            \node at (-2,0.5) {\large{$-E_{\vec{p}}$}};
            \draw[fill=black] (2,0) circle [radius=0.07cm];
            \node at (2,-0.5) {\large{$+E_{\vec{p}}$}};
            \draw[ultra thick, blue, ->] (-0.5,0) -- (-0.49,0);
            \draw[ultra thick, blue, ->] (0.5,0) -- (0.51,0);
            \draw[ultra thick, blue] (-3,0) -- (-2.5,0) .. controls (-2.25,-0.5) and (-1.75,-0.5) .. (-1.5,0)-- (1.5,0) .. controls (1.75,0.5) and (2.25,0.5) .. (2.5,0) -- (3,0);
        \etik 
    \end{center}
\er 

\bbox 
    Show that the diagram in the above remark will lead to the same result for the Feynman propagator. \textit{Hint: If you get stuck, this is how Prof. Tong does it. But I advise you try it first using the calculation done in the proof above for guidance.}
\ebox  

The Feynman propagator is, in fact, also a Green's function\footnote{Basically a function that gives a delta function when acted on by a differential operator.} for the Klein-Gordan equation. We see this easily using \Cref{eqn:FeynmanPropagatorIntegral} (with $\epsilon=0$, as we don't need to use a contour integral here):
\bse 
    \begin{split}
        (\p^2 +m^2)\Delta_F(x-y) & = \int \frac{d^4p}{(2\pi)^4} \frac{i}{p^2-m^2} (\p^2 + m^2) e^{-ip(x-y)} \\
        & = \int \frac{d^4p}{(2\pi)^4} \frac{i}{p^2-m^2} \big(-p^2+m^2\big) e^{-ip(x-y)} \\
        & = -i \int \frac{d^4p}{(2\pi)^4} e^{-ip(x-y)} \\
        & = -i\del^{(4)}(x-y).
    \end{split}
\ese
This is a really nice result as it allows us to `invert' the Klein-Gordan equation, i.e. turn it from a differential equation to a integral one.\footnote{For more info on why this is the case, just Google what a Green's function is.}

\section{Building Interacting Theories}

As we have seen, in free theories our equations of motion are linear in the fields, e.g. the Klein-Gordan equation
\bse 
    (\p^2 + m^2)\phi = 0.
\ese 
We solved these theories by a Fourier analysis and saw that all the different modes decoupled. That is we never had a $a_{\vec{p}} \, a_{\vec{q}}$ term in our Hamiltonian etc. This lead to us showing that the number operator in such theories in conserved, \Cref{eqn:NumberOperatorConserved}. This is why we called them free theories.

Obviously the in the real world things do interact, and we have hinted a few times that this will lead to the number operator not being conserved anymore. We need some way, then, to construct interacting theories and the rest of this lecture is dedicated to exactly that. 

We get interactions in our theory when the equations of motion are not linear in the fields. Recalling that the equations of motion come from varying the Lagrangian (i.e. they're the Euler-Lagrange equations). So if we want interaction terms in our equations of motion, we're going to have to add interaction terms into the Lagrangian:
\bse 
    \cL = \cL_{\text{free}} + \cL_{\text{int}}.
\ese
This will also lead to interaction terms in the Hamiltonian, which we want. 

However, in QFT we cannot just arbitrarily change our Lagrangian. That is, we need to keep our `suitable Lagrangian' conditions we introduced previously. We shall state them again here, but now in terms of the additional $\cL_{\text{int}}$:
\ben[label=(\roman*)]
    \item We only want local interactions, in order to preserve causality. That is something of the form 
    \bse 
        \int d^3 y \phi(x)\phi^2(y)
    \ese
    is \textit{not} allowed. 
    \item $\cL_{\text{int}}$ is a Lorentz scalar. 
    \item $\cL_{\text{int}}$ respects the internal symmetries. For example, if our free theory was U(1) invariant (i.e. $\psi \to e^{i\a}\psi$) then we cannot have something like $\psi\psi$ in $\cL_{\text{int}}$.
    \item Renormalisability. We explained this before, but for a reminder, in 4-dimensions, we would have to truncate the following Lagrangian at the $\phi^4$ term, as $[\l_{i>4}]<0$ which is non-renormalisable.
    \bse 
        \cL = \frac{1}{2}(\p\phi)^2 - \frac{1}{2} m^2\phi^2 - \frac{\l_3}{3!} \phi^3 - \frac{\l_4}{4!}\phi^4 - \frac{\l_5}{5!}\phi^5 - ...
    \ese 
    Note here we can see $\cL_{\text{free}}$ as the real scalar field.
\een 

\subsection{Examples}

\bex 
    As an example consider the popular $\phi^4$ theory. This has Lagrangian
    \be 
    \label{eqn:phi4Lagrangian}
        \cL = \underbrace{\frac{1}{2}(\p\phi)^2 - \frac{1}{2}m^2\phi^2}_{\cL_{\text{free}}} - \underbrace{\frac{\l}{4!}\phi^4}_{\cL_{\text{int}}}.
    \ee 
    In terms of Feynman diagrams (to come shortly) the interaction Lagrangian will correspond to something of the form 
    % To anyone reading the code... I tried to use TikzFeynman and FeynMF but neither seemed to work very well, so decided to just do it manually in Tikz myself. See the shortcuts tab for some things I defined.
    \begin{center}
        \btik 
            \draw[thick] (-1,1) -- (0,0);
            \draw[->] (-0.9,0.7) -- (-0.4,0.2);
            \node at (-1.2,1.2) {$p_1$};
            \draw[thick] (-1,-1) -- (0,0);
            \draw[->] (-0.9,-0.7) -- (-0.4,-0.2);
            \node at (-1.2,-1.2) {$p_2$};
            \draw[thick] (0,0) -- (1,1);
            \draw[->] (0.4,0.2) -- (0.9,0.7);
            \node at (1.2,1.2) {$p_4$};
            \draw[thick] (0,0) -- (1,-1);
            \draw[->] (0.4,-0.2) -- (0.9,-0.7);
            \node at (1.2,-1.2) {$p_3$};
            \draw[fill=black] (0,0) circle [radius=0.07cm];
            \node at (0,-0.3) {$\l$};
        \etik  
    \end{center}
\eex 


\bex 
    In the above theory we only had one field $\phi$, but we can have (and we will definitely need) interactions between different fields. As an example, we have the \textit{scalar Yukawa} theory, with Lagrangian 
    \be 
    \label{eqn:ScalarYukawaLagrangian}
        \cL = \underbrace{\frac{1}{2}(\p\phi)^2 - m^2\phi^2}_{\text{Klein-Gordan}} + \underbrace{(\p_{\mu}\psi)(\p^{\mu}\psi^*) - M^2\psi\psi^*}_{\text{Complex Scalar}} - \underbrace{g \psi^*\psi\phi}_{\text{interaction}}.
    \ee 
    This is the interaction between a Klein-Gordan field of mass $m$ and a complex scalar field of mass $M$. 
    \begin{center}
        \btik 
            \midarrow (-1,1) -- (0,0);
            \draw[->] (-0.9,0.7) -- (-0.4,0.2);
            \node at (-0.8,0.3) {$p_1$};
            \node at (-1.2,1.3) {$\psi$};
            \midarrow (0,0) -- (-1,-1);
            \draw[->] (-0.9,-0.7) -- (-0.4,-0.2);
            \node at (-1.2,-1.2) {$\overline{\psi}$};
            \node at (-0.8,-0.25) {$p_2$};
            \draw[thick,dashed] (0,0) -- (1.5,0);
            \node at (1.7,0) {$\phi$};
            \draw[->] (0.4,0.2) -- (1,0.2);
            \node at (0.7,0.5) {$(p_1+p_2)$};
            \draw[fill=black] (0,0) circle [radius=0.07cm];
            \node at (0.15,-0.3) {$g$};
        \etik 
    \end{center}
\eex 

\bbox 
    Find the equations of motion for \Cref{eqn:phi4Lagrangian,eqn:ScalarYukawaLagrangian}. 
\ebox


As the result of the previous exercise shows, our equations of motion are non-linear and so we can no longer solve them by Fourier decomposition. This is exactly what we wanted, and will lead to different modes coupling together. As we will see, we solve these theories in a perturbative way; we assume the couplings $\l$/$g$ are small and expand around the free theory. We are saved from hideously long expression by the famous \textit{Feynman diagrams}. These things really are a field theorist's best friend.\footnote{Sorry Harmonic Oscillators, you've been pipped at the post.}


\chapter{Interaction Picture \& Dyson's Formula}

So far we have used the Schr\"{o}dinger picture and the Heisenberg picture. Recall that the difference between these two pictures was where the time dependence sat: for the Schr\"{o}dinger picture the states are time dependant, whereas in the Heisenberg picture it is the operators that are time dependant. There is a third picture, the \textit{interaction picture},  which is a sort of hybrid of the other two. As the name suggests, it is useful when studying interacting systems. What we do is consider small perturbations from some well-understood Hamiltonian, $H_0$. That is, we write the Hamiltonian as
\bse 
    H = H_0 + H_{\text{int}}.
\ese 
We treat $H_{\text{int}}$ like a Schr\"{o}dinger Hamiltonian, and $H_0$ like a Heisenberg one. That is, the time dependence of the states is governed by $H_{\text{int}}$ while the time dependence of the operators is governed by $H_0$. As we did for the other pictures, we denote the states/operators in the interaction picture with a subscript $I$, and they are given as follows:
\be 
\label{eqn:InteractionPictureStatesAndOperators}
    \ket{\psi(t)}_I = e^{iH_0t}\ket{\psi(t)}_S, \qand \cO_I(t) = e^{iH_0t}\cO_Se^{-iH_0t}.
\ee 
From this we ca work out how the states $\ket{\psi}_I$ depend on time. We simply use the Schr\"{o}dinger equation
\bse 
    i\frac{d}{dt}\ket{\psi}_S = H_S\ket{\psi}_s
\ese 
along with $H_S=(H_0+H_{\text{int}})_S$. We get 
\bse 
    \begin{split}
        i\frac{d}{dt} \Big( e^{-iH_0t}\ket{\psi}_I\Big) & = \big(H_0+H_{\text{int}}\big)_S e^{-iH_0t}\ket{\psi}_I \\
        H_0\Big(e^{-iH_0t}\ket{\psi}_I\Big) + e^{-iH_0t}\bigg(i\frac{d}{dt}\ket{\psi}_I\bigg) & = \big(H_0+H_{\text{int}}\big)_S e^{-iH_0t}\ket{\psi}_I,
    \end{split}
\ese
so we get
\be
\label{eqn:InteractionSchrodigerEquation}
    i\frac{d}{dt}\ket{\psi}_I = (H_{\text{int}})_I\ket{\psi}_I, \qquad \text{with} \qquad (H_{\text{int}})_I := e^{iH_ot} (H_{\text{int}})_S e^{-iH_0t}.
\ee 

\bnn 
    From now on we shall simply write $H_I(t)$ instead of $(H_{\text{int}})_I(t)$ for obvious reasons. However we should be careful not to confuse $H_I(t)$ with $H_{\text{int}}$ itself: the former is the latter in the interaction picture, while the latter is a small perturbation to our Hamiltonian. Note in particular that the former is \textit{always} a function of time whereas the latter is time independent in the Schr\"{o}dinger picture. 
\enn 

\bbox 
    Convince yourself that $[H_I(t_1),H_I(t_2)]\neq0$ in general. 
\ebox 

Ok this looks great but it's useless unless we can actually solve it, so what do we do? The answer is we derive what is known as \textit{Dyson's formula}.

\section{Dyson's Formula}

We start with the ansatz\footnote{As with all good physics derivations do, our guess will turn out to miraculously be correct...}
\bse 
    \ket{\psi(t)}_I = U(t,t_0) \ket{\psi(t_0)}_I,
\ese 
where $U(t,t_0)$ is a \textit{unitary time evolution operator} satisfying 
\bse 
    U(t_1,t_2)U(t_2,t_3) = U(t_1,t_3), \qand U(t,t) = \b1.
\ese 
Using \Cref{eqn:InteractionSchrodigerEquation}, we can easily show 
\be 
\label{eqn:USchrodingerEquation}
    i\frac{d}{dt}U(t,t_0) = H_I(t)U(t,t_0).
\ee 
Now if we were just considering wavefunctions here instead of operators then the solution to this differential equation would just be 
\be
\label{eqn:UWrong}
    U(t,t_0) = \exp\bigg( -i\int_{t_0}^t H_I(t') dt'\bigg).
\ee
However we are considering non-commuting operators, and so we have to be careful about ordering. To be more explicit, consider the Taylor expansion of the above formula, we get 
\bse 
    U(t,t_0) = \b1 - i\int_{t_0}^t H_I(t')dt' + \frac{(-i)^2}{2}\bigg(\int_{t_0}^t H_I(t')dt'\bigg)^2 + ...,
\ese 
and we want the derivative to leave us with 
\bse 
    i\frac{d}{dt}U(t,t_0) = H_I(t) \Bigg[ \b1 - i\int_{t_0}^t H_I(t')dt' + \frac{(-i)^2}{2}\bigg(\int_{t_0}^t H_I(t')dt'\bigg)^2 + ...\Bigg].
\ese 
However if we take the derivative, we get terms like
\bse 
    -\frac{1}{2}\bigg(\int_{t_0}^t H_I(t')dt'\bigg)H_I(t) - \frac{1}{2}H_I(t) \bigg(\int_{t_0}^t H_I(t')dt'\bigg).
\ese 
The second term is the kind of thing we want but the first term has the $H_I(t)$ on the wrong side! We cannot simply `pull it through' as the operators are non-commuting. So we need to alter \Cref{eqn:UWrong} somehow to account for this. The obvious\footnote{Or perhaps only obvious when you know its the answer...} thing to try is taking a time ordering. 

\bcl 
    The solution of \Cref{eqn:USchrodingerEquation} is given by \textit{Dyson's formula}
    \be 
    \label{eqn:DysonsFormula}
        U(t,t_0) = \cT \exp\bigg( -i\int_{t_0}^t H_I(t') dt'\bigg),
    \ee 
    where $\cT$ is the time ordering operator defined above, \Cref{eqn:TimeOrdering}.
\ecl 

\br 
    Before presenting the proof, let's make a quick remark. Recall that if the result of an integral is finite, then we can use\footnote{If you haven't seen this before, it's worth convincing yourself why this is true.}
    \bse 
        \bigg(\int_a^b f(x) dx\bigg)^2 = \int_a^bdx \int_a^b dy f(x)f(y).
    \ese 
    Now in QFT we consider bounded operators, and so the integral over the Hamiltonian is finite, so we can turn our exponential expansion above into a sum of higher order integrals, rather then powers of integrals. This is obviously a \textit{massive} help. In particular, our expansion becomes
    \bse 
        U(t,t_0) = \b1 - i\int_{t_0}^t dt'H_I(t') + \frac{(-i)^2}{2}\Bigg[ \int_{t_0}^t dt' \int_{t'}^t dt'' H_I(t'')H_I(t') + \int_{t_0}^t dt' \int_{t_0}^{t'} dt'' H_I(t')H_I(t'')\Bigg] + ...
    \ese 
    Note the integration limits on the quadratic terms. The first one is the condition that $t''>t'$ as the lower limit for $t''$ is $t'$, while the second one is the condition that $t'>t''$ as the upper limit for $t''$ is $t'$. These two terms are in fact equal as
    \bse
        \begin{split}
            \int_{t_0}^t dt' \int_{t'}^t dt'' H_I(t'')H_I(t') & = \int_{t_0}^t dt'' \int_{t_0}^{t''} dt' H_I(t'')H_I(t') \\
            & = \int_{t_0}^t dt' \int_{t_0}^{t'} dt'' H_I(t')H_I(t''),
        \end{split}
    \ese 
    where the left had side says $t''$ has lower limit $t'$ while the first expression on the right hand side says $t'$ has supper limit $t''$. These are clearly the same statement. The last expression is obtained by simple change of variables $t' \longleftrightarrow t''$. We can also see this diagrammatically as the two following triangle areas: the blue area corresponds to the left hand side above (i.e. $t''>t')$, whereas the red area is the final expression (i.e. $t'>t''$).
    \begin{center}
        \btik 
            \draw[thick, ->] (-0.5,0) -- (4,0);
            \node at (4,-0.3) {$t'$};
            \draw[thick, ->] (0,-0.5) -- (0,4);
            \node at (-0.3,4) {$t''$};
            \draw[fill = blue, opacity = 0.5] (1,1) -- (3,3) -- (1,3) -- (1,1);
            \draw[fill = red, opacity = 0.5] (1,1) -- (3,3) -- (3,1) -- (1,1);
            \draw[thick] (1,1) -- (3,1) -- (3,3) -- (1,3) -- (1,1);
            \draw (1,1) -- (3,3);
            \draw[dashed] (0,1) -- (1,1) -- (1,0);
            \node at (-0.3,1) {$t_0$};
            \node at (1,-0.3) {$t_0$};
            \draw[dashed] (0,3) -- (1,3);
            \node at (-0.3,3) {$t$};
            \draw[dashed] (3,0) -- (3,1);
            \node at (3,-0.3) {$t$};
        \etik 
    \end{center}
    So our expansion becomes 
    \bse 
        U(t,t_0) = \b1 -i \int_{t_0}^t dt'H_I(t') + (-i)^2 \int_{t_0}^tdt'\int_{t_0}^{t'} dt'' H_I(t')H_I(t'') + ...
    \ese 
\er 

\bq 
    Given the final result from the above remark, the proof of Dyson's formula is almost trivial. We simply note that our left most integral is \textit{always} the one with upper limit $t$, and so when we act with the time derivative we are always just getting the latest time result (as the expression is time ordered). We can therefore pull it out of the time ordering and get the result. In terms of maths, that is 
    \bse 
        \begin{split}
            i\frac{d}{dt}\Bigg[ \cT\exp \bigg(-i\int_{t_0}^t H_I(t')dt'\bigg)\Bigg] & = \cT\bigg[H_I(t) \exp \bigg(-i\int_{t_0}^t H_I(t')dt'\bigg)\Bigg] \\
            & = H_I(t)\Bigg[\cT\exp \bigg(-i\int_{t_0}^t H_I(t')dt'\bigg)\Bigg],
        \end{split}
    \ese 
    which is the result we want. 
\eq 

Dyson's formula might look a bit daunting, but when we remember that in the interaction picture we're considering small perturbations, we can truncate our expansion and it becomes a lot nicer. Now, recalling that we said in order to study interacting filed theories we were going to take a perturbation around a free theory, we see why Dyson's formula is so useful. It will allow us to use something called \textit{Wick's theorem} and then define the all-so-useful Feynman diagrams. But first let's talk about scattering. 

\section{Scattering}

\bd
    We define the \textit{scattering matrix}, or \textit{S-matrix} for short, to be the amplitude to go from some initial state $\ket{i}$ to a final state $\ket{f}$:
    \be 
    \label{eqn:SMatrixI}
        S_{fi} := \bra{f}S\ket{i} := \lim_{t_{\pm}\to\infty}\bra{f(t_+)} U(t_+,t_-) \ket{i(t_-)}_I,
    \ee 
    where the states are in the interaction picture
\ed 

When we study scattering in QFT we make one big assumption:
\mybox{
    \begin{center}
        The initial and final states, which we collectively call \textit{asymptotic states}, are eigenstates of the free Hamiltonian. .
    \end{center}
}
\noindent The basic idea is we want to say that the initially the fields/particles are so far apart that they do not interact with each other at all, and so do not feel in the interaction Hamiltonian. As Prof. Tong points out,\footnote{Pages 54-55.} at first this seems like a reasonable thing to do, but then we realise it's actually not as solid an assumption as we might think. We shall illustrate one of the problems here. 

Consider the case of a real Klein-Gordan field, $\phi$, in some potential $V(x)$, the Lagrangian is 
\bse 
    \cL = \cL_{KG} - V(x)\phi,
\ese 
then the equations of motion become 
\bse 
    (\p^2+m^2)\phi = -V(x).
\ese 
So our asymptotic states condition tells us that we need $V(x)\to 0$ as $|x|\to\infty$. This still seems somewhat reasonable. However now let's consider the interaction term to be a $\phi^4$ term. Recall the Lagrangian is 
\bse 
    \cL = \cL_{KG} - \frac{\l}{4!}\phi^4,
\ese
and you should have shown in the exercise that the equation of motion is
\bse 
    (\p^2+m^2)\phi = -\frac{\l}{3!}\phi^3.
\ese 
This poses a more significant problem, as we can't `turn off' our interaction at $|t|\to\infty$; the field is defined everywhere! Despite this we continue on with our assumption that the asymptotic states are eigenstates of the free theory. We therefore rewrite \Cref{eqn:SMatrix} as
\mybox{
    \be 
    \label{eqn:SMatrix}
        S_{fi} := \bra{f}S\ket{i} := \lim_{t_{\pm}\to\infty}\bra{f(t_+)} U(t_+,t_-) \ket{i(t_-)},
    \ee 
}
\noindent where the states are now eigenstates for the free theory.

We can actually define a different matrix, called the \textit{transition matrix}, by 
\bse 
    T_{fi} := \bra{f}(S-\b1)\ket{i},
\ese 
which removes the $\b1$ term in the Taylor expansion of $U(t,t_0)$.

\subsection{Particle Decay}

Ok let's actually derive our first interaction result. Consider scalar Yakawa theory, and recall that the Lagrangian is 
\bse 
    \cL = \frac{1}{2}(\p\phi)^2 -m^2\phi^2 + \p_{\mu}\psi^*\p^{\mu}\psi - M\psi^*\psi - g\psi^*\psi\phi.
\ese 
This has\footnote{This notation is meant to mean everything in the bracket is a function of $x$. Just want to save writing all the brackets and clutterig notation.} 
\bse 
    H_{\text{int}} = g \int d^3x \, \psi^{\dagger}_x\psi_x\phi_x.
\ese 
We're going to `flip' the diagram we gave for this at the end of last lecture, and consider the decay process 
\bse 
    \phi(p) \to \overline{\psi}(q_2)\psi(q_1),
\ese 
where $\phi$ is the particle created by $a^{\dagger}$, $\psi$ the one created by $b^{\dagger}$ and $\overline{\psi}$ the \textit{anti}particle created by $c^{\dagger}$. The diagram\footnote{We will understand how to draw these soon.} looks like 
\begin{center}
    \btik 
        \midarrow (0,0) -- (1,1);
        \draw[->] (0.4,0.2) -- (0.9,0.7);
        \node at (0.8,0.3) {$q_1$};
        \node at (1.2,1.3) {$\psi$};
        \midarrow (1,-1) -- (0,0);
        \draw[->] (0.4,-0.2) -- (0.9,-0.7);
        \node at (1.2,-1.2) {$\overline{\psi}$};
        \node at (0.8,-0.25) {$q_2$};
        \draw[thick,dashed] (-1.5,0) -- (0,0);
        \node at (-1.7,0) {$\phi$};
        \draw[->] (-1,0.2) -- (-0.4,0.2);
        \node at (-0.7,0.5) {$p$};
        \draw[fill=black] (0,0) circle [radius=0.07cm];
        \node at (-0.15,-0.3) {$g$};
    \etik 
\end{center}
Our asymptotic states are given by 
\bse 
    \begin{split}
        \ket{i} & = \sqrt{2E_{\vec{p}}} \, a_{\vec{p}}^{\dagger} \, \ket{0} =: \ket{\phi(p)} \\
        \ket{i} & = \sqrt{4E_{\vec{q}_1}E_{\vec{q}_2}} \,  b_{\vec{q}_1}^{\dagger} \, c_{\vec{q}_2}^{\dagger} \, \ket{0} =: \ket{\psi(q_1)\overline{\psi}(q_2)}
    \end{split}
\ese 
So what is our scattering amplitude? Well we are just considering the first order expansion, i.e. one power of $H_I$ (higher powers would correspond to terms like $\overline{\psi}(p_1)\psi(p_2) \to \phi(q) \to \overline{\psi}(p_3)\psi(p_4)$ etc), so we have
\bse 
    U(-\infty,\infty) = \b1 - ig\int d^4 x \psi^{\dagger}_x\psi_x\phi_x + \cO(g^2),
\ese
and so the transition matrix to leading order in $g$ is 
\bse 
    T_{fi} = ig \bra{f} \int d^4 x \, \psi^{\dagger}_x \psi_x \phi_x \ket{i}. 
\ese 
We now need to substitute in $\psi_x^{\dagger}$, $\psi_x$ and $\phi_x$. First let's just do $\phi_x$. We use the definition 
\bse 
    \phi(x) = \int \frac{d^3\vec{k}}{(2\pi)^3} \frac{1}{\sqrt{2E_{\vec{k}}}} \Big( a_{\vec{k}} \, e^{-ikx} + a^{\dagger}_{\vec{k}} \, e^{ikx}\Big),
\ese 
and note that the $a_{\vec{k}}^{\dagger}$ term will act on $\ket{i}$ and produce a two particle state. That is (up to factors of $\sqrt{2E}$)
\bse 
    a^{\dagger}_{\vec{k}} \ket{i} = a^{\dagger}_{\vec{k}} \, a^{\dagger}_{\vec{p}} \ket{0} = \ket{\phi(p)\phi(k)}.
\ese
Then noting that the $\psi/\psi^{\dagger}$ terms are not going to be able to `undo' this, we will get a zero projection onto the final state $\ket{f}$. So we just need to consider the $a_{\vec{k}}$ term. So using the usual trick of commuting $a_{\vec{k}}$ with the $a_{\vec{p}}^{\dagger}$ and picking up a $(2\pi)^3\del^{(3)}(\vec{k}-\vec{p})$, we get 
\bse 
    T_{fi} = -ig \bra{f} \int d^4x \, \psi^{\dagger}_x\psi_x e^{-ipx}\ket{0}.
\ese 

\bbox 
    Using the definitions for $\psi$ and $\psi^{\dagger}$ show that 
    \bse 
        \bra{f}\psi^{\dagger}_x\psi_x \, ``=" \, e^{i(q_1+q_2)x}\bra{0}.
    \ese 
    \textit{Hint: Again you will argue away terms that will vanish when taking the inner product. This is why I have put ``=''.}
\ebox 

Using the result of the above exercise we get 
\bse 
    T_{fi} = -ig \int d^4 x \, e^{i(q_1+q_2-p)x} = -ig(2\pi)^4\del^{(4)}(p-q_1-q_2).
\ese 
This result is basically just telling us that 4-momentum must be conserved at the vertex, a very nice physical result. What kind of restrictions does this put on the process? Well the result is Lorentz invariant so we can pick the decaying $\phi$s rest frame to evaluate it. In this frame we have $p=(m,0,0,0)$, so the delta function splits into 
\bse 
    \del^{(4)}(p-q_1-q_2) = \del(m-q_1^0-q_2^0)\del^{(3)}(\vec{q}_1+\vec{q}_2),
\ese 
and so we see the decay can only happen if $m \geq 2M$.

Obviously we expect momentum conservation to be something that comes up a lot in our interaction calculations, and indeed it turns out that in general
\be 
\label{eqn:TransitionMatrixMatrixElements}
    T_{fi} = i (2\pi)^4 M_{fi} \, \del^{(4)}(p_i - p_f),
\ee 
where $p_i$ and $p_f$ are the total 4-momentum initial and final state, respectively. So essentially our goal is to find the $M_{fi}$s, and these are given the, somewhat disappointing, name \textit{matrix elements}. We often also group the $i$ with it and call $iM_{fi}$ the matrix elements.

\subsection{$\phi^4$ Scattering}

Ok so we've computed our first scattering matrix, but it corresponded to a decay, so now let's find the first \textit{scattering} scattering matrix. That is we want something like 
\bse 
    F(p_1)G(p_2) \to F(p_3)G(p_4),
\ese 
where $F$ and $G$ are some fields. We could consider higher order terms in the expansion for the scalar Yakawa theory as mentioned above, however as we will see these are incredibly hard to compute and will require the machinery of Wick contractions. So we want something that to first order in $H_I$ will give us 4 fields (two in, two out). The easiest example is of course $\phi^4$ theory, which has 
\bse 
    H_{\text{int}} = \frac{\l}{4!} \int d^3 x \, \phi_x\phi_x\phi_x\phi_x,
\ese 
and corresponds simply to the scattering process 
\bse 
    \phi(p_1)\phi(p_2) \to \phi(p_3)\phi(p_4).
\ese 
The calculation of the transition matrix is left as an exercise. 

\bbox 
    Show that the above interaction Hamiltonian leads to 
    \bse 
        T_{fi} = -i\l \del^{(4)}(p_1+p_2-p_3-p_4).
    \ese 
\ebox  

Note that the factor of $4!$ has gone in the result to the above exercise. This is actually the reason it's included in the original expression, so that the matrix elements don't have any additional factors. Its an example of a what we call \textit{symmetry factors}, which will become more clear after studying Wick contractions. 

\br 
    Note that for the scalar Yakawa decay we could have actually used the S-matrix instead of the transition matrix. This is because the two only differ by the inner product between the initial and final state, which for the scalar Yakawa theory vanishes. However for the $\phi^4$ scattering, this inner product is non-vanishing and doesn't actually correspond to a scattering. Diagrammatically it looks like,
    \begin{center}
        \btik 
            \draw[thick] (-2,0.75) -- (2,0.75);
            \node at (-2.3,0.75) {$\phi$};
            \node at (2.3,0.75) {$\phi$};
            \draw[->] (-1,1) -- (1,1) node [midway, above] {$p_1=p_3$};
            \draw[thick] (-2,-0.75) -- (2,-0.75);
            \node at (-2.3,-0.75) {$\phi$};
            \node at (2.3,-0.75) {$\phi$};
            \draw[->] (-1,-1) -- (1,-1) node [midway, below] {$p_2=p_4$};
        \etik 
    \end{center}
    It is therefore important to actually distinguish between $S_{fi}$ and $T_{fi}$, however it is often the case that we write $S_{fi}$ and really mean $T_{fi}$. In these notes I shall try be careful here, but its quite possible that I will forget and make a mistake, so keep on your toes. 
\er 


\chapter{Wick's Theorem}

Last lecture we showed how we could use Dyson's formula to find the S-matrix for a simple decay and simple scattering. We made a comment about how would could have got a scattering in the scalar Yakawa theory, but that it would involve a quadratic term from Dyson's formula, and so would be very hard to do in the manner above. However, obviously we want to be able to calculate such terms (otherwise we never pass first order in perturbation theory, and we loose \textit{a lot} of information), so how do we do it? The answer is we be smart about how we write things. 

\section{Normal Ordering}

Ok we just said we "be smart", but what does that actually entail? We shall use the scalar Klein-Gordan field to explain. Recall that we can write this field as the expansion 
\bse 
    \phi(x) = \int \frac{d^3\Vec{p}}{(2\pi)^3} \frac{1}{\sqrt{2E_{\Vec{p}}}} \Big( a_{\vec{p}}\, e^{-ipx} + a_{\vec{p}}^{\dagger} \, e^{ipx} \Big).
\ese 
Well let's define 
\be 
\label{eqn:phi+-}
    \begin{split}
        \phi^+(x) & = \int \frac{d^3\Vec{p}}{(2\pi)^3} \frac{1}{\sqrt{2E_{\Vec{p}}}} a_{\vec{p}}\, e^{-ipx} \\
        \phi^-(x) & = \int \frac{d^3\Vec{p}}{(2\pi)^3} \frac{1}{\sqrt{2E_{\Vec{p}}}} a_{\vec{p}}^{\dagger} \, e^{ipx},
    \end{split}
\ee 
so that we simply have 
\bse 
    \phi(x) = \phi^+(x) + \phi^-(x).
\ese 

\br 
    It might seem strange that we define $\phi^+(x)$ to be the one that has the negative exponential. The reason for this is that $e^{-ipx}$ corresponds to a mode with positive frequency, and so we use $\phi^+(x)$. Also note that the dagger does \textit{not} come with $\phi^+(x)$, as you might think. This is actually important because it is fact that $\phi^{\pm}(x)$ correspond to annihilation/creation operators, respectively, that we will use. 
\er 

Now, as we just mentioned in the remark above, $\phi^+(x)$ comes with an annihilation operator, so if it were to appear to the far right in a vacuum expectation value, the result would be zero. The same thing is true for $\phi^-(x)$ being to the far left. That is 
\bse 
    \phi^+(x)\ket{0} = 0 = \bra{0}\phi^-(x).
\ese 
Similarly, we can see that 
\bse 
    \phi^+(x)\ket{\vec{p}} = e^{-ipx} \ket{0}, \qand \bra{\vec{p}}\phi^-(x) = e^{ipx}\bra{0}.
\ese 

We call the procedure of putting all the annihilation operators to the right \textit{normal ordering}. 

\mybox{
    \bd[Normal Ordering]
        We denote the \textit{normal ordered} string of operators $\cO_1(x_1)...\cO_n(x_n)$ as 
        \be 
        \label{eqn:NormalOrdering}
            \cl \cO_1(x_1) ... \cO_n(x_n) \cl,
        \ee 
        which tells us to put all the annihilation operators to the right.
    \ed 
}

\br
    Normal ordering is `linear' in the sense that 
    \bse 
        \cl AB + CD \cl = \cl AB \cl + \cl CD \cl.
    \ese 
\er 

\br 
    Note we could have introduced this (as Prof. Tong does) when talking about the energy of the ground state being infinite. That is we could write the Hamiltonian for the single free scalar field as 
    \bse 
        \cl H \cl = \int \frac{d^3\vec{p}}{(2\pi)^3} E_{\vec{p}} a_{\vec{p}}^{\dagger} \, a_{\vec{p}},
    \ese 
    so that
    \bse 
        \cl H \cl \ket{0} = 0.
    \ese 
\er 

\bnn 
    The colon notation used here for normal ordering is very common, however some authors, including Peskin and Schroeder, denote normal ordering by an $N$. So \Cref{eqn:NormalOrdering} would be written $N(\cO_1(x_1)...\cO_n(x_n))$. 
\enn 

\section{Wick's Theorem}

Ok normal ordering seems great, but Dyson's formula comes in the form of time ordered products, not normal ordered ones. So the question becomes "can we convert one into the other?" Well let's look at an example and see. 

\subsection{Recovering The Propagator}

Let's use \Cref{eqn:phi+-} to try investigate the link. Let's start with the simple case of the product of two fields. Let's define 
\bse 
    \phi^{\pm}(x_i) = \phi^{\pm}_i,
\ese 
and let's consider $x_1^0>x_2^0$. Then the time ordered product is\footnote{Be careful not to confuse the colon in $:=$, which means "is defined to be equal to" to the colons for normal ordering. I considered using another type font for the colon, but decided against it.}
\bse 
    \begin{split}
        \cT\big(\phi_1\phi_2\big) & :=  \phi_1\phi_2 \\
        & = (\phi_1^+\phi_1^-) (\phi_2^+\phi_2^-) \\
        & = \phi_1^+\phi_2^+ + \phi_1^+\phi_2^- + \phi_1^-\phi_2^+ + \phi_1^-\phi_2^- \\
        & = \phi_1^+\phi_2^+ + \phi_2^-\phi_1^+ + [\phi_1^+,\phi_2^-] + \phi_1^-\phi_2^+ + \phi_1^-\phi_2^- \\
        & = \cl \phi_1\phi_2 \cl + [\phi_1^+,\phi_2^-] \\
        & = \cl \phi_1\phi_2\cl + D(x_1-x_2),
    \end{split}
\ese 
where we have used the commutator to make every term have a $\phi^+$ to the far right and the definition of the propagator $D(x-y)$.

\bbox 
    Following a similar method to the above, show that if $x_2^0 > x_1^0$, then
    \bse 
        \cT(\phi_1\phi_2) = \cl \phi_1\phi_2 \cl + D(x_2-x_1).
    \ese
\ebox

So together with the result of the above exercise, we conclude that 
\bse 
    \cT(\phi_1\phi_2) = \cl \phi_1\phi_2 \cl + \begin{cases}
        D(x_1-x_2) & \text{if } x_1^0>x_2^0 \\
        D(x_2-x_1) & \text{if } x_2^0>x_1^0
    \end{cases}.
\ese 
Hang on, this looks familiar... the cases term is just the Feynman propagator, \Cref{eqn:FeynmanPropagator}. So we get 
\be
\label{eqn:TimeOrderdNormalOrderedTwoRealFields}
    \cT(\phi_1\phi_2) = \cl \phi_1\phi_2\cl + \Delta_F(x_1-x_2).
\ee 

\br 
    Note that once again we have an expression that mixes operators and $\C$-numbers: both $\cT(\phi_1\phi_2)$ and $\cl \phi_1\phi_2 \cl$ are operators, but the Feynman propagator is a $\C$-number. We have made this comment several times throughout this course, but make it here again to keep us on our toes.
\er

\subsection{Wick Contractions \& Wick's Theorem}

So we obtained a really nice expression relating the time ordering and normal ordering for two real scalar fields, and their difference was the Feynman propagator. We now want to try and generalise this result to a general relation between time ordering and normal ordering. In order to do that, it is convenient to introduce a notation for \textit{Wick contractions}. 

\mybox{
    \bd[Wick Contraction]
        We define the \textit{Wick contraction}, or just \textit{contraction}, of a pair of fields, $\phi_i$ and $\phi_j$, in a string of $n$-operators, $\phi_1...\phi_i... \phi_j...\phi_n$, to mean "replace those operators with the Feynman propagator, and leave everything else alone." We use the following notation.
        \be 
        \label{eqn:WickContraction}
            \contraction{\phi_1...}{\phi_i}{...}{\phi_j}
            \phi_1...\phi_i...\phi_j...\phi_n = \Delta_F(x_i-x_j) \phi_1...\phi_n,
        \ee 
        where the $...$ on the right hand side does \textit{not} include $\phi_i$ and $\phi_j$.\footnote{If you like, in standard "this element is missing" notation, we could write the right hand side as $\Delta_F(x_i-x_j) \phi_1...\hat{\phi_i}...,\hat{\phi_j}...\phi_n$, but I thought this would be confusing, as hats are used for operators in QFT.} 
    \ed 
}
\bex 
    So we can write \Cref{eqn:TimeOrderdNormalOrderedTwoRealFields} simply as 
    \bse
        \cT(\phi_1\phi_2) = \cl \phi_1\phi_2\cl +  \contraction{}{\phi_1}{}{\phi_2} \phi_1\phi_2.
    \ese 
\eex

\bcl 
    A similar argument/proof can be made for complex fields, where we have
    \bse 
        \contraction{}{\psi_1}{}{\psi_2^{\dagger}} \psi_1\psi_2^{\dagger} = \Delta_F(x_1-x_2), \qand \contraction{}{\psi_1}{}{\psi_2} \psi_1\psi_2 = 0 = \contraction{}{\psi_1^{\dagger}}{}{\psi_2^{\dagger}} \psi_1^{\dagger}\psi_2^{\dagger},
    \ese 
    so that 
    \bse 
        \cT(\psi_1\psi_2^{\dagger}) = \cl \psi_a\psi^{\dagger}_2 \cl + \contraction{}{\psi_1}{}{\psi_2^{\dagger}} \psi_1\psi_2^{\dagger}.
    \ese
\ecl 

So we see that for the case of two real scalar fields, and claimed it for two complex scalar field, the difference between the time ordered and normal ordered products is a Wick contraction, but what about if we have more than two fields? Well this is where Wick's theorem comes in. 

\mybox{
    \bt[Wick's Theorem] 
       For any collection of fields $\phi_1,...,\phi_n$ the following holds 
       \be 
        \cT(\phi_1...\phi_n) = \cl \phi_1...\phi_n \cl + \cl \text{all possible contractions} \cl.
       \ee 
    \et
}

\bq 
    We do not prove Wick's theorem here, but an explanation of the proof (its an inductive proof) is given in Prof. Tong's notes, page 58.
\eq 

It is often very useful to break a calculation using Wick's theorem into categories, given by the number of contractions done. Every calculation will have the $0$-contraction term, which is just the fully normal ordered term, and then you consider the terms in increasing number of contractions.\footnote{Well it depends what you want to do. In string theory it is often the term which is fully contracted that contains most of the useful information. For more details, see my notes on Dr. Shiraz Minwalla's string theory course.}

\bex
    Let's consider 4 real Klein-Gordan fields. Then we can have either $0$-contractions, $1$-contraction or $2$-contractions. The respective terms are:
    \begin{itemize}
        \item $0$-contractions
            \bse 
                \cl \phi_1\phi_2\phi_3\phi_4 \cl, 
            \ese 
        \item $1$-contraction:
            \bse 
                \cl \contraction{}{\phi_1}{}{\phi_2} \phi_1\phi_2\phi_3\phi_4 + \contraction{}{\phi_1}{\phi_2}{\phi_3} \phi_1\phi_2\phi_3\phi_4 + \contraction{}{\phi_1}{\phi_2\phi_3}{\phi_4} \phi_1\phi_2\phi_3\phi_4 + \contraction{\phi_1}{\phi_2}{}{\phi_3} \phi_1\phi_2\phi_3\phi_4 + \contraction{\phi_1}{\phi_2}{\phi_3}{\phi_4} \phi_1\phi_2\phi_3\phi_4 + \contraction{\phi_1\phi_2}{\phi_3}{}{\phi_4} \phi_1\phi_2\phi_3\phi_4 \cl 
            \ese 
        \item $2$-contractions
            \bse 
                \cl \contraction{}{\phi_1}{}{\phi_2} \contraction{\phi_1\phi_2}{\phi_3}{}{\phi_4} \phi_1\phi_2\phi_3\phi_4 + \contraction{}{\phi_1}{\phi_2}{\phi_3} \contraction[2ex]{\phi_1}{\phi_2}{\phi_3}{\phi_4} \phi_1\phi_2\phi_3\phi_4 + \contraction[2ex]{}{\phi_1}{\phi_2\phi_3}{\phi_4} \contraction{\phi_1}{\phi_2}{}{\phi_3} \phi_1\phi_2\phi_3\phi_4 \cl 
            \ese 
    \end{itemize}
    This result implies that 
    \bse 
        \bra{0}\cT(\phi_1\phi_2\phi_3\phi_4)\ket{0} = \Delta_F(x_1-x_2)\Delta_F(x_3-x_4) + \Delta_F(x_1-x_3)\Delta_F(x_2-x_4) + \Delta_F(x_1-x_4)\Delta_F(x_2-x_3),
    \ese 
    which is a Green's function for the Klein-Gordan operator (recall that the Feynman integral is a Green's function).
\eex 

\subsection{Wick's Theorem For Matrix Elements}

Wick's theorem is defined for time ordered products of operators, but we can adapt it to allow us to contract with states. We do this by defining 
\mybox{
    \be 
    \label{eqn:WickContractionOnStates}
        \contraction{}{\phi(x)}{}{\ket{\vec{p}}} \phi(x)\ket{\vec{p}} = e^{-ipx}, \qand \contraction{}{\bra{\vec{p}}}{}{\phi(x)} \bra{\vec{p}}\phi(x) = e^{ipx},
    \ee 
}
\noindent and then Wick's theorem is simply adapted to include contractions between fields and states. This allows us to calculate the S-matrix completely using Wick's theorem. 

\section{Scalar Yakawa Scattering}

We can now return to the problem of trying to find the $\cO(g^2)$ term in scalar Yakawa theory. Recall the Lagrangian is,\footnote{Using $\psi^\dagger$ now instead of $\psi^*$ because we will only consider the quantum theory from now on.}
\bse 
    \cL = \frac{1}{2}(\p\phi)^2 - \frac{1}{2}m^2\phi^2 + \frac{1}{2}\p_{\mu}\psi^{\dagger}\p^{\mu}\psi - M^2\psi^{\dagger}\psi - g\phi\psi^{\dagger}\psi. 
\ese 
Similarly to \Cref{eqn:phi+-} we define 
\bse 
    \begin{split}
        \phi_b^+(x) & = \int \frac{d^3 \vec{p}}{(2\pi)^3} \frac{1}{\sqrt{2E_{\vec{p}}}} b_{\vec{p}} \, e^{-ipx}, \qand \phi_b^-(x) = \int \frac{d^3 \vec{p}}{(2\pi)^3} \frac{1}{\sqrt{2E_{\vec{p}}}} b_{\vec{p}}^{\dagger} \, e^{ipx} \\
        \phi_c^+(x) & = \int \frac{d^3 \vec{p}}{(2\pi)^3} \frac{1}{\sqrt{2E_{\vec{p}}}} c_{\vec{p}} \, e^{-ipx} \qand \phi_c^-(x) = \int \frac{d^3 \vec{p}}{(2\pi)^3} \frac{1}{\sqrt{2E_{\vec{p}}}} c_{\vec{p}}^{\dagger} \, e^{ipx}
    \end{split}
\ese 
so that 
\bse 
    \psi(x) = \phi_b^+(x) + \phi_c^-(x), \qand \psi^{\dagger}(x) = \phi_c^+(x) + \phi_b^-(x),
\ese 
where we note that each term contains a $b$ and a $c$, in correspondance with \Cref{eqn:ComplexFieldExpansion}. Then, we recalling the definitions 
\bse 
    \ket{\psi(p)} := \sqrt{2E_{\vec{p}}} \, b^{\dagger}_{\vec{p}}\ket{0}, \qand \ket{\overline{\psi}(p)} := \sqrt{2E_{\vec{p}}} \, c^{\dagger}_{\vec{p}}\ket{0},
\ese 
our Wick contractions are
\bse 
    \contraction{}{\psi(x)}{}{\ket{\psi(p)}} \psi(x)\ket{\psi(p)} = \contraction{}{\psi^{\dagger}(x)}{}{\ket{\overline{\psi}(p)}} \psi^{\dagger}(x)\ket{\overline{\psi}(p)} = e^{-ipx}, \qand \contraction{}{\bra{\psi(p)}}{}{\psi^{\dagger}(x)} \bra{\psi(p)}\psi^{\dagger}(x) = \contraction{}{\bra{\overline{\psi}(p)}}{}{\psi(x)} \bra{\overline{\psi}(p)}\psi(x) = e^{ipx},
\ese 
and all others vanishing. 

So what we're looking at is $\psi\psi \to \phi \to \psi \psi$, so our initial and final states are 
\bse 
    \ket{i} = \ket{\psi(p_1)\psi(p_2)}, \qand \bra{f} = \bra{\overline{\psi}(p_3)\overline{\psi}(p_4)}.
\ese 
The transition matrix is then 
\bse
    \begin{split}
        T_{fi} = & -ig \bra{\overline{\psi}(p_3)\overline{\psi}(p_4)} \int d^4 x \, \psi^{\dagger}_x\psi_x\phi_x \ket{\psi(p_1)\psi(p_2)} \\
        & + \frac{(-ig)^2}{2!} \bra{\overline{\psi}(p_3)\overline{\psi}(p_4)} \cT \bigg( \int d^4x \, d^4y \,  \psi^{\dagger}_x\psi_x \phi_x \psi^{\dagger}_y\psi_y\phi_y \bigg) \ket{\psi(p_1)\psi(p_2)} + \cO(g^3).
    \end{split}
\ese 
We can use Wick's theorem to evaluate this. First we note that the first term vanishes. Why? Well because we have nothing to contract the $\phi_x$ with, so we have to leave it as a normal ordered term, which then acts on the vacuum giving a vanishing result. We will see in terms of Feynman diagrams soon why such a term can't exist. Now what about the second term? Well we get three different types of contractions which are conceptually different they are 
\ben[label=(\roman*)]
    \item Fully Connected,
    \item Connected, and 
    \item Disconnected.
\een 

We shall consider the fully connected diagrams in this lecture and derive the corresponding Feynman rules, and then talk about connected and disconnected diagrams next lecture.\footnote{Just because that's what happened on the course.}

\subsection{Fully Connected}

There are two unique types of fully connected contractions for the term above, they are 
\bse 
    -\frac{g^2}{2!} \int d^4x \, d^4y \,  \contraction{}{\overline{\psi}(p_3)}{\overline{\psi}(p_4)\cl \,\,\, \psi_x}{psi^{\dagger}_y} \bcontraction[2ex]{\psi(p_3)}{\overline{\psi}(p_4)}{ \cl\quad }{\psi^{\dagger}_x} \bra{\overline{\psi}(p_3)\overline{\psi}(p_4)} \cl  \psi^{\dagger}_x \psi^{\dagger}_y \contraction{}{\phi_x}{}{\phi_y} \phi_x\phi_y \contraction{}{\psi_x}{\psi_y\cl }{\ket{\psi(p)}} \bcontraction[2ex]{\psi_x}{\psi_y}{\cl\psi(p) \,}{\psi(p)} \psi_x\psi_y \cl  \ket{\psi(p_1)\psi(p_2)}  + (x\longleftrightarrow y),
\ese 
and 
\bse 
    -\frac{g^2}{2!} \int d^4x \, d^4y \,  \contraction{}{\overline{\psi}(p_3)}{\ket{\overline{\psi}(p_4)}}{psi^{\dagger}_x} \bcontraction[2ex]{\psi(p_3)}{\overline{\psi}(p_4)}{ \cl \psi_x \quad }{\psi^{\dagger}_y} \bra{\overline{\psi}(p_3)\overline{\psi}(p_4)} \cl  \psi^{\dagger}_x \psi^{\dagger}_y \contraction{}{\phi_x}{}{\phi_y} \phi_x\phi_y \contraction{}{\psi_x}{\psi_y\cl }{\ket{\psi(p)}} \bcontraction[2ex]{\psi_x}{\psi_y}{\cl\psi(p) \,}{\psi(p)} \psi_x\psi_y \cl  \ket{\psi(p_1)\psi(p_2)}  + (x\longleftrightarrow y).
\ese 
The $(x\longleftrightarrow y)$ terms are symmetries factors that will remove the $1/2$ factors, just like ones we mentioned before that got rid of the $1/4$ in $\phi^4$ theory. We will see soon in terms of Feynman diagrams why they are symmetry factors.

Evaluating the contractions using the rules we have above, these two terms become, respectively
\be 
\label{eqn:ScalarYakawaFullyConnectedA}
    \int d^4 x \, d^4 y \, \Big[ e^{i(p_3-p_1)x} \,  (-ig) \Delta_F^{\phi}(x-y) (-ig) \, e^{i(p_4-p_2)y} \Big],
\ee 
and 
\be 
\label{eqn:ScalarYakawaFullyConnectedB}
    \int d^4 x \, d^4 y \, \Big[ e^{i(p_4-p_1)x} \,  (-ig) \Delta_F^{\phi}(x-y) (-ig) \, e^{i(p_3-p_2)y} \Big].
\ee 
 

\bbox 
    Check that the Wick contractions do indeed result in the above two expressions.
\ebox 

\subsection{Feynman Rules --- Position Space}

We are finally at a place to draw a Feynman diagram and derive the corresponding Feynman rules. Really its a flip of a coin which you do first (draw the diagram or define the rules), here I've decided to present the rules and then use them to draw the diagram as I think this is perhaps more instructive. 

The Feynman rules are a set of instructions on how to draw the different terms in the transition matrix, so that you can manipulate the expressions diagrammatically, and then at the end convert them back to equations should you need to. 

\mybox{
The \textit{position space} Feynman rules for scalar Yakawa scattering $\psi\psi\to \phi \to \psi\psi$ are:
\begin{center}
	\begin{tabular}{@{} C{4cm} C{4cm} C{4cm} @{}}
		\toprule
		 Type & Diagram & Math Expression \\
		\midrule 
		Incoming Particle & \btik 
            \midarrow (0,0) -- (2,0);
            \draw[->] (0.5,0.2) -- (1.5,0.2) node [midway, above] {$p$};
            \draw[fill=black] (2,0) circle [radius=0.07];
            \node at (2.3,0) {$x$};
        \etik & $e^{-ipx} = \contraction{}{\psi_x}{}{\ket{\vec{p}}} \psi_x\ket{\vec{p}}$ \\
        Outgoing Particle & \btik 
            \midarrow (0,0) -- (2,0);
            \draw[->] (0.5,0.2) -- (1.5,0.2) node [midway, above] {$p$};
            \draw[fill=black] (0,0) circle [radius=0.07];
            \node at (-0.3,0) {$x$};
        \etik & $e^{ipx} = \contraction{}{\bra{\vec{p}}}{}{\psi_x} \bra{\vec{p}}\psi_x$ \\ \\
        Propagator & \btik 
            \draw[thick, dashed] (0,0) -- (2,0);
            %\draw[->] (0.5,0.2) -- (1.5,0.2) node [midway, above] {$k$};
            \draw[fill=black] (0,0) circle [radius=0.07];
            \draw[fill=black] (2,0) circle [radius=0.07];
            \node at (-0.3,0) {$x$};
            \node at (2.3,0) {$y$};
        \etik & $\Delta_F^{\phi}(x-y) = \contraction{}{\phi_x}{}{\phi_y} \phi_x\phi_y$ \\
        Vertex & \btik 
            \midarrow (0,0) -- (1,0);
            \midarrow (1,0) -- (2,0);
            \draw[thick, dashed] (1,0) -- (1,-0.8);
            \draw[fill=black] (1,0) circle [radius=0.07];
            \node at (1,0.3) {$x$};
        \etik & $ig\int d^4x$ \\
		\bottomrule
	\end{tabular}
\end{center}
}
\noindent Let's clarify the notation/terminology of the rules.
\begin{itemize}
    \item We take time to run from left to right. So the initial state is to the far left and final state is to the far right. Note, though, that some authors choose to have time run vertically. Peskin and Schroeder use this latter convention.
    \item We use solid lines to indicate complex fields (or nucleons) and dotted lines to indicate real fields (or mesons).
    \item We use a filled in circle to indicate a vertex, and we label it by its integration variable.
    \item The arrow above the lines indicate the value of the 4-momentum and the direction it flows.
    \item The arrows on the nucleon line tells us about the flow of the charge (e.g. Fermion number), if the line points in the same direction as the momentum arrow we have a particle, and if they point oppositely, then we have an antiparticle. This is why we don't draw an arrow on the dotted line; a real scalar field is its own antiparticle. 
    \item The arrows for a charge must always be conserved at a vertex. By this we mean one must always go in and one must always leave. 
    \item  The incoming and outgoing particles are known as \textit{external line} because they only connect to one vertex, whereas the propagator is called an \textit{internal line} because it is connected to two vertices. 
    \item We usually denote the 4-momentum on the external lines with $p_i$, and in the position space picture we don't need to label the momentum of the internal lines.  
    \item Sometimes we also include the symbols of the particles next to the lines. This is particularly the case when we have multiple different species particle of the same type (e.g. electrons and muons) in the same diagram.
\end{itemize}

\noindent Using these rules, we can construct the Feynman diagrams for \Cref{eqn:ScalarYakawaFullyConnectedA,eqn:ScalarYakawaFullyConnectedB}. They are, respectively:
\begin{center}
    \btik 
        \begin{scope}[xshift=-3.5cm]
            \midarrow (-2,0) -- (0,0);
            \draw[->] (-1.5,0.2) -- (-0.5,0.2) node [midway, above] {$p_1$};
            \draw[fill=black] (0,0) circle [radius=0.07] node [above] {$x$};
            \midarrow (0,0) -- (2,0);
            \draw[->] (0.5,0.2) -- (1.5,0.2) node [midway, above] {$p_3$};
            \draw[thick, dashed] (0,0) -- (0,-2);
            \midarrow (-2,-2) -- (0,-2);
            \draw[->] (-1.5,-2.2) -- (-0.5,-2.2) node [midway, below] {$p_2$};
            \draw[fill=black] (0,-2) circle [radius=0.07] node [below] {$y$};
            \midarrow (0,-2) -- (2,-2);
            \draw[->] (0.5,-2.2) -- (1.5,-2.2) node [midway, below] {$p_4$};
            \node at (-2.3,0) {$\psi$};
            \node at (-2.3,-2) {$\psi$};
            \node at (2.3,0) {$\psi$};
            \node at (2.3,-2) {$\psi$};
            \node at (-0.3,-1) {$\phi$};
        \end{scope}
        \begin{scope}[xshift=3.5cm]
            \midarrow (-2,0) -- (0,0);
            \draw[->] (-1.5,0.2) -- (-0.5,0.2) node [midway, above] {$p_1$};
            \draw[fill=black] (0,0) circle [radius=0.07] node [above] {$x$};
            \aftermidarrow (0,0) -- (2,-2);
            \draw[->, rotate around={-45:(0,0)}] (0.2,0.2) -- (1.2,0.2);
            \node at (0.9,-0.2) {$p_3$};
            \draw[thick, dashed] (0,0) -- (0,-2);
            \midarrow (-2,-2) -- (0,-2);
            \draw[->] (-1.5,-2.2) -- (-0.5,-2.2) node [midway, below] {$p_2$};
            \draw[fill=black] (0,-2) circle [radius=0.07] node [below] {$y$};
            \aftermidarrow (0,-2) -- (2,0);
            \draw[->, rotate around={45:(0,-2)}] (0.2,-2.2) -- (1.2,-2.2);
            \node at (0.9,-1.8) {$p_4$};
            \node at (-2.3,0) {$\psi$};
            \node at (-2.3,-2) {$\psi$};
            \node at (2.3,0) {$\psi$};
            \node at (2.3,-2) {$\psi$};
            \node at (-0.3,-1) {$\phi$};
        \end{scope}
    \etik 
\end{center}

\br
    The symmetry factor $x\longleftrightarrow y$ in terms of the Feynman diagrams is seen by the fact that if we relabel $x\longleftrightarrow y$ then the diagram looks exactly the same apart from the relabelling. So there is a symmetry in the diagram, hence the name. 
\er 

\br 
    Note that we can see from the Feynman diagrams that the order $g$ term wont contribute, as we said before. That's just because there is no way to start with two $\psi$ lines and end up with two $\psi$ lines using only one vertex.\footnote{If you don't see this, try and do it and you'll quickly see what I mean.} This illustrates one of the incredible powers of Feynman diagrams: it allows us to see almost instantly which terms do not contribute to the transition matrix. 
\er 

\bbox 
    Convince yourself that these diagrams are correct. \textit{Hint: This exercise should make it clear why we collected the terms in \Cref{eqn:ScalarYakawaFullyConnectedA,eqn:ScalarYakawaFullyConnectedB}  ad we did.}
\ebox 

\subsection{Feynman Rules --- Momentum Space}

As it says in the box above, the above rules are in position space (our integrals are over $x$ and $y$), but the Feynman propagator is given by \Cref{eqn:FeynmanPropagatorIntegral},
\bse 
    \Delta_F(x-y) = \int \frac{d^4 p}{(2\pi)^4} \frac{i e^{-ip(x-y)}}{p^2-m^2 + i\epsilon} ,
\ese 
and so we can replace the position integrals with momentum ones. This gives us the \textit{momentum space} Feynman rules. Let's do this explicitly for \Cref{eqn:ScalarYakawaFullyConnectedA}. Plugging the Feynman propagator in, and using $k$ as the momentum variable in correspondance with the rules above, we have 
\bse 
    \int d^4 x \, d^4 y \, \Big[ e^{i(p_3-p_1)x} \,  (-ig) \int \frac{d^4 k}{(2\pi)^4} \frac{i e^{-ik(x-y)}}{k^2-m^2 + i\epsilon} (-ig) \, e^{i(p_4-p_2)y} \Big],
\ese 
combining the $x$ and $y$ exponentials and then integrating over $x$ and $y$ gives us two delta functions, that is we get 
\bse 
    \int \frac{d^4k}{(2\pi)^4} (2\pi)^4 \del^{(4)}\big(p_1 - p_3 - k\big) (-ig)\frac{i}{k^2-m^2+i\epsilon}(-ig) (2\pi)^4 \del^{(4)}\big(p_2 - p_4 + k\big),
\ese 
and then integrating over $k$ leaves us with\footnote{Note we could also use $(p_2-p_4)^2$ in the denominator.} 
\bse 
    (2\pi)^4 \del^{(4)}\big(p_1+p_2 - p_3-p_4\big) \bigg[ (-ig) \frac{i}{(p_1-p_3)^2 - m^2 +i\epsilon}(-ig) \bigg],
\ese 
which, comparing to \Cref{eqn:TransitionMatrixMatrixElements}, tells us 
\bse 
    iM = (-ig) \frac{i}{(p_1-p_3)^2 - m^2 +i\epsilon}(-ig).
\ese 

\bbox 
    Show that \Cref{eqn:ScalarYakawaFullyConnectedB} leads to 
    \bse 
        (2\pi)^4 \del^{(4)}\big(p_1+p_2 - p_3-p_4\big) \bigg[ (-ig) \frac{i}{(p_1-p_4)^2 - m^2 +i\epsilon}(-ig) \bigg],
    \ese 
    but note the difference of the intermediate steps. \textit{Hint: Be careful about delta functions.}
\ebox 

So how do the Feynman rules change? Well we note now that $x$ and $y$ don't appear in the final expression so we don't need to label the nodes any more. Then we note in the second step, the two delta functions simply correspond to conservation of 4-momentum at the vertices, i.e. $\del^{(4)}(p_1-p_3-k)$ is the momentum flowing into the $x$ vertex above. Then to get the final answer to integrated over the internal momentum $k$, which gave use the final delta function. 

\mybox{
The \textit{momentum space} Feynman rules for scalar Yakawa scattering $\psi\psi\to \phi \to \psi\psi$ are:
\begin{center}
	\begin{tabular}{@{} C{4cm} C{4cm} C{4cm} @{}}
		\toprule
		 Type & Diagram & Math Expression \\
		\midrule 
		External lines & \btik 
            \midarrow (0,0) -- (2,0);
            \draw[->] (0.5,0.2) -- (1.5,0.2) node [midway, above] {$p$};
            \draw[fill=black] (2,0) circle [radius=0.07];
            \node at (1,-0.5) {and};
            \midarrow (0,-1) -- (2,-1);
            \draw[->] (0.5,-1.2) -- (1.5,-1.2) node [midway, below] {$p$};
            \draw[fill=black] (0,-1) circle [radius=0.07];
        \etik & $1$ \\
        Propagator & \btik 
            \draw[thick, dashed] (0,0) -- (2,0);
            \draw[->] (0.5,0.2) -- (1.5,0.2) node [midway, above] {$k$};
            \draw[fill=black] (0,0) circle [radius=0.07];
            \draw[fill=black] (2,0) circle [radius=0.07];
        \etik & \bse \int \frac{d^4k}{(2\pi)^4} \frac{i}{k^2 - m^2 +i\epsilon}\ese  \\ 
        Vertex & \btik 
            \midarrow (0,0) -- (1,0);
            \midarrow (1,0) -- (2,0);
            \draw[thick, dashed] (1,0) -- (1,-0.8);
            \draw[fill=black] (1,0) circle [radius=0.07] node [above] {$-ig$};
        \etik & $-(ig)$ \\
		\bottomrule
	\end{tabular}
\end{center}
}


So we have the rough\footnote{We will need to be more clear for the case of spin-1/2 particles later.} following procedure to obtain the maths expression from a momentum-space Feynman diagram:
\begin{itemize}
    \item Take the position space Feynman diagram, and given every internal line a momentum labelled by $k_i$.
    \item Next to each vertex write down the coupling strength (in the case of scalar Yakawa this is $-ig$).
    \item Impose momentum conservation at each vertex with a factor of the coupling strength, i.e. write $(-ig)(2\pi)^4\del^{(4)}(\sum_i k_i)$, where the sum is taken over the momentum flowing \textit{into} the vertex. That is, if the momentum flows out, put a minus sign in the sum, as we did in $\del^{(4)}(p_1-p_3-k)$ and $\del^{(4)}(p_2-p_4+k)$.
    \item Include a factor of 
    \bse 
        \int \frac{d^4k}{(2\pi)^4}\frac{i}{k_i^2 - m^2 +i\epsilon}
    \ese 
    for all the internal $\phi$ lines, and the same for internal $\psi$ lines but now with $m\to M$.
\end{itemize}

\br 
    Some authors do not include the integral factor $\int\frac{d^4k}{(2\pi)^4}$ in the propagator term expression, but instead add another step at the end that says "finally integrate over all internal momentum". Of course this amounts to exactly the same thing, but this remark is just included to avoid confusion when comparing to other literature. 
\er 

\bbox 
    Draw the scalar Yakawa scattering Feynman diagrams above in momentum space.
\ebox 
\chapter{LSZ Theorem, Loop \& Amputated Diagrams}

Last lecture we discussed fully connected diagrams and derived the Feynman rules for $\psi\psi \to \psi\psi$ scattering. We also said that there were two other types of conceptually different diagrams, namely \textit{connected} and \textit{disconnected} diagrams. As diagrams they look like:
\begin{center}
    \btik 
        \begin{scope}[xshift=-3.25cm]
            \draw (-3,1.5) -- (3,1.5) -- (3,-3.5) -- (-3,-3.5) -- (-3,1.5);
            \midarrow (-2,0) -- (-1,0);
            \node at (-2.2,0) {$\psi$};
            \midarrow (-1,0) -- (1,0);
            \draw[thick, dashed] (-1,0) .. controls (-0.5,1) and (0.5,1) .. (1,0);
            \node at (0,1) {$\phi$};
            \midarrow (1,0) -- (2,0);
            \node at (0,-0.3) {$\psi$};
            \node at (2.2,0) {$\psi$};
            \midarrow (-2,-2) -- (2,-2);
            \node at (-2.2,-2) {$\psi$};
            \node at (2.2,-2) {$\psi$};
            \node at (0,-3) {\large{Connected}};
        \end{scope}
        %
        \begin{scope}[xshift=3.25cm]
            \draw (-2,1.5) -- (6,1.5) -- (6,-3.5) -- (-2,-3.5) -- (-2,1.5);
            \midarrow (-1,0) -- (1,0);
            \node at (-1.2,0) {$\psi$};
            \node at (1.2,0) {$\psi$};
            \midarrow (-1,-2) -- (1,-2);
            \node at (-1.2,-2) {$\psi$};
            \node at (1.2,-2) {$\psi$};
            \beforemidarrow (3,-1) circle [radius=0.5cm];
            \node at (3,-0.25) {$\psi$};
            \draw[thick, dashed] (3.5,-1) -- (4.5,-1);
            \node at (4,-1.3) {$\phi$};
            \beforemidarrow (5,-1) circle [radius=0.5cm];
            \node at (5,-0.25) {$\psi$};
            \node at (2,-3) {\large{Disconnected}};
        \end{scope}
    \etik 
\end{center}

Intuitively, we would say these don't really correspond to scattering as the initial states don't interact with each other. However we need some way to actually prove this, and this is where the so-called LSZ theorem\footnote{Named after Harry Lehmann, Kurt Symanzik and Wolfhart Zimmermann.} comes in. First we need some definitions and a proposition.

\bd[$n$-Point Green's Function]
    We define the \textit{$n$-point Green's function} to be the time ordered vacuum expectation value of $n$ Heisenberg picture\footnote{Note we need this, otherwise time-ordering doesn't mean anything.} field operator. That is 
    \be  
    \label{eqn:nPointGreensFunction}
        G_n(x_1,...,x_n) := \bra{0} \cT\big[\phi(x_1) ... \phi(x_n)\big] \ket{0}.
    \ee 
\ed 

\bd[Wavefunction Normalisation]
    We define the so-called \textit{wavefunction normalisation} to be 
    \bse 
        Z := |\bra{p_i} \phi(x) \ket{0}|^2,
    \ese 
    so that 
    \be 
    \label{eqn:ZDefinition}
        \bra{p_1,...,p_n} \phi(t=\pm\infty,\Vec{x}) \ket{q_1, ..., q_m} = \lim_{t\to\pm\infty} Z^{-1/2} \, \bra{p_1,...,p_n} \phi(x) \ket{q_1, ..., q_m}.
    \ee 
\ed 

\br 
    What \Cref{eqn:ZDefinition} does is it allows us to turn an asymptotic field $\phi(t=\pm\infty,\Vec{x})$ into a field at a general time under a limit. The factor of $\sqrt{Z}$ is included just to normalise the result, hence the name. \textcolor{red}{For Michael: Is this the correct interpretation of $Z$?}
\er 

\bd[Left-Right Derivative Action]\footnote{This might not be a technical name, it's just what I've decided to call it.}
    Let $\phi_1(x)$ and $\phi_2(x)$ be to fields on the spacetime. Then we define the \textit{left-right derivative action} as 
    \be 
    \label{eqn:LeftRightDerivativeAction}
        \phi_1(x) \lra{\p}_{\mu} \, \phi_2(x) := \phi_1(x) \p_{\mu} \phi_2(x) - \big(\p_{\mu}\phi_1(x)\big)\phi_2(x).
    \ee 
\ed 

\bp 
    For the real scalar field, \Cref{eqn:HeisenbergPhi}, we have 
    \be 
    \label{eqn:aaDaggerLeftRightDerivativeAction}
        \begin{split}
            a_{\Vec{p}} & = \frac{i}{\sqrt{2E_{\vec{p}}}} \, \int d^3 \Vec{x} \, e^{ip\cdot x} \lra{\p}_0 \, \phi(x) \\
            a_{\Vec{p}}^{\dagger} & = -\frac{i}{\sqrt{2E_{\vec{p}}}} \, \int d^3 \Vec{x} \, e^{-ip\cdot x} \lra{\p}_0 \, \phi(x)
        \end{split}
    \ee 
\ep 

\bbox 
    Prove \Cref{eqn:aaDaggerLeftRightDerivativeAction}. \textit{Hint: You will want to use the constraint $p_0=E_{\vec{p}}$ present in \Cref{eqn:HeisenbergPhi} along with $E_{\vec{p}} = \sqrt{\vec{p}^2+m^2} = E_{-\vec{p}}$ at the end.}
\ebox 

We choose here to first state the result of the LSZ theorem and then derive it, as this way the heavy manipulations that follow have some guiding light. 

\bt[LSZ Reduction Formula]
    We can express the S-matrix between an initial $m$-particle state and a $n$-particle final particle state in terms of the $(n+m)$-point Green's function as follows:
    \be 
    \label{eqn:LSZ}
        \begin{split}
            S_{fi} = & \big(iZ^{-1/2}\big)^{n+m} \int dy_1...dy_n \int dx_1...dx_m \, e^{i(p_j\cdot y_j - q_i\cdot x_i)} \\
            & \times \bigg(\prod_{i=1}^m (\p_{x_i}^2+m^2)\bigg) \bigg(\prod_{j=1}^n (\p_{y_j} +m^2) \bigg) G_{n+m}(y_1,...,y_n,x_1,...,x_m) \\
            & + \text{all non-fully connected stuff},
        \end{split}
    \ee 
    where we assume an implicit sum in the exponential.
\et 

Before proving this theorem it is worth stating what it allows us to do. Basically we note that the right-hand side of \Cref{eqn:LSZ} splits into a contribution that contains only information about the fully connected diagrams (we will prove this), and a part that contains only the non-fully connected stuff. We can therefore simply restrict our analysis to processes that are only fully connected in a well defined way. That is, a priori there is no clear way to see that the result of the full S-matrix calculation wont `mix' fully connected and non-fully connected stuff, but the LSZ theorem tells us that it indeed does. This statement does \textit{not} mean that the connected diagrams do not contribute to the full S-matrix value, but simply that we are allowed to restrict ourselves to considering only fully-connected diagrams.

\bq 
    Unfortunately for those who don't like long proofs, this one comes from sheer brute force,\footnote{I'm following the one given by Prof. Weigand, starting on page 49.} so prepare yourself. Firstly we need to clarify the notation that follows: by Dysons formula, the S-matrix is given by a bunch of time evolution operators, and we also take our states to asymptotic (i.e. at $t\to \pm\infty$). Therefore the creation operators that we extract from our initial states \textit{do not} simply act on the final state and annihilate them, as they are separated in time. For this reason we shall use subscripts `in' and `out' on the creation/annihilation operators to avoid confusion. Similarly we will denote the initial/final states using the following notation 
    \bse 
        \bra{p_1...p_n;out}, \qand \ket{q_1...q_m;in}.
    \ese 
    We use this notation as it is the one used by Prof. Weigand in his notes. 
    
    We therefore have 
    \bse 
        \begin{split}
            S_{fi} & = \braket{p_1...p_n;out}{q_1...q_m;in} \\
            & = \sqrt{2E_{\vec{q}_1}}\bra{p_1...p_n;out} a^{\dagger}_{\vec{q}_1} \ket{q_2...q_m;in} \\
            & = \frac{1}{i} \int d^3\vec{x}_1 \bra{p_1...p_n;out} e^{-iq_1\cdot x_1} \lra{\p}_0 \, \phi(t=-\infty,\vec{x}) \ket{q_2...q_m;in} \\
            & = \frac{1}{i} \lim_{t\to-\infty} Z^{-1/2} \, \int d^3\vec{x}_1 \bra{p_1...p_n;out} e^{-iq_1\cdot x_1} \lra{\p}_0 \, \phi(x) \ket{q_2...q_m;in} \\
            & = \frac{1}{i\sqrt{Z}} \lim_{t\to-\infty}  \, \int d^3\vec{x}_1 e^{-iq_1\cdot x_1} \lra{\p}_0  \bra{p_1...p_n;out} \phi(x) \ket{q_2...q_m;in}
        \end{split}
    \ese 
    where we have used the definition $\ket{q} = \sqrt{2E_{\vec{q}}}a^{\dagger}_{\vec{q}}\ket{0}$ and \Cref{eqn:aaDaggerLeftRightDerivativeAction,eqn:ZDefinition}. Note on the third line we have $\phi(t=-\infty,\vec{x})$, this is because we want to consider an initial state, which we take to be at $t=-\infty$. Indeed for this exact reason, we should really denote the creation operator on the second line as something like 
    \bse 
        (a_{\text{in}})^{\dagger}_{\vec{q}_1},
    \ese 
    so that we know it only acts on in/initial states.
    
    What we now want to do is get the $\phi(x)$ to act on a final state, in which case we need to change the limit from $t\to-\infty$ to $t\to + \infty$. We achieve this by using the simple result 
    \bse 
        \lim_{t\to-\infty} f(t) = \lim_{t\to+\infty} f(t) - \lim_{T\to\infty} \int_{-T}^{T} \frac{d}{dt}f(t),
    \ese 
    so our S-matrix splits into two terms:
    \bse 
        \begin{split}
            S_{fi} = & \frac{1}{i\sqrt{Z}} \lim_{t\to+\infty} \, \int d^3\vec{x}_1 e^{-iq_1\cdot x_1} \lra{\p}_0 \bra{p_1...p_n;out} \phi(x) \ket{q_2...q_m;in} \\
            & \qquad + \frac{1}{i\sqrt{Z}} \, \int d^4x_1 \, \p_0\Big( e^{-iq_1\cdot x_1} \lra{\p}_0 \bra{p_1...p_n;out} \phi(x) \ket{q_2...q_m;in}\Big),
        \end{split}
    \ese
    where we notice that the second term contains a $4$-integral. Let's call the first $A$ and the second $B$, and consider them in turn. In $A$ we can use \Cref{eqn:ZDefinition} backwards to get 
    \bse 
        A = \frac{1}{i\sqrt{Z}} \int d^3\vec{x}_1 e^{-iq_1\cdot x_1} \lra{\p}_0  \bra{p_1...p_n;out} \phi(t=+\infty,\vec{x}) \ket{q_2...q_m;in},
    \ese 
    and then we can use \Cref{eqn:aaDaggerLeftRightDerivativeAction} backwards to give 
    \bse 
        A = \sqrt{2E_{\vec{q}_1}}  \bra{p_1...p_n;out} (a_{\text{out}})^{\dagger}_{\vec{q}_1} \ket{q_2...q_m;in},
    \ese 
    where, in agreement with the comment made above, we have included a subscript `out' to tell us this acts on the final states. Now the action of a creation operator to the left gives a delta function, and so we get 
    \bse 
        A = 2E_{\vec{q}_1}(2\pi)^3 \sum_{k=1}^{n}\del^{(3)}(\vec{p}_k -\vec{q}_1) \braket{p_1...\hat{p}_k...p_n;out}{q_2...q_m;in},
    \ese 
    where we have used the standard maths notation where a hatted entry in a string of terms is missing (i.e. $\hat{p}_k$ is missing from the final states). This corresponds to a connected, \textit{but not fully connected}, diagram as we have the $k$-th final particle just corresponding to the $1^{\text{st}}$ initial particle (it is just a straight line in a diagram). 
    
    Now what about $B$? Well we use \Cref{eqn:LeftRightDerivativeAction} to get
    \bse
        \begin{split}
            B & = \frac{1}{i\sqrt{Z}}\, \int d^4x_1 \, \p_0\Big( e^{-iq_1\cdot x_1} \p_0 \la...\ra - \p_0\big(e^{-iq_1\cdot x_1}\big)\la...\ra \Big) \\
            & = \frac{1}{i\sqrt{Z}} \, \int d^4x_1 \, \Big( e^{-iq_1\cdot x_1} \p_0^2 \la...\ra - \p_0^2\big(e^{-iq_1\cdot x_1}\big)\la...\ra \Big)
        \end{split}
    \ese 
    Next note that 
    \bse 
        \p_0^2 e^{-ip\cdot x} = -p_0^2 e^{-ip\cdot x} = (p^2 +(\vec{p})^2) e^{-ip\cdot x} = (m^2 - \nabla^2)e^{-ip\cdot x},
    \ese 
    which we can apply to the second term on the second line above. We then integrate by parts twice\footnote{Note, as always, we assume boundary terms vanish at spatial infinity.} to move the $\nabla^2$ to act on $\la...\ra$ and then use 
    \bse 
        \p^2 := \p_0^2 - \nabla^2
    \ese 
    to give us 
    \bse 
        B = \frac{1}{i\sqrt{Z}} \int d^4 x_1 e^{-iq_1\cdot x_1} \big(\p^2+m^2)\bra{p_1...p_n;out} \phi(x) \ket{q_2...q_m;in}
    \ese 
    This is exactly of the form of one of the product terms in \Cref{eqn:LSZ}! We can then apply a similar argument to all the other initial states, noting that 
    \bse 
        [a^{\dagger}_{\vec{q}_i},a^{\dagger}_{\vec{q}_j}] = 0 \qquad \forall i,j \in \{1,...,m\},
    \ese 
    to obtain the full product from $1$ to $m$ in \Cref{eqn:LSZ}. 
    
    We now need to do the same thing for the final states. However now we have 
    \bse 
        \bra{p} = \sqrt{2E_{\vec{p}}} \bra{0} (a_{\text{out}})_{\vec{p}},
    \ese
    and from 
    \bse 
        [a_{\vec{p}_i},a^{\dagger}_{\vec{q}_j}] \neq 0
    \ese
    we cannot simply commute the annihilation operators past $\phi(x)$ to get them to act on the initial state, that is 
    \bse 
        \begin{split}
            \bra{p_2...p_n;out}\phi(y_1^0=\infty,&\vec{y}_1)\phi(x_1^0=-\infty,\vec{x}_1)\ket{q_1...q_m;in} \\
            & \neq \bra{p_2...p_n;out}\phi(x_1^0=-\infty,\vec{x}_1)\phi(y_1^0=\infty,\vec{y}_1)\ket{q_1...q_m;in}.
        \end{split}
    \ese 
    This is when we notice that \Cref{eqn:LSZ} contains a Green's function, which contains time-ordering and so we realise we are saved! More explicitly, we will get something of the form 
    \bse 
        \begin{split}
            \lim_{x_1^0\to-\infty} \lim_{y_1^0\to+\infty} \int d^3 \vec{x}_1 \int d^3 \vec{y}_1 e^{-iq_1\cdot x_1} \lra{\p}_{x_1^0} \, e^{ip_1\cdot y_1} \lra{\p}_{y_1^0} \, \bra{p_2...p_n;out}\phi(y_1)\phi(x_1)\ket{q_2...q_m;in} \\
            = \lim_{x_1^0\to-\infty} \int d^3 \vec{x}_1 e^{-iq_1\cdot x_1} \lra{\p}_{x_1^0} \, \lim_{y_1^0\to+\infty} \int d^3 \vec{y}_1  e^{ip_1\cdot y_1} \lra{\p}_{y_1^0} \, \bra{p_2...p_n;out}\phi(y_1)\phi(x_1)\ket{q_2...q_m;in}
        \end{split}
    \ese 
    where we notice that the $x_1^0$ limit is taken after the $y_1^0$ one, so during the latter we can take $x_1^0$ to be some fixed finite value, so we can just consider the $y_1$ part of the expression. Now consider 
    \bse 
        \begin{split}
            \lim_{T\to\infty} \int_{-T}^T dy_1^0 \p_{y_1^0} & \bigg(\int d^3\vec{y}_1 e^{ip_1\cdot y_1} \lra{\p}_{y_1^0} \bra{p_2...p_n;out}\cT\big[ \phi(y_1)\phi(x_1)\big] \ket{q_2...q_m;in}\bigg) \\
            = & \lim_{y_1^0\to+\infty} \int d^3\vec{y}_1 e^{ip_1\cdot y_1} \lra{\p}_{y_1^0} \bra{p_2...p_n;out} \phi(y_1)\phi(x_1) \ket{q_2...q_m;in} \\
            & \qquad  - \lim_{y_1^0\to-\infty} \int d^3\vec{y}_1 e^{ip_1\cdot y_1} \lra{\p}_{y_1^0} \bra{p_2...p_n;out} \phi(x_1)\phi(y_1) \ket{q_2...q_m;in}.
        \end{split}
    \ese 
    The first term on the right hand side is what we want, (compare it to the proof above), while the second term on the right hand side will again give a connected, but not fully connected, diagram as we get a $a_{\vec{p}_1}$ term acting on the initial state. So we finally see that we can replace the expression we had above for the Green's function (the time-ordered expression) at the cost of picking up a connected, but not fully connected, term. Doing this for all the final states will then give us all the terms in the product from $j=1$ to $j=n$ needed in \Cref{eqn:LSZ}. Note also that all the other factors needed for \Cref{eqn:LSZ} also appear, i.e. all the exponentials and the $iZ^{-1/2}$ factors.
\eq 

Ok so we have that the fully connected part of the scattering process is given by a $(n+m)$-point Green's function. However the LSZ formula is quite a large formula and so we want to simplify it somehow. We do this by introducing the \textit{truncated Green's function} as follows
\be 
\label{eqn:TruncatedGreensFunction}
    \begin{split}
        G_{n+m}(y_1,...,y_n,x_1,...,x_m) = \int & d^4 z_1 ... d^4 dz_{m+n} \Delta_F(x_1-z_1) ... \Delta_F(x_m-z_m) \\
        & \times \Delta(y_1-z_{m+1}) ... \Delta_F(y_n-z_{n+m}) \widetilde{G}_{n+m}(z_1,...,z_{n+m}),
    \end{split}
\ee
where $\Delta_F$ is the Feynman propagator and $\widetilde{G}_{n+m}$ is our truncated $(n+m)$-point Green's functions. Essentially what we're doing is saying is propagate our initial states to the points $\{z_1,...,z_m\}$ and similarly the final states with $\{z_{m+1},...,z_{m+n}\}$, and characterise all the interaction stuff in the middle using some Green's function, $\widetilde{G}$, as we've tried to indicate below diagrammatically. 
\begin{center}
    \btik[scale=0.8]
        \draw[thick, fill = gray!40, opacity = 0.8] (0,0) circle [radius=2cm];
        \node at (0,0) {\Large{$\widetilde{G}_{5+4}$}};
        \node at (-6.3,0) {$x_3$};
        \node at (-1.7,0) {$z_3$};
        \draw[thick] (-6,0) -- (-2,0);
        \draw[fill=black] (-2,0) circle [radius=0.07cm];
        \begin{scope}[rotate around={20:(0,0)}]
            \node at (-6.3,0) {$x_4$};
            \node at (-1.7,0) {$z_4$};
            \draw[thick] (-6,0) -- (-2,0);
            \draw[fill=black] (-2,0) circle [radius=0.07cm];
        \end{scope}
        \begin{scope}[rotate around={40:(0,0)}]
            \node at (-6.3,0) {$x_5$};
            \node at (-1.7,0) {$z_5$};
            \draw[thick] (-6,0) -- (-2,0);
            \draw[fill=black] (-2,0) circle [radius=0.07cm];
        \end{scope}
        \begin{scope}[rotate around={-20:(0,0)}]
            \node at (-6.3,0) {$x_2$};
            \node at (-1.7,0) {$z_2$};
            \draw[thick] (-6,0) -- (-2,0);
            \draw[fill=black] (-2,0) circle [radius=0.07cm];
        \end{scope}
        \begin{scope}[rotate around={-40:(0,0)}]
            \node at (-6.3,0) {$x_1$};
            \node at (-1.7,0) {$z_1$};
            \draw[thick] (-6,0) -- (-2,0);
            \draw[fill=black] (-2,0) circle [radius=0.07cm];
        \end{scope}
        %
        \node at (6.3,0) {$y_3$};
        \node at (1.7,0) {$z_8$};
        \draw[thick] (6,0) -- (2,0);
        \draw[fill=black] (2,0) circle [radius=0.07cm];
        \begin{scope}[rotate around={20:(0,0)}]
            \node at (6.3,0) {$y_2$};
            \node at (1.7,0) {$z_7$};
            \draw[thick] (6,0) -- (2,0);
            \draw[fill=black] (2,0) circle [radius=0.07cm];
        \end{scope}
        \begin{scope}[rotate around={55:(0,0)}]
            \node at (6.3,0) {$y_1$};
            \node at (1.7,0) {$z_6$};
            \draw[thick] (6,0) -- (2,0);
            \draw[fill=black] (2,0) circle [radius=0.07cm];
        \end{scope}
        \begin{scope}[rotate around={-30:(0,0)}]
            \node at (6.3,0) {$y_4$};
            \node at (1.7,0) {$z_9$};
            \draw[thick] (6,0) -- (2,0);
            \draw[fill=black] (2,0) circle [radius=0.07cm];
        \end{scope}
    \etik 
\end{center}

Why would we introduce the truncated Green's function? Well we recall that the Feynman propagator is a Green's function of the Klein-Gordan operator:
\bse 
    (\p^2+m^2)\Delta_F(x-z) = 0,
\ese 
so in the LSZ formula all of the Klein-Gordan operators will just give us delta functions, which we can integrate over, leaving us with 
\be
\label{eqn:SMatrixTruncatedGreensFunction}
    \begin{split}
        S_{if}^{F.C.} & = \big(iZ^{-1/2}\big)^{n+m} \int dy_1...dy_n\int dx_1...dx_m e^{i(p_j\cdot y_j - q_i\cdot x_i)} \widetilde{G}_{n+m}(x_1...,x_m, y_1,...,y_n) \\
        & = \big(iZ^{-1/2}\big)^{n+m} \widetilde{G}_{n+m} (q_1,...,q_m, p_1, ... , p_n),
    \end{split}
\ee 
where the second line follows from the delta functions we get from the exponentials, and where the superscript F.C. means "fully connected". This result tells us that we can compute the entire fully connected S-matrix by just considering the momentum space truncated Green's function! This is a massively convenient result. 

\br 
    It is worth stressing again that the LSZ theorem does \textit{not} tell us that the connected, but not fully connected, diagrams vanish but that we can ignore them if we want to just consider a scattering process where everything interacts. In fact the former diagrams will contribute to renormalisation and so are very important. 
\er 

\section{Loop Diagrams}

So far we have only considered what are known as \textit{tree level diagrams}. This essentially just means that all of our vertices are connected to our external lines. However this is clearly not the only type of diagram we can have and we can have vertices internally. For example in the scalar Yakawa theory the following is a valid diagram to order $g^2$:
\begin{center}
    \btik 
        \midarrow (-3,0) -- (0,0);
        \node at (-3.3,0) {$\psi$};
        \draw[->] (-2,0.3) -- (-1,0.3) node [midway, above] {$p_1$};
        \midarrow (0,0) -- (3,0);
        \node at (3.3,0) {$\psi$};
        \draw[->] (1,0.3) -- (2,0.3) node [midway, above] {$p_3$};
        \draw[thick, dashed] (0,0) -- (0,-1);
        \draw[->] (-0.3,-0.15) -- (-0.3,-0.85) node [midway, left] {$k_1$};
        \midarrow (0,-1.5) circle [radius=0.5cm];
        \draw[->] (0.8,-1.5) arc (0:60:0.7cm);
        \node at (1,-1) {$\ell_1$};
        \draw[->, rotate around={180:(0,-1.5)}] (0.8,-1.5) arc (0:60:0.7cm);
        \node at (-1,-2) {$\ell_2$};
        \midarrow (0,-1.5) circle [radius=-0.5cm];
        \draw[thick, dashed] (0,-2) -- (0,-3);
        \draw[->] (0.3,-2.15) -- (0.3,-2.85) node [midway, right] {$k_2$};
        \midarrow (-3,-3) -- (0,-3);
        \node at (-3.3,-3) {$\psi$};
        \draw[->] (-2,-3.3) -- (-1,-3.3) node [midway, below] {$p_2$};
        \midarrow (0,-3) -- (3,-3);
        \node at (3.3,-3) {$\psi$};
        \draw[->] (1,-3.3) -- (2,-3.3) node [midway, below] {$p_4$};
    \etik 
\end{center}
\noindent where we haven't labelled the internal particles and the $(-ig)$ factors to avoid cluttering the diagram. This is called a \textit{loop diagram} for obvious reasons. Loop diagrams are a bit of a pain because we don't have enough delta functions to remove all of our integrals. That is, recall that in the momentum space Feynman rules we include a factor of\footnote{Note that before we only defined this for $\phi$ propagators, but here we say we get the same factor for $\psi$ ones but we simply replace the mass with the $\psi$ mass. From now on we shall use $m_{\phi}=m$ and $m_{\psi}=M$.}
\bse 
    \int \frac{d^4k_i}{(2\pi)^2} \frac{i}{k_i^2-m_{\phi/\psi}^2 + i\epsilon}
\ese
for each internal line and a delta function of the momentum flowing in at each vertex. For example, in the above diagram the $k_1$ and $k_2$ momentum are obtained from the delta functions with the external momenta, i.e. $k_1=p_1-p_3$ and $k_2=p_4-p_2$. However we cannot find the values of $\ell_1$ and $\ell_2$, but at best can obtain a relation between them. This is just because $\ell_1$ and $\ell_2$ will both appear in two delta functions, namely 
\bse 
    \del^{(4)}(k_1+\ell_1-\ell_2), \qand \del^{(4)}(\ell_2-k_2-\ell_1),
\ese
and so when we integrate over one of the two $\ell$s we kill both delta functions and we are left with the integral over the other $\ell$ and no delta function to remove it! To be more explicit, if we did the integral over $\ell_2$, say, we would get 
\bse 
    \ell_2 = k_1 + \ell_1, \qand \ell_2 = k_2 + \ell_1,
\ese 
so if we equate the two terms we just get 
\bse 
    k_1 = k_2 \qquad \implies \qquad p_1 + p_2 = p_3+p_4,
\ese
where the implication comes from the momentum conservation on the $k$s. This just tells us the $4$-momentum initially and finally are equivalent, and does not allow us to conclude what $\ell_1$ is. 

\br 
    As we said when we introduced the momentum space Feynman rules, it is common to alter the rules slightly to so just include the factor 
    \bse 
        \frac{i}{k^2-m^2+i\epsilon}
    \ese 
    for each propagator and then impose momentum conservation at each vertex. If we use these rules, then when we have loops we insert the additional rule of "and then integrate over all undetermined momenta", i.e. $\ell_1$ in the above example. 
\er 

We also get loop diagrams in connected, but not fully connected, diagrams, as in the next exercise. 

\bbox 
    Show that this diagram
    \begin{center}
        \btik 
            \midarrow (-3,0) -- (-1,0);
            \node at (-3.2,0) {$\psi$};
            \midarrow (-1,0) -- (1,0);
            \draw[thick, dashed] (-1,0) .. controls (-0.5,1) and (0.5,1) .. (1,0);
            \draw[->, rotate around={-30:(-1.2,0.2)}] (-1.2,0.2) arc (0:-30:-2) node [midway, above] {$k$};
            \draw[->] (-2.8,-0.3) -- (-1.8,-0.3) node [midway,below] {$p_1$};
            \draw[->] (1.8,-0.3) -- (2.8,-0.3) node [midway,below] {$p_2$};
            \draw[->] (-0.5,-0.3) -- (0.5,-0.3) node [midway,below] {$\ell$};
            \midarrow (1,0) -- (3,0);
            \node at (3.2,0) {$\psi$};
        \etik 
    \end{center}
    \noindent corresponds to 
    \bse 
        (2\pi)^4 \del^{(4)}(p_2-p_1)\Bigg[ \int \frac{d^4\ell}{(2\pi)^4} (-ig)^2 \bigg( \frac{i}{(p_1-\ell)^2- m^2 +i\epsilon} \bigg) \bigg( \frac{i}{\ell^2 - M^2 +i\epsilon} \bigg) \Bigg]
    \ese
\ebox  

\br 
    Note that the result of the previous exercise diverges logarithmically in $\ell$. This is easily seen because we essentially have $1/\ell^4$ and we do $d^4\ell$. This kind of divergence is known as a \textit{UV divergence}.
\er 

\section{Amputated Diagrams}

Note in all the Feynman diagrams we have just drawn, we always draw the ingoing and outgoing particles as `pure' lines, by which I mean nothing is attached to them. An example of something we haven't drawn is the following:
\begin{center}
    \btik 
        \midarrow (-2,0) -- (-1,0);
        \midarrow (-1,0) -- (1,0);
        \draw[thick, dashed] (-1,0) .. controls (-0.5,1) and (0.5,1) .. (1,0);
        \midarrow (1,0) -- (3,0);
        \midarrow (3,0) -- (5,0);
        \draw[thick, dashed] (3,0) -- (3,-2);
        \midarrow (-2,-2) -- (3,-2);
        \midarrow (3,-2) -- (5,-2);
    \etik 
\end{center}

The reason we have never considered these is because what we've been drawing are so-called \textit{amputated diagrams}. You get amputated diagrams by saying "can I cut through the diagram and remove something by only cutting through one line?" If the answer is yes, do it. So for the above example we would cut along the dotted red line below:
\begin{center}
    \btik 
        \midarrow (-2,0) -- (-1,0);
        \midarrow (-1,0) -- (1,0);
        \draw[thick, dashed] (-1,0) .. controls (-0.5,1) and (0.5,1) .. (1,0);
        \midarrow (1,0) -- (3,0);
        \midarrow (3,0) -- (5,0);
        \draw[thick, dashed] (3,0) -- (3,-2);
        \midarrow (-2,-2) -- (3,-2);
        \midarrow (3,-2) -- (5,-2);
        \draw[ultra thick, red, dashed] (3,1) -- (1,-1);
    \etik 
\end{center}
\noindent which gives us the scattering diagram we drew before. It is important that you only cut through $1$ line, i.e. we can not amputate the following 
\begin{center}
    \btik 
        \midarrow (-2,0) -- (-1,0);
        \midarrow (-1,0) -- (0,0);
        \midarrow (0,0) -- (1,0);
        \draw[thick, dashed] (-1,0) .. controls (-0.5,1) and (0.5,1) .. (1,0);
        \midarrow (1,0) -- (2,0);
        \draw[thick, dashed] (0,0) -- (0,-2);
        \midarrow (-2,-2) -- (0,-2);
        \midarrow (0,-2) -- (2,-2);
    \etik 
\end{center}
The reason we amputate diagrams is because the LSZ theorem tells us to consider only amputated diagrams. This is seen from the truncated Greens function expression: we consider just propagating our initial and final states in/out without any interactions and then only consider the interactions that happen between all the particles in the internal parts, see the truncated $\widetilde{G}$ picture above. 

% -------------------------------------------------------------------
% Bibliography/Further Readings
% -------------------------------------------------------------------

\chapter*{Useful Texts \& Further Readings}

\section*{Quantum Field Theory}
\begin{itemize}
    \item Peskin, Michael E. \textit{An introduction to quantum field theory}. CRC Press, 2018.
    \item Zee, Anthony. \textit{Quantum field theory in a nutshell}. Vol. 7. Princeton university press, 2010.
    \item Schwartz, Matthew D. \textit{Quantum field theory and the standard model}. Cambridge University Press, 2014.
\end{itemize}




%\bibliographystyle{agsm} 
%\bibliography{mybibliography} 
%\printbibliography[heading=bibintoc]


% -------------------------------------------------------------------
% Appendices
% -------------------------------------------------------------------

%\begin{appendices}
%\chapter{An Appendix of Some Kind}

\lipsum[1-3]  % Replace with your text

%\end{appendices}

\end{document}
